\documentclass[
headings=optiontohead,              % allows double headers
12pt,                               % fontsize 
DIV=13,                             % koma script diveider amount. tells koma how much of the site can be written to
twoside=false,                      % if set to true, automatically formats as book style with different left and right pages
open=right,                         % starting page on twosided texts 
BCOR=00mm,                          % correction that accounts for the center of the pages being glued in
toc=bibliographynumbered            % bibliography gets a number and is listed in the table of contents
]{scrreport}

\usepackage[utf8]{inputenc}                     % correct encoding of output, technically not needed anymore
\usepackage[T1]{fontenc}                        % correct encoding of output, technically not needed anymore
\usepackage[english]{babel}                     % font that supports English
\usepackage{upgreek}                            % non-cursive Greek letters
\usepackage[stretch=10,shrink=10,protrusion=true,expansion=true,final]{microtype} % prettier block format
\usepackage{hyperref}                           % links for everything
\usepackage{color}                              % allows for setting in different colors
\usepackage[autooneside=false,automark]{scrlayer-scrpage} % page-style with "Kolumnentitel" (title of current chapter is displayed at the top)
\usepackage[sb]{libertinus}                     % use the font libertinus (needs to be installed from the web)
\usepackage[slantedGreek]{libertinust1math}     % math mode improvement for libertinus
\usepackage{siunitx}                            % physical units setting
\usepackage{icomma}                             % commas in lists get extra space if needed                        
\usepackage{amsfonts,amssymb,amstext,amsmath,amsthm,bbm} % better math mode (\mathrm and \text) and symbols
\usepackage{xspace}                             % works to improve own commands and provides "\xspace"-command, that puts a space if needed
\usepackage{ifthen}                             % more control over non-obligatory parameters
\usepackage{titling}                            % get title values as macros
\usepackage[onehalfspacing]{setspace}           % control the spacing between lines and in enumeration lists
\usepackage[backend=biber, style=phys, biblabel=brackets]{biblatex} % citations with "modern" backend and an physics-accepted citation style
\usepackage{graphicx}                           % work with graphics 
\usepackage{ragged2e}                           % ragged-commands (when no block format is wanted)
\usepackage{pdfpages}                           % allows including of pdfs into this pdf
\usepackage{booktabs}                           % better table formatting
\usepackage{multicol}                           % allows for the definition of multi-columns in tables
\usepackage{multirow}                           % allows for the definition of multi-row-tables instead of just multi-column
\usepackage[section]{placeins}                  % provides the command "\FloatBarrier" to control the end of floatable regions for figures/tables
\usepackage{float}                              % provides the "H" option for forcing placement of a figure
\usepackage{floatpag}                           % make it possible for float-pages to not have a page number
\usepackage{url}                                % sometimes needed by biblatex, technically no longer needed
\usepackage{minted}                             % nice code highlighting (needs Python Package to compile!!)
\usepackage{accents}                            % better control over accents
\usepackage{mathtools}                          % more math control possibilities
\usepackage[autostyle=true]{csquotes}           % context-sensitive-quotes -> quotation marks that are set correctly for the context
\usepackage{physics}                            % bra-ket and more
\usepackage{nicematrix}                         % label row/cols on matrix
\usepackage{caption}                            % caption of different environments
\usepackage{tikz}
\usetikzlibrary{calc}
  % required packages

\title{Classical networks for the Hubbard model with a tilted potential}
\author{Jonas Kell}
\date{15$^\text{th}$ December 2024} %TODO
\newcommand{\pdfkeywords}{neural networks, ai, physics, quantum mechanics, machine learning, analytical calculations, classical networks, hubbard model}
\newcommand{\pdfsubject}{Master Thesis}

\graphicspath{{./../shared-latex/images/}}              % custom paths for folders in that graphics can be found

\sisetup                                % setup for siunitx
{
detect-all,
locale=US,                              % language setup for siunitx
range-phrase={ \text{to} },             % word that is put into an si range
range-units = single,                 % better display of error ranges
per-mode=symbol-or-fraction,            % more dynamic frac usage in inline/displaymath mode
separate-uncertainty,                   % for better +- , \pm when including an error range 
}

\AtBeginDocument{
	\let\mathbb\relax
	\DeclareMathAlphabet\PazoBB{U}{fplmbb}{m}{n}
	\newcommand{\mathbb}{\PazoBB}
}       %more options to the \mathbb command

\setminted[]{
    xleftmargin=0cm,
    xrightmargin=0cm,
    frame=single,
    framesep=.25cm,
    linenos,
    tabsize=2,
    breaklines,
    breakafter=.],
    breakaftersymbolpre= ,
}           %configure the minted code-highlighting style

\addbibresource{literature.bib}              %initialize bibtex with correct file

\NiceMatrixOptions{
code-for-first-row = \color{dblue} ,
code-for-last-row = \color{dblue} ,
code-for-first-col = \color{dblue} ,
code-for-last-col = \color{dblue}
}
                % another file that holds the package/document configuration
\linespread{1.1}                                    % line-spacing can be controlled here

\clubpenalty10000                                   % Schusterjunge, orphan
\widowpenalty10000                                  % Hurenkind, Witwe
\displaywidowpenalty=10000                          % Make document obey stricter rules considering "Schusterjungen" and "Hurenkinder"
\setcounter{biburlnumpenalty}{100}
\setcounter{biburlucpenalty}{100}
\setcounter{biburllcpenalty}{100}
\renewcommand{\topfraction}{0.8}                    % allows for more chilled "text to image ratio" 
\renewcommand{\bottomfraction}{0.8}
\renewcommand{\textfraction}{0.1}
\renewcommand{\floatpagefraction}{0.8}

% \renewcaptionname{ngerman}{\figurename}{Abb.}       %"Figure" becomes "Fig." in English
\setcapindent{0cm}                                  %useful if image captions have multiple lines. Removers indentation below "Fig."
\setlength{\parindent}{0cm}                         %removes indentation at start of new paragraphs
\setlength{\abovecaptionskip}{0.0cm}

%bibliography slots are redefined/modified here
\DeclareFieldFormat{journaltitle}{\textsl{#1}\isdot}
\DeclareFieldFormat{titlecase}{{#1}}


%COLORS
\definecolor{dblue}{rgb}{0,0,0.5} % #000080
\definecolor{dgray}{rgb}{0.5,0.5,0.5}
\definecolor{lgray}{rgb}{0.8,0.8,0.8}


%overwrite the coma-script definitions
\addtokomafont{pagehead}{\normalfont\color{dgray}}                  %overwrite the coma-script definitions
\addtokomafont{sectioning}{\rmfamily\color{dblue}\boldmath}         %rmfamily puts headings in "normal" "serif-font" instead of "sans-serif"  boldmath ensures a bold math font in subscripts
\addtokomafont{captionlabel}{\bfseries\footnotesize}                %better Fig. format
\addtokomafont{caption}{\footnotesize}                          


% headline spacing
\RedeclareSectionCommand[beforeskip=0cm,afterskip=1cm]{chapter}                % another file that holds format information
%! Ref-Commands 
\newcommand*{\fullref}[1]{\hyperref[{#1}]{\textit{\autoref*{#1} \nameref*{#1}}}}
\newcommand*{\fullpage}[1]{\hyperref[{#1}]{Seite \pageref*{#1}}}
\newcommand*{\fullpages}[1]{\hyperref[{#1}]{Seiten \pageref*{#1}ff}}

%! Math operators and other small conveniences
\newcommand\thickbar[1]{\accentset{\rule{.6em}{.8pt}}{#1}}
\DeclareMathOperator{\ggt}{ggT}
\DeclareMathOperator{\kgv}{kgV}
\DeclarePairedDelimiter\ceil{\lceil}{\rceil}
\DeclarePairedDelimiter\floor{\lfloor}{\rfloor}
\renewcommand*{\arraystretch}{0.8}

\newcommand{\filepath}[2]{\colorbox{lgrey}{#1: \texttt{#2}}\xspace}

%! replace German quotation marks with English ones (I always write the left and right quotes explicitly with the German command variants as they are nice to debug and already in my macros. Then in the end, I can swap them with the desired replacement quotes of my choice)
\renewcommand*{\glqq}{``}
\renewcommand*{\grqq}{''}

\newcommand{\bigo}[1]{\ensuremath{\mathcal{O}(#1)}}


%! qm commands
\newcommand{\up}{\ensuremath{\uparrow}\xspace}
\newcommand{\down}{\ensuremath{\downarrow}\xspace}

\newcommand{\difft}[1]{\ensuremath{\frac{\text{d}#1}{\text{d}t}}\xspace}

\newcommand{\schroedingerPicture}{\text{S}}
\newcommand{\interactionPicture}{\text{I}}
\newcommand{\heisenbergPicture}{\text{H}}

\newcommand{\generalop}[7]{\ensuremath{\hat{\text{#1}}^{#2\xspace #3}_{#4#5#6}#7}\xspace}
\newcommand{\withspinop}[4][]{\generalop{c}{#4}{#1}{#2}{,\,}{#3}{}}
\newcommand{\withspinhcop}[4][]{\generalop{h}{#4}{#1}{#2}{,\,}{#3}{}}
\newcommand{\cop}[3][]{\generalop{c}{#3}{#1}{#2}{}{}{}}
\newcommand{\bop}[3][]{\generalop{b}{#3}{#1}{#2}{}{}{}}
\newcommand{\hop}[3][]{\generalop{h}{#3}{#1}{#2}{}{}{}}
\newcommand{\dop}[3][]{\generalop{d}{#3}{#1}{#2}{}{}{}}
\newcommand{\nop}[3][]{\generalop{n}{}{#1}{#2}{,\,}{#3}{}}
\newcommand{\Aop}[1][]{\generalop{A}{}{#1}{}{}{}{}}
\newcommand{\Bop}[1][]{\generalop{B}{}{#1}{}{}{}{}}
\newcommand{\pauli}[2][]{\generalop{$\sigma$}{\text{#2}}{}{#1}{}{}{}}
\newcommand{\ObservableOp}{\generalop{$\mathcal{O}$}{}{}{}{}{}{}}
\newcommand{\copOfT}[3][]{\generalop{c}{#3}{#1}{#2}{}{}{(t)}}
\newcommand{\hopOfT}[3][]{\generalop{h}{#3}{#1}{#2}{}{}{(t)}}
\newcommand{\dopOfT}[3][]{\generalop{d}{#3}{#1}{#2}{}{}{(t)}}
\newcommand{\nopOfT}[3][]{\generalop{n}{}{#1}{#2}{,\,}{#3}{(t)}}
\newcommand{\AopOfT}[1][]{\generalop{A}{}{#1}{}{}{}{(t)}}
\newcommand{\localObservable}[2]{\ensuremath{{\ObservableOp}_\text{loc}(#1,\,#2)}\xspace}
\newcommand{\doubleOccupationOperator}[1]{\ensuremath{{\ObservableOp}_\text{do}(#1)}\xspace}
\newcommand{\singleOccupationOperator}[2]{\ensuremath{{\ObservableOp}_\text{so}(#1,\,#2)}\xspace}
\newcommand{\spinPolarizedKineticsOperator}[3]{\ensuremath{{\ObservableOp}_\text{sp-kin}(#1 \leftrightarrow #2,\,#3)}\xspace}
\newcommand{\spinPolarizedKineticsOperatorDir}[3]{\ensuremath{{\ObservableOp}_\text{sp-kin, dir}(#1 \rightarrow #2,\,#3)}\xspace}
\newcommand{\one}[1][]{\ensuremath{\mathbbm{1}_{#1}}\xspace}


\newcommand{\hamiltonian}{\ensuremath{\mathcal{H}}\xspace}
\newcommand{\pictureHamiltonian}[1][]{\ensuremath{\hamiltonian^{#1}}\xspace}
\newcommand{\Vhamiltonian}[1][]{\ensuremath{\hat{\text{V}}^{#1}}\xspace}
\newcommand{\HzeroHamiltonian}[1][]{\ensuremath{\hamiltonian_0^{#1}}\xspace}
\newcommand{\VhamiltonianOf}[2][]{\ensuremath{\hat{\text{V}}^{#1}(#2)}\xspace}
\newcommand{\VhamiltonianOfT}[1][]{\VhamiltonianOf[#1]{t}}
\newcommand{\HNOft}[1][N]{\ensuremath{\hamiltonian_{#1}(t)}\xspace}
\newcommand{\HNOftOrder}[2][N]{\ensuremath{\hamiltonian_{#1}^{#2}(t)}\xspace}
\newcommand{\Heff}[3]{\ensuremath{{\hamiltonian_\text{eff}}^{#3}({#1}#2)}\xspace}
\newcommand{\HeffOft}[1][N]{\Heff{#1}{,\,t}{}}
\newcommand{\Hvcn}[2][N]{\ensuremath{{\hamiltonian_\text{VCN}}({#1},\,{#2})}\xspace}
\newcommand{\HeffOftStar}[1][N]{\Heff{#1}{,\,t}{\ast}}
\newcommand{\Heffvcn}[2]{\ensuremath{{\hamiltonian_\text{eff}^\text{VCN}}({#1},\,#2)}\xspace}
\newcommand{\VhamiltonianPartA}[2]{\ensuremath{\hat{\text{V}}_\text{Part A}\left(#1,\,#2\right)}\xspace}
\newcommand{\VhamiltonianPartB}[2]{\ensuremath{\hat{\text{V}}_\text{Part B}\left(#1,\,#2\right)}\xspace}
\newcommand{\VhamiltonianPartC}[2]{\ensuremath{\hat{\text{V}}_\text{Part C}\left(#1,\,#2\right)}\xspace}
\newcommand{\VhamiltonianOperatorPart}[4][]{\ensuremath{\hat{\text{F}}^{#1}_\text{#2}\left(#3,\,#4\right)}\xspace}
\newcommand{\VhamiltonianAnalyticalPart}[5][]{\ensuremath{\Lambda^{#1}_\text{#2}\left(#3,\,#4,\,#5\right)}\xspace}
\newcommand{\VhamiltonianAnalyticalPartIntegrated}[5][]{\ensuremath{\Pi^{#1}_\text{#2}\left(#3,\,#4,\,#5\right)}\xspace}

\newcommand{\timeOrderingOperator}{\ensuremath{\mathbb{T}}\xspace}
\newcommand{\timeEvolutionOperator}[1][]{\ensuremath{\hat{\text{P}}^{#1}(t)}\xspace}

\newcommand{\lsum}[1][l]{\ensuremath{\sum\limits_{#1}}\xspace}
\newcommand{\neighborsum}[3][]{\ensuremath{\sum\limits_{\left\langle#2, #3\right\rangle#1}}\xspace}
\newcommand{\neighborsumWSpin}[3]{\neighborsum[,\,#3]{#1}{#2}}
\newcommand{\fullneighborsum}[3][]{\ensuremath{\sum\limits_{\left[#2, #3\right]#1}}\xspace}

\newcommand{\epsl}[1][l]{\ensuremath{\varepsilon_{#1}}\xspace}
\newcommand{\ueps}[3]{\ensuremath{\tilde{U}_\text{#1}\left(#2,\,#3\right)}\xspace}

\newcommand{\picturePsi}[2][]{\ensuremath{\varPsi^{#1}{#2}}\xspace}
\newcommand{\psiN}[1][N]{\ensuremath{\varPsi_{#1}}\xspace}
\newcommand{\psiEta}[2]{\ensuremath{\varPsi_{#1}(#2)}\xspace}
\newcommand{\psiNStar}[1][N]{\ensuremath{\varPsi^\ast_{#1}}\xspace}
\newcommand{\brapsi}[2][]{\ensuremath{\bra{\picturePsi[#1]{#2}}}\xspace}
\newcommand{\ketpsi}[2][]{\ensuremath{\ket{\picturePsi[#1]{#2}}}\xspace}
\newcommand{\ketpsiEta}[2]{\ensuremath{\ket{\psiEta{#1}{#2}}}\xspace}
\newcommand{\brapsiof}[2][]{\brapsi[#1]{(#2)}}
\newcommand{\ketpsiof}[2][]{\ketpsi[#1]{(#2)}}
\newcommand{\psiOfT}[1][]{\ensuremath{\picturePsi[#1]{}(t)}\xspace}
\newcommand{\braketHelper}[2]{\ensuremath{\braket{#1}{#2}}\xspace}
\newcommand{\bracketHelper}[3]{\ensuremath{\matrixel{#1}{#2}{#3}}\xspace}

\newcommand{\braN}[1][N]{\ensuremath{\bra{#1}}\xspace}
\newcommand{\ketN}[1][N]{\ensuremath{\ket{#1}}\xspace}

\newcommand{\absSquare}[1]{\ensuremath{{\left\vert{#1}\right\vert}^2}\xspace}
\newcommand{\timeExpectationVal}[1]{\ensuremath{{\left\langle {#1}\right\rangle}_{(t)}}\xspace}
\newcommand{\timelessEtaExpectationVal}[1]{\ensuremath{{\left\langle {#1}\right\rangle}_{(\vec{\eta})}}\xspace}
\newcommand{\etaExpectationVal}[1]{\ensuremath{{\left\langle {#1}\right\rangle}_{(\vec{\eta}(t))}}\xspace}
\newcommand{\probabilityOf}[2]{\ensuremath{\text{P}(#1,\,#2)}\xspace}

\newcommand{\equivalentOperator}{\ensuremath{\hat{=}}\xspace}
\newcommand{\interactionTimeEvolution}[1]{\ensuremath{\left\{#1\right\}^{\interactionPicture}(t)}\xspace}
\newcommand{\heisenbergTimeEvolution}[1]{\ensuremath{\left\{#1\right\}^{\heisenbergPicture}(t)}\xspace}
\newcommand{\discreteGradient}{\ensuremath{\vec{\grad}_\text{disc}}\xspace}
\newcommand{\generaljvecop}[4]{\ensuremath{\vec{j}_{#2,\,#3}^{\,\,#1}#4}\xspace}
\newcommand{\currentvec}[3][]{\generaljvecop{#1}{#2}{#3}{}}
\newcommand{\currentvecOft}[3][]{\generaljvecop{#1}{#2}{#3}{(t)}}

\newcommand{\biggerNeighbor}[2]{\ensuremath{   {\vphantom{{\stackrel{+}{<}}} \left\langle#1 \smash{\stackrel{+}{<}}                         #2\right\rangle} }\xspace}
\newcommand{\smallerNeighbor}[2]{\ensuremath{  {\vphantom{{\stackrel{+}{>}}} \left\langle#1 \smash{\stackrel{+}{>}}                         #2\right\rangle} }\xspace}
\newcommand{\biggerNeighborX}[2]{\ensuremath{  {\vphantom{{\stackrel{+(\mathrm{x})}{<}}} \left\langle#1 \smash{\stackrel{+(\mathrm{x})}{<}} #2\right\rangle} }\xspace}
\newcommand{\smallerNeighborX}[2]{\ensuremath{ {\vphantom{{\stackrel{+(\mathrm{x})}{>}}} \left\langle#1 \smash{\stackrel{+(\mathrm{x})}{>}} #2\right\rangle} }\xspace}
\newcommand{\biggerNeighborY}[2]{\ensuremath{  {\vphantom{{\stackrel{+(\mathrm{y})}{<}}} \left\langle#1 \smash{\stackrel{+(\mathrm{y})}{<}} #2\right\rangle} }\xspace}
\newcommand{\smallerNeighborY}[2]{\ensuremath{ {\vphantom{{\stackrel{+(\mathrm{y})}{>}}} \left\langle#1 \smash{\stackrel{+(\mathrm{y})}{>}} #2\right\rangle} }\xspace}

\newcommand{\partialTrace}[2]{\ensuremath{\text{Tr}_{#1}\left(#2\right)}\xspace}

\DeclareRobustCommand{\Chi}{{\mathpalette\irchi\relax}}
\newcommand{\irchi}[2]{\raisebox{\depth}{$#1\chi$}} % inner command, used by \rchi

\let\oldtilde\tilde
\newcommand{\tildetilde}[1]{{\smash[b]{\overset{
    \mathchoice {\approx}% doesn't do anything in any of my formulas
                {\approx}% doesn't do anything in any of my formulas
                {   \vphantom{\raisebox{0pt}{\scaleobj{0.3}{\approx}}}
                    \smash{\raisebox{-3.0pt}{\scaleobj{1.1}{\approx}}}
                }
                {   \vphantom{\raisebox{0pt}{\scaleobj{0.3}{\approx}}}
                    \smash{\raisebox{-1.5pt}{\scaleobj{0.6}{\approx}}}
                }
    }{#1}}}}
\renewcommand{\tilde}[1]{{\smash[b]{\overset{
    \mathchoice {\sim} % doesn't do anything in any of my formulas
                {\sim} % doesn't do anything in any of my formulas
                {   \vphantom{\raisebox{0pt}{\scaleobj{0.22}{\sim}}}
                    \smash{\raisebox{-3.4pt}{\scaleobj{1.1}{\sim}}}
                }
                {   \vphantom{\raisebox{0pt}{\scaleobj{0.22}{\sim}}}
                    \smash{\raisebox{-1.9pt}{\scaleobj{0.6}{\sim}}}
                }
    }{#1}}}}

\newcommand{\energyOft}[1][t]{\ensuremath{E\left(#1\right)}\xspace}
\newcommand{\varianceOft}[1][t]{\ensuremath{\text{Var}_E(#1)}\xspace}
\newcommand{\matvec}[1]{\ensuremath{{\scalebox{0.65}{$\mathrlap{\overleftrightarrow{\phantom{#1}\vphantom{\scalebox{1.54}{$#1$}}}}$}{#1}}}\xspace}              % another file that holds predefined commands

%! hyphenation commonly used words can be spelled here to provide latex with the correct places to make line breaks
\hyphenation{Li-pid-mono-lage}
\hyphenation{Me-tro-po-lis-Al-go-rithm}

           % another file that holds predefined hyphenation helpers for unknown words


\begin{document}

% ! Bibliography, page numbering and Title setups
\thispagestyle{empty}                           % make sure title page is not numbered or anything else


\newcommand{\mail}{jonas.kell@student.uni-augsburg.de}



\begin{titlepage}
\makebox[\textwidth][c]{\includegraphics[width=0.5\textwidth]{logo_uni_augsburg.jpg}}
    
\color{dblue}

\begin{center}
    \vspace*{2cm}
    \Huge
    \textbf{\thetitle}

    \vspace*{1.5cm}
    \color{black}
    \textbf{Practical Training Report}

    \vspace*{1cm}
    \normalsize
    submitted by\\
    \LARGE
    \theauthor\\\vspace*{0.3cm}
    \normalsize
    on \thedate

    \vspace{1.8cm}
    \color{black}
    \emph{Augsburg University}\\
    \emph{Faculty of Mathematics, Natural Sciences, and Materials Engineering}\\
    \emph{Institute of Physics}\\
    \emph{Chair for theoretical Physics III}

    \vfill

    \begin{tabular}{rl}
        1$^\text{st}$ Corrector: &Prof. Dr. Markus Heyl\\
        2$^\text{nd}$ Corrector: &??\\ %TODO
    \end{tabular}
\end{center}

\end{titlepage}
                               % include title-page
\cleardoublepage                                % make sure, that if double-page is active, to reset the double page counter
\pagestyle{scrheadings}                         % puts current chapters title into the header in small gray font
\pagenumbering{roman}                           % number the pages of the table of contents in roman numerals
\renewcommand{\contentsname}{Table of Contents} % title of table of contents
\tableofcontents                                % table of contents
\noindent\\\\

\addsec*{List of Abbreviations}
\begin{tabular}[h]{p{3cm}|l}
	Abbreviation & Meaning\\        
	\midrule 
	TP III & Chair for theoretical Physics III\\
	MC & Monte Carlo\\
\end{tabular}\\\\

\newpage                                   % list of abbreviations, figures, etc
\cleardoublepage                                % make sure, that if double-page is active, to reset the double page counter
\pagenumbering{arabic}                          % number the pages of the main document in Arabic numerals

% After this, the redefinition of the "Kolumnentitel" takes place
\clearpairofpagestyles
\ihead{\leftmark}
\ohead{\Ifstr{\leftmark}{\rightmark}{}{\rightmark}}
\cfoot*{\pagemark}
% End of the "Kolumnentitel" redefinition


% ! Main Document Body
\chapter{Introduction}
\label{sec:introduction}
When performing time-evolution with a perturbative approach, one ought to watch out for resonances.
In the case of degenerate energy solutions, so-called \emph{secular terms} \cite{secularTermsPerturbation} might be introduced, if non-degenerate perturbation theory is used to find the solution to a degenerate problem. 
Often inaccuracies manifest as linear terms in a function where only oscillating terms would be expected.
The linear term constantly grows over time, guaranteeing the approximation is only valid for short times and then diverges further and further from the exact solution.

For the examined Hamiltonian the first occurrence of such terms can be spotted in the second order cumulant expansion of the Hamiltonian.
\autoref{eq:time-ordered-integral-part1} and \autoref{eq:time-ordered-integral-part2} specifically list extra solutions for the cases where the two energies are the same.
This resonance causes a linear term to appear in the expansion, while all other pre-factors so far were complex exponentials -- meaning they are oscillating.

The idea to combat this is to replace the perturbative parameters that were calculated analytically with variational parameters instead.
This has the motive to generate a physically motivated structure, but the fine details of the parameter weights are adjusted by optimization.
E.g. for the \emph{neural-network quantum state} \cite{neuralNetworkQuantumStates} approach, a neural network structure is used to represent a wave-function.
Reason for this being the hope that the exponentially complicated wave-function can be expressed as a \glqq higher truth\grqq{} only requiring few parameters to describe a seemingly complicated wave-function.
While it would be basically impossible to always find the few-parameter-solution, \emph{approximating} it close enough with optimization akin to machine-learning should be a suitable replacement.

The following section aims to achieve better approximations than the ones found by cumulant expansion by letting the parameters be optimized during the step to reflect the situation more closely.
\FloatBarrier

\chapter{Theory}
\label{sec:theory}

% Disclaimer, this is partly taken from report
Some parts of \fullref{sec:theory} were first written for a \emph{practical training report}\\ \filepath{\cite{selfDocument}}{/practical-training-latex} that was part of the works on the thesis-subject.
The calculations and text was supplemented and adapted for this final thesis. 
How the parts have been adapted, as well as the original author being the thesis author, can be verified in the version history of the repository \cite{selfDocument}.

    \section{Physical Basics}
    \label{sec:theory-physics}
    When performing time-evolution with a perturbative approach, one ought to watch out for resonances.
In the case of degenerate energy solutions, so-called \emph{secular terms} \cite{secularTermsPerturbation} might be introduced, if non-degenerate perturbation theory is used to find the solution to a degenerate problem. 
Often inaccuracies manifest as linear terms in a function where only oscillating terms would be expected.
The linear term constantly grows over time, guaranteeing the approximation is only valid for short times and then diverges further and further from the exact solution.

For the examined Hamiltonian the first occurrence of such terms can be spotted in the second order cumulant expansion of the Hamiltonian.
\autoref{eq:time-ordered-integral-part1} and \autoref{eq:time-ordered-integral-part2} specifically list extra solutions for the cases where the two energies are the same.
This resonance causes a linear term to appear in the expansion, while all other pre-factors so far were complex exponentials -- meaning they are oscillating.

The idea to combat this is to replace the perturbative parameters that were calculated analytically with variational parameters instead.
This has the motive to generate a physically motivated structure, but the fine details of the parameter weights are adjusted by optimization.
E.g. for the \emph{neural-network quantum state} \cite{neuralNetworkQuantumStates} approach, a neural network structure is used to represent a wave-function.
Reason for this being the hope that the exponentially complicated wave-function can be expressed as a \glqq higher truth\grqq{} only requiring few parameters to describe a seemingly complicated wave-function.
While it would be basically impossible to always find the few-parameter-solution, \emph{approximating} it close enough with optimization akin to machine-learning should be a suitable replacement.

The following section aims to achieve better approximations than the ones found by cumulant expansion by letting the parameters be optimized during the step to reflect the situation more closely.
    \FloatBarrier
        \subsection{Different Particle-Types, their Operators and Mappings}
        \label{sec:particles}
        To describe different particle-classes in the setting of many-body-physics on discrete lattices, the \emph{ladder operator formalism} is a useful tool. 
The inherit properties of the respective particle are guaranteed to be upheld by the \emph{(anti-) commutation relations}.

In some cases, transformations may be used to map problems formulated in one particle type onto equivalent formulations in a different particle type.
For the examples a discrete lattice of arbitrary enumeration -- which also takes care of potential additional spin-degrees -- is assumed.

\subsubsection*{Bosons}

\begin{equation}
    \label{eq:boson-commutators}
    \begin{split}
        \left[\bop{l}{},\,\bop{m}{}\right] &= 0\\
        \left[\bop{l}{\dagger},\,\bop{m}{\dagger}\right] &= 0\\
        \left[\bop{l}{},\,\bop{m}{\dagger}\right] &= \delta_{l,\,m}\\
    \end{split}
\end{equation}

The bosonic operators follow the commutation-relations described in \autoref{eq:boson-commutators}, with $\delta$ the \emph{Kronecker-Delta} \cite{schwablBookII}, which allows for occupation numbers of all natural numbers or $0$ per site.

\subsubsection*{Fermions}

For fermionic systems it is important, that their wave-function anti-symmetrizes. This results in the anti-commutation relations of  \autoref{eq:fermion-commutators}.

\begin{equation}
    \label{eq:fermion-commutators}
    \begin{split}
        \left\{\cop{l}{},\,\cop{m}{}\right\} &= 0\\
        \left\{\cop{l}{\dagger},\,\cop{m}{\dagger}\right\} &= 0\\
        \left\{\cop{l}{},\,\cop{m}{\dagger}\right\} &= \delta_{l,\,m}\\
    \end{split}
\end{equation}

From \autoref{eq:fermion-commutators}, one can derive that the operator combination $\cop{l}{\dagger}\cop{l}{}$, the \emph{number operator}, is \emph{idempotent}: 

\begin{equation}
    \label{eq:fermion-counting-op-idempotent}
    \left(\cop{l}{\dagger}\cop{l}{}\right)^n = 
    \cop{l}{\dagger}\cop{l}{} \qquad \forall n \in \mathbb{N}
\end{equation}

This is a property that is necessary for performing a specific calculation later and makes sense, once one realizes, this restricts the  fermionic occupations to the values $0$ and $1$.

\subsubsection*{Hard-Core Bosons}

While we require the number operator for the investigated particle type to be idempotent, it would be convenient to work with an operator that commutes and doesn't anti-commute like the fermions. 
In this regard, the hard-core bosons behave like a typical boson, but with extra restrictions that limit the occupation to either $0$ or $1$, meaning $\hop{l}{\dagger}\hop{l}{\dagger} = \hop{l}{}\hop{l}{} = 0$ \cite{hardCoreBosonsBasics}.
The required properties are fulfilled, if the particles obey the commutation-relations from \autoref{eq:hard-core-commutators} \cite{hardCoreBosonsFullCommutationRelation}, which in turn makes it possible to derive the idempotence-relation in \autoref{eq:hc-counting-op-idempotent}.

\begin{equation}
    \label{eq:hard-core-commutators}
    \begin{split}
        \left[\hop{l}{},\,\hop{m}{}\right] &= 0\\
        \left[\hop{l}{\dagger},\,\hop{m}{\dagger}\right] &= 0\\
        \left[\hop{l}{},\,\hop{m}{\dagger}\right] &= \left(1 - 2 \cdot \hop{m}{\dagger}\hop{m}{} \right) \cdot \delta_{l,\,m}
    \end{split}
\end{equation}

\begin{equation}
    \label{eq:hc-counting-op-idempotent}
    \left(\hop{l}{\dagger}\hop{l}{}\right)^n = 
    \hop{l}{\dagger}\hop{l}{} \qquad \forall n \in \mathbb{N}
\end{equation}

This additionally grants the helpful property $\hop{l}{\dagger}\hop{l}{} + \hop{l}{}\hop{l}{\dagger}  = 1$.

\subsubsection*{Spins / Pauli-Matrices}

Basic Hamiltonians, like the \emph{Ising-Model} are often expressed in the \emph{spin basis}. 
The operators that act on this basis are the \emph{Pauli matrices}, which have the useful property of being a basis for complex matrices.
Most \emph{quantum computing} formalisms are expressed in the spin basis and heavily rely on the Pauli matrices \cite{quantumBookPaulisAndBasics}:

\begin{equation}
    \label{eq:pauli-matrices}
    \begin{split}
        \pauli{0} = \one = \left(\begin{matrix}
            1& 0 \\
            0& 1
        \end{matrix}\right) 
        \qquad & \qquad
        \pauli{1} = \pauli{x} = \left(\begin{matrix}
            0& 1 \\
            1& 0
        \end{matrix}\right) \\
        \pauli{2} = \pauli{y} = \left(\begin{matrix}
            0& -i \\
            i& 0
        \end{matrix}\right) 
        \qquad & \qquad
        \pauli{3} = \pauli{z} = \left(\begin{matrix}
            1& 0 \\
            0& -1
        \end{matrix}\right)
    \end{split}
\end{equation}

For spin-particles on a lattice, the operators need to obey the relations in \autoref{eq:pauli-commutators}.

\begin{equation}
    \label{eq:pauli-commutators}
    \begin{split}
        \left\{\pauli[l]{$\alpha$},\,\pauli[l]{$\beta$}\right\} &= 2 \cdot \delta_{\alpha,\,\beta}\\
        \left[\pauli[l]{$\alpha$},\,\pauli[m]{$\beta$}\right] &= 2 \cdot i \cdot \varepsilon_{\alpha,\,\beta,\,\gamma} \cdot \delta_{l,\,m}\cdot \pauli[l]{$\gamma$}
    \end{split}
\end{equation}

With the complex unit $i$ and the \emph{Levi-Civita-Symbol} $\varepsilon_{\alpha,\,\beta,\,\gamma}$.

\subsubsection*{Transformation: Express Spins as Fermions}
For mapping between fermions and spins, there exists a well known transformation, the \emph{Jordan-Wigner-Transformation} \cite{jordanWignerBaseCase} (\autoref{eq:jordan-wigner-classical}).

The following transformations use \autoref{eq:jordan-wigner-p-factor}, where the equality follows from the fact that these operators are all different representations of the number operator.

\begin{equation}
    \label{eq:jordan-wigner-p-factor}
\left(-1\right)^{\sum\limits_{k=0}^{l-1} \left( \pauli[k]{z} +1 \right)/2 } = 
\left(-1\right)^{\sum\limits_{k=0}^{l-1}\cop{k}{\dagger}\cop{k}{} } =
\left(-1\right)^{\sum\limits_{k=0}^{l-1}\hop{k}{\dagger}\hop{k}{} } = 
\text{P}(l)
\end{equation}

With the easily verifiable property $\text{P}(l)\cdot \text{P}(l) = 1$. Also $\text{P}(l)$ commutes with all $\cop{m}{}$, $\cop{m}{\dagger}$, $\hop{m}{}$ and $\hop{m}{\dagger}$ if $l\leq m$ and furthermore $\left[\text{P}(l),\, \text{P}(m)\right] = 0$, always.

\begin{equation}
    \label{eq:jordan-wigner-classical}
    \begin{split}
        \cop{l}{\dagger} &= \text{P}(l) \cdot \frac{\pauli[l]{x} + i\cdot \pauli[l]{y}}{2} = \text{P}(l) \cdot \pauli[l]{+}\\
        \cop{l}{} &= \text{P}(l) \cdot \frac{\pauli[l]{x} - i\cdot \pauli[l]{y}}{2} = \text{P}(l) \cdot \pauli[l]{-}\\
        \cop{l}{\dagger}\cop{l}{} &= \left( \pauli[l]{z} +1 \right)/2\\
        &\Longleftrightarrow \\
        \pauli[l]{x} &= \left(\cop{l}{} + \cop{l}{\dagger}\right) \cdot \text{P}(l) \\
        \pauli[l]{y} &= i \cdot \left(\cop{l}{} - \cop{l}{\dagger}\right) \cdot \text{P}(l)\\
        \pauli[l]{z} &= 2 \cdot \cop{l}{\dagger}\cop{l}{} -1
    \end{split}
\end{equation}

For performing commutation-relation calculations, a helper program was written in conjunction with this thesis and related submissions.
The program is called \emph{Math-Manipulator} \cite{selfMathManipulator} and the transformations in this chapter have been validated with it in a supplementary repository:\\ \filepath{\cite{selfMathManipulatorCalculations}}{/jordan-wigner-transformation}.


\subsubsection*{Transformation: Express Fermions as Hard-Core Bosons}
The tensor-network library \emph{ITensor} provides in their documentation an example for a mapping between fermions and hard-core bosons, used to simplify their calculations \cite{itensorFermionizationLibrary}.
The transformations may be called \emph{fermionization/bosonization} or respectively also Jordan-Wigner-Transformation.

\autoref{eq:jordan-wigner-hcb-fermions} presents the transformation:

\begin{equation}
    \label{eq:jordan-wigner-hcb-fermions}
    \begin{split}
        \cop{l}{\dagger} &= \text{P}(l) \cdot \hop{l}{\dagger}\\
        \cop{l}{} &= \text{P}(l) \cdot \hop{l}{}\\
        &\Longleftrightarrow \\
        \hop{l}{\dagger} &= \text{P}(l) \cdot \cop{l}{\dagger}\\
        \hop{l}{} &= \text{P}(l) \cdot \cop{l}{}\\
    \end{split}
\end{equation}

\subsubsection*{Transformation: Express Spins as Hard-Core Bosons}
Finally this confirms the transformation that will be required later: expressing a measurement in the spin-basis in terms of hard-core bosonic operators.

\autoref{eq:jordan-wigner-spin-hcb} can be directly obtained by plugging \autoref{eq:jordan-wigner-hcb-fermions} into \autoref{eq:jordan-wigner-classical}.

\begin{equation}
    \label{eq:jordan-wigner-spin-hcb}
    \begin{split}
        \pauli[l]{x} &= \hop{l}{} + \hop{l}{\dagger} \\
        \pauli[l]{y} &= i \cdot \left(\hop{l}{} - \hop{l}{\dagger}\right) \\
        \pauli[l]{z} &= 2 \cdot \hop{l}{\dagger}\hop{l}{} -1
    \end{split}
\end{equation}

This connection shows, that spin-$\frac{1}{2}$ systems and hard-core bosonic systems can be used to investigate the same problems in different notations.
\autoref{eq:spin-hcb-state-mapping} shows, which states are equivalent for these two notations.

\begin{equation}
    \label{eq:spin-hcb-state-mapping}
    \begin{split}
        \ketN[\down] &= \left(\begin{matrix}
            0 \\
            1
        \end{matrix}\right) = \ketN[0] \\
        \pauli[]{+}\ketN[\down] = \ketN[\up] &= \left(\begin{matrix}
            1 \\
            0
        \end{matrix}\right) = \ketN[1] = \hop{}{\dagger}\ketN[0] \\
    \end{split}
\end{equation}

This exact mapping is important to conserve sign-consistency across all calculations. 
E.g. swapping the $\ketN[0] = \ketN[\down]$ correspondence to $\ketN[0] = \ketN[\up]$, would introduce a hidden minus sign in observables depending on $\pauli[]{y}$, because of the implicit order of basis states in the definition of $\pauli[]{y}$.
        \FloatBarrier

        \subsection{Hubbard Model and discussed Hamiltonian}
        \label{sec:theory-hubbard-hamiltonian}
        The system in question is described by a \emph{Hubbard Model Hamiltonian}, that includes an electrical field. It is shown in \autoref{eq:main-hamiltonian}.


\begin{equation}
    \label{eq:main-hamiltonian}
    \pictureHamiltonian = \HzeroHamiltonian + \Vhamiltonian
\end{equation}

\begin{equation}
    \label{eq:main-hamiltonian-h0}
    \HzeroHamiltonian = U \cdot \lsum \nop{l}{\up}\nop{l}{\down} + \lsum[l,\,\sigma] \underbrace{\left(\vec{E} \cdot \vec{r_l}\right)}_{\text{\normalsize \epsl}} \nop{l}{\sigma}
\end{equation}

\begin{equation}
    \label{eq:main-hamiltonian-perturbation}
    \Vhamiltonian =  - J \cdot \neighborsumWSpin{l}{m}{\sigma} \left(\withspinhcop{l}{\sigma}{\dagger} \withspinhcop{m}{\sigma}{} + \withspinhcop{m}{\sigma}{\dagger} \withspinhcop{l}{\sigma}{} \right)
\end{equation}

$\vec{E}$ describes the vector of said electrical field and $\vec{r_l}$ the position of the site with index $l$.
The number operators $\nop{l}{\sigma} = \withspinhcop{l}{\sigma}{\dagger}\withspinhcop{l}{\sigma}{}$ measure the occupation on site $l$ with the respective spin $\sigma$. 

$U$, $J$ and \epsl[] are constants that describe the interaction strength. 
They all have the unit of energy. 
The scalar Product that is defined as \epsl[] can be evaluated, based on the system geometry. 
In this case, the system described by a regular square pattern that can be seen in \autoref{fig:geometry-of-square-system}. 
For such a system, the energy difference that is acquired from hopping from site $l$ to $m$ ($\Delta E_{l \rightarrow m}$) (provided, $l$ and $m$ are nearest neighbors) is described by \autoref{eq:energy-difference-hopping}. By assuming a default value for \epsl[0] one gets the relation from \autoref{eq:epsl} for \epsl{} (using the telescoping sum over nearest-neighbors-hopping along the path from $\vec{r_0}$ to $\vec{r_l}$). 

\begin{figure}[htbp]
    \centering
            
\def\getFromArray#1{%
    \ifcase#1 0\or
    1\or
    2\or
    \dots\or
    (M-1)\or
    M\or
    M+1\or
    M+2\or
    \dots\or
    M+(M-1)\or
    2\cdot M\or
    2\cdot M+1\or
    2\cdot M+2\or
    \dots\or
    2\cdot M+(M-1)\or
    \dots\or
    \dots\or
    \dots\or
    \dots\or
    \dots\or
    (M-1)\cdot M\or
    M^2-M+1\or
    M^2-M+2\or
    \dots\or
    M^2 -1\or
    \else\fi
}


\begin{tikzpicture}[scale=0.38]
    \def\m{5} % Change this value to adjust the grid size (m x m)
    \def\side{6} % Change this value to adjust the square side length
    \def\labelsize{\side/6} % Adjust label size
    
    % Draw dots and labels with custom text using a loop
    \foreach \x in {0,...,\numexpr\m-1\relax} {
        \foreach \y in {0,...,\numexpr\m-1\relax} {
            \pgfmathtruncatemacro{\index}{(\m-\y-1) * \m + \x}
            \node[scale=\labelsize] at (\x*\side + \side/2,\y*\side + \side/2) {$\getFromArray{\index}$};
            \fill[color=dblue] (\x*\side + \side/2,\y*\side + \side/2 + \side/4) circle (0.4);
        }
    }

    % Store coordinates for reference
    \coordinate (top) at (current bounding box.north west);

    % Draw horizontal black line
    \draw[line width=2pt,overlay,lightgray,dashed] (top) ++ (0,-0.55*\side) -- ++(\m*\side,0);
        
    % Draw green arrow pointing downwards
    \draw[green, ->, overlay, line width=2pt,>=stealth] (top) ++ (1.3*\side,-0.55*\side) -- ++(3*\side,-3*\side) node[pos=0.05,right] {$\varphi$};

    % Draw angle circle
    \draw[overlay, green, line width=1.5pt] (top) ++ (1.3*\side,-0.55*\side) +(0:0.5*\side) arc (0:-45:0.5*\side);
\end{tikzpicture}
    \vspace{0.8cm}
    \caption{Graphical representation of how the examined square system is laid out. The sites are labeled with the index $l$ they later can be identified by. Each of the sites has 4 nearest neighbors, except the ones on the borders, that have 3 or only 2 nearest neighbors (as the system is not periodic). A system of side-length $M$ with size $N$ is depicted, which means that it has $M^2 = N$ sites.
    In green, the Electrical field vector $\vec{E}$ is depicted, it is parametrized by the field strength $E$ and the angle $\varphi$.}
    \label{fig:geometry-of-square-system}
\end{figure}

\begin{equation}
    \label{eq:energy-difference-hopping}
    \begin{split}
        \Delta E_{l \rightarrow m} &=  \epsl[m] - \epsl[l] = \vec{E} \cdot \left(\vec{r_m}-\vec{r_l}\right)\\
        &= E \cdot \left[\cos(\varphi)\cdot \left(m \% M-l \% M\right) + \sin(\varphi)\cdot \left(\floor[\bigg]{\frac{m}{M}}-\floor[\bigg]{\frac{l}{M}}\right)\right]
    \end{split}
\end{equation}
\vspace{0.5cm} % equations are not well enough separated otherwise
\begin{equation}
    \label{eq:epsl}
    \begin{split}
        \epsl[0] = 0 \Rightarrow
        \epsl[l] &= \epsl[l] - \epsl[0] = \Delta E_{0 \rightarrow x} + \dots + \Delta E_{y \rightarrow l}\\
        &= E \cdot \left[\cos(\varphi)\cdot l \% M + \sin(\varphi)\cdot \floor[\bigg]{\frac{l}{M}}\right]
    \end{split}
\end{equation}

During measurements in the \nameref{sec:numerical-experiments} section, a 1-dimensional version of this is used that will be labeled \emph{chain}. 
It will be obtained by cutting the first row of the respective square geometry and is consistent with basically all mathematical derivations performed for the \emph{square} case, because of the nearest-neighbor notation that is independent of the underlying geometry. Only in the derivation of \fullref{sec:theory-observables-current} appear differences based on the dimensionality, that however do not affect the end result.
The reason for the truncation is to reduce thy system size for the verification of convergence towards exact comparisons.
While such a modification necessarily affects the physical properties, it is fine for verification of convergence.
Also interesting physical - system-size independent - effects only manifest in experiments with lattice site numbers far beyond any possible exact diagonalization anyway.

The Hubbard Model \cite{hubbardModelOriginalDerivation} is a spin-dependent quantum-mechanical model. 
Because of this, the Hamiltonian includes terms for both the spin directions $\up$ and $\down$ (most of the time denoted by a summation over $\sigma$).

In principle, the operators $\withspinhcop[]{l}{\sigma}{(\dagger)}$ can be any type of quantum-mechanical ladder operator, like the aforementioned bosonic or fermionic ones.
In this application, the decision was made to use \emph{hard-core bosons} \cite[]{hardCoreBosonsBasics}, for they happen to obey multiple useful properties.

Typically for bosons, the Hubbard Model is transformed into the \emph{Bose-Hubbard Model} \cite{boseHubbardModelOriginalDerivation}.
However hard-core bosons have the added property, that the can only have the occupation numbers of $0$ and $1$ per-spin-and-site in the investigated energy range.
This means they behave basically like fermions, only that their wave-function symmetrizes instead of anti-symmetrizes \cite{schwablBook}.
Because of the restriction, one term of the Bose-Hubbard Model drops out and the studied model in equations \ref{eq:main-hamiltonian}, \ref{eq:main-hamiltonian-h0} and \ref{eq:main-hamiltonian-perturbation} is produced.
Also the typical Bose-Hubbard Model is not spin-dependent, while this property is kept here.
Hard-core bosons on their own have relevant properties that makes them worth studying. 
Experimental research on hard-core bosons can give insight into e.g. \emph{ultra-cold gasses} and \emph{Bose-Einstein condensation} \cite{hardCoreBosonsBasics}, in this case it is important for the following derivations, that the operators on different sites commute, the number operator is idempotent (\ref{eq:hc-counting-op-idempotent}) and can be mapped to spin-basis notation (\ref{eq:jordan-wigner-spin-hcb}).

When adjusting the commutation-relations in \autoref{eq:hard-core-commutators} for spin-dependency, one obtains \autoref{eq:commutation-relations}.

\begin{equation}
    \label{eq:commutation-relations}
    \begin{split}
        \left[\withspinhcop[]{l}{\sigma}{},\,\withspinhcop[]{l'}{\sigma'}{}\right] &= 
        \left[\withspinhcop[]{l}{\sigma}{\dagger},\,\withspinhcop[]{l'}{\sigma'}{\dagger}\right] = 0\\
        \left[\withspinhcop[]{l}{\sigma}{},\,\withspinhcop[]{l'}{\sigma'}{\dagger}\right] &= 
        \left(1 - 2 \cdot \withspinhcop[]{l}{\sigma}{\dagger}\withspinhcop[]{l}{\sigma}{} \right) \cdot \delta_{l,\,l'}\cdot\delta_{\sigma,\sigma'}
    \end{split}
\end{equation}

For ease of notation, the two spin degrees of freedom may be described by the alternate naming scheme provided in \autoref{eq:alt-naming-scheme}.
The related code implementation, comments and some calculations in linked material \cite{selfCode} use the letters c and d, as that used to be the first designation wile the thesis was developed.

\begin{equation}
    \label{eq:alt-naming-scheme}
    \withspinhcop{l}{\up}{(\dagger)} \leftrightarrow \hop{l}{(\dagger)}
    \qquad\qquad
    \withspinhcop{l}{\down}{(\dagger)} \leftrightarrow \dop{l}{(\dagger)}
\end{equation}

For simplification purposes, the two spin directions are assumed to have the same coupling constants ($U$, $J$ and \epsl).
This means, all later results \emph{must} be symmetrical in terms of the spin directions $\up$ and $\down$.


Translated into the alternative notation, \autoref{eq:main-hamiltonian-h0} reads like \autoref{eq:main-hamiltonian-h0-alt-notation} and \autoref{eq:main-hamiltonian-perturbation} reads like \autoref{eq:main-hamiltonian-perturbation-alt-notation}.

\begin{equation}
    \label{eq:main-hamiltonian-h0-alt-notation}
    \HzeroHamiltonian = U \cdot \lsum \hop[]{l}{\dagger}\hop[]{l}{} \dop[]{l}{\dagger}\dop[]{l}{} + \lsum\epsl \hop[]{l}{\dagger}\hop[]{l}{} + \lsum\epsl \dop[]{l}{\dagger}\dop[]{l}{}
\end{equation}

\begin{equation}
    \label{eq:main-hamiltonian-perturbation-alt-notation}
    \Vhamiltonian =  - J \cdot \neighborsum{l}{m} \left(\hop[]{l}{\dagger}\hop[]{m}{} +\hop[]{m}{\dagger}\hop[]{l}{} + \dop[]{l}{\dagger}\dop[]{m}{} +\dop[]{m}{\dagger}\dop[]{l}{} \right)
\end{equation}

Several of the previous equations feature the \emph{nearest neighbor} summing notation with pointy brackets $\langle\rangle$.
This is supposed to express, that in the case of $\neighborsum[]{l}{m}$ for each possible $l$, $m$ will take on the indices of all neighbor sites of $l$.
However, if a pair $l=x,\,m=y$ is part of the sum, then the inverse $l=y,\,m=x$ will not be.
To be more in-line with the code implementation, \autoref{eq:fullneigborsum} gives an alternative notation using $[]$.
In such sums, each pair $l=x,\,m=y$ and $l=y,\,m=x$ will be included in the sum for each nearest-neighbor-relation $x,\,y$.

\begin{equation}
    \label{eq:fullneigborsum}
    \begin{split}
        &\neighborsum[]{l}{m} \left( F(l,m)\, \hop[]{l}{\dagger}\hop[]{m}{} + F(m,l)\, \hop[]{m}{\dagger}\hop[]{l}{} \right) \stackrel{\stackrel{\text{\scriptsize rename-swap}}{\text{$l$ and $m$}}}{=} \\
        &\neighborsum[]{l}{m} F(l,m)\, \hop[]{l}{\dagger}\hop[]{m}{} + \neighborsum[]{m}{l} F(l,m)\, \hop[]{l}{\dagger}\hop[]{m}{} = \\
        &\fullneighborsum[]{l}{m} F(l,m)\, \hop[]{l}{\dagger}\hop[]{m}{} \stackrel{\text{! caution !}}{\neq}
        2 \cdot \neighborsum[]{l}{m} F(l,m)\, \hop[]{l}{\dagger}\hop[]{m}{}
    \end{split}
\end{equation}

With that convention, it is possible to rewrite \autoref{eq:main-hamiltonian-perturbation-alt-notation} to \autoref{eq:main-hamiltonian-perturbation-full-sum}.

\begin{equation}
    \label{eq:main-hamiltonian-perturbation-full-sum}
    \begin{split}
        \Vhamiltonian =  - J \cdot \fullneighborsum{l}{m} \left(\hop[]{l}{\dagger}\hop[]{m}{} + \dop[]{l}{\dagger}\dop[]{m}{} \right)
    \end{split}
\end{equation}

        \FloatBarrier

        
        \subsection{Interaction Picture and Time Evolution}
        \label{sec:theory-interaction-picture}
        Solving the quantum-mechanical many-body time-evolution problem in general requires at least a time-complexity of \bigo{2^n}.
For some special cases it is possible to solve it analytically.
E.g. for one-dimensional systems, the Hubbard model can be solved analytically even with the addition of an electrical field \cite{exactSolutionExampleHubbardEvenElectricalField}.

For arbitrary geometry in two or even more dimensions, this is no longer possible.
In that case, approximations are required to lower the computational complexity to a reasonable level to obtain the required measurements.
Here, the calculation is performed in the \emph{Interaction Picture} (relevant equations from \cite{schwablBook}), along the lines of the calculation performed in \cite{variationalClassicalNetworksPaper}.

The equations provided so far show operators from the \emph{Schrödinger Picture}. 
For this section, the operators and states will be labeled whether they are Schrödinger- or Interaction Picture. All un-labeled elements are assumed to be Schrödinger Picture.

A general state in the Schrödinger Picture can be expanded in a basis as shown in \autoref{eq:base-expansion-state}. It is time-evolved using the full Hamiltonian in the Schrödinger Picture (\autoref{eq:time-evolution-schroedinger}, all units so that $\hbar = 1$).

\begin{equation}
    \label{eq:base-expansion-state}
    \ketpsi[\schroedingerPicture]{} = \ketpsiof[\schroedingerPicture]{t=0} = 
    \lsum[N] \ketN{}\underbrace{\braketHelper{N}{\picturePsi[\schroedingerPicture]{}}}_{\text{\normalsize\psiN}} = \lsum[N] \psiN \ketN{}
\end{equation}

\begin{equation}
    \label{eq:time-evolution-schroedinger}
    \ketpsiof[\schroedingerPicture]{t} = e^{-i \pictureHamiltonian[\schroedingerPicture] t} \ketpsiof[\schroedingerPicture]{t=0} = e^{-i \pictureHamiltonian[\schroedingerPicture] t} \ketpsi[\schroedingerPicture]{}
\end{equation}

In the Interaction Picture, operators have a time evolution, that depends only on the \HzeroHamiltonian[\schroedingerPicture]{}-part, shown in \autoref{eq:time-evolution-operator-interaction} (assuming the operators are time-independent in the Schrödinger Picture). 

\begin{equation}
    \label{eq:time-evolution-operator-interaction}
    \begin{split}
        \AopOfT[\interactionPicture] &= e^{i \HzeroHamiltonian[\schroedingerPicture]{} t} \Aop[\schroedingerPicture] e^{-i \HzeroHamiltonian[\schroedingerPicture]{} t}
        \quad \Rightarrow \quad 
        \HzeroHamiltonian[\schroedingerPicture]{} = \HzeroHamiltonian[\interactionPicture]{}= \HzeroHamiltonian[]{}\\
        \text{especially:} \qquad \VhamiltonianOfT[\interactionPicture] &=  e^{i \HzeroHamiltonian[\schroedingerPicture]{} t} \Vhamiltonian[\schroedingerPicture] e^{-i \HzeroHamiltonian[\schroedingerPicture]{} t}
    \end{split}
\end{equation}

The time-evolution of a general state in the Interaction Picture on the other hand is obtained as shown in \autoref{eq:time-evolution-interaction}.

\begin{equation}
    \label{eq:time-evolution-interaction}
    \ketpsiof[\interactionPicture]{t} = e^{i \HzeroHamiltonian{} t} \ketpsiof[\schroedingerPicture]{t} \stackrel{\ref{eq:time-evolution-schroedinger}}{=} e^{i \HzeroHamiltonian{} t} e^{-i \pictureHamiltonian[\schroedingerPicture]{} t} \ketpsi[\schroedingerPicture]{} 
    \quad \Rightarrow \quad 
    \ketpsiof[\schroedingerPicture]{t} = e^{-i \HzeroHamiltonian{} t} \ketpsiof[\interactionPicture]{t}
\end{equation}

Differentiating \autoref{eq:time-evolution-interaction} and substituting the special case in \autoref{eq:time-evolution-operator-interaction} along the way, one gets \autoref{eq:equation-of-motion-interaction-picture}.

\begin{equation}
    \label{eq:equation-of-motion-interaction-picture}
    \begin{split}
        \difft{\ketpsiof[\interactionPicture]{t}} &\stackrel{\ref{eq:time-evolution-interaction}}{=}
        \difft{} \left(e^{i \HzeroHamiltonian{} t} e^{-i \pictureHamiltonian[\schroedingerPicture]{} t} \right)\ketpsi[\schroedingerPicture]{} \\
        &\stackrel{\phantom{\ref{eq:time-evolution-interaction}}}{=} -i \cdot e^{i \HzeroHamiltonian{} t}\left( \pictureHamiltonian[\schroedingerPicture]{} - \HzeroHamiltonian{} \right) e^{-i \pictureHamiltonian[\schroedingerPicture]{} t} \ketpsi[\schroedingerPicture]{}\\
        &\stackrel{\ref{eq:main-hamiltonian}}{=} -i \cdot e^{i \HzeroHamiltonian{} t} \Vhamiltonian[\schroedingerPicture] \cdot \one \cdot e^{-i \pictureHamiltonian[\schroedingerPicture]{} t} \ketpsi[\schroedingerPicture]{}\\
        &\stackrel{\phantom{\ref{eq:main-hamiltonian}}}{=} -i \cdot e^{i \HzeroHamiltonian{} t} \Vhamiltonian[\schroedingerPicture] e^{-i \HzeroHamiltonian{} t}e^{i \HzeroHamiltonian{} t} e^{-i \pictureHamiltonian[\schroedingerPicture]{} t} \ketpsi[\schroedingerPicture]{}\\
        &\stackrel{\ref{eq:time-evolution-interaction}}{=} -i \cdot e^{i \HzeroHamiltonian{} t} \Vhamiltonian[\schroedingerPicture] e^{-i \HzeroHamiltonian{} t}\ketpsiof[\interactionPicture]{t}\\
        &\stackrel{\ref{eq:time-evolution-operator-interaction}}{=} -i \cdot \VhamiltonianOfT[\interactionPicture] \ketpsiof[\interactionPicture]{t}
    \end{split}
\end{equation}

The equation of motion in \autoref{eq:equation-of-motion-interaction-picture} is solved by the ansatz in \autoref{eq:solve-equation-of-motion} with the  \emph{time-ordering-operator} \timeOrderingOperator and the \emph{time-evolution-operator} \timeEvolutionOperator[\interactionPicture].

\begin{equation}
    \label{eq:solve-equation-of-motion}
    \begin{split}
        \ketpsiof[\interactionPicture]{t} &= \timeOrderingOperator \left\lbrace e^{-i\int\limits_0^t \text{d} t' \VhamiltonianOf[\interactionPicture]{t'}}\right\rbrace \ketpsiof[\interactionPicture]{0}\\
        &= \timeEvolutionOperator[\interactionPicture] \ketpsiof[\interactionPicture]{0} \stackrel{\ref{eq:time-evolution-interaction}}{=} \timeEvolutionOperator[\interactionPicture] \ketpsiof[\schroedingerPicture]{0}
    \end{split}
\end{equation}

Plugging this in, one finally obtains the general time-evolution for the state in the Schrödinger Picture in \autoref{eq:time-evolution-target}.
This requires the \emph{base-energies} $E_0(N)$, where $e^{-i \HzeroHamiltonian{} t}\ketN{} = e^{-i E_0(N) t}\ketN{}$.

\begin{equation}
    \label{eq:time-evolution-target}
    \begin{split}
        \ketpsiof[\schroedingerPicture]{t} &\stackrel{\ref{eq:time-evolution-interaction}}{=} e^{-i \HzeroHamiltonian{} t} \ketpsiof[\interactionPicture]{t} 
        \stackrel{\ref{eq:solve-equation-of-motion}}{=}
        e^{-i \HzeroHamiltonian{} t} \timeEvolutionOperator[\interactionPicture] \ketpsi[\schroedingerPicture]{} \\
        &\stackrel{\phantom{\ref{eq:solve-equation-of-motion}}}{=}
        e^{-i \HzeroHamiltonian{} t} \cdot \one \cdot  \timeEvolutionOperator[\interactionPicture] \ketpsi[\schroedingerPicture]{}\cdot \one\\
        &\stackrel{\phantom{\ref{eq:solve-equation-of-motion}}}{=}
        e^{-i \HzeroHamiltonian{} t} \lsum[N] \ketN{} \bracketHelper{N}{\timeEvolutionOperator[\interactionPicture]}{\picturePsi[\schroedingerPicture]{}} \cdot \frac{\braketHelper{N}{\picturePsi[\schroedingerPicture]{}}}{\braketHelper{N}{\picturePsi[\schroedingerPicture]{}}}\\
        &\stackrel{\ref{eq:base-expansion-state}}{=}
        \lsum[N] \underbrace{e^{-i \HzeroHamiltonian{} t}}_{\text{\normalsize{$e^{-i E_0 t}$}}}  \underbrace{\frac{\bracketHelper{N}{\timeEvolutionOperator[\interactionPicture]}{\picturePsi[\schroedingerPicture]{}}}{\braketHelper{N}{\picturePsi[\schroedingerPicture]{}}}}_{\text{\normalsize{$e^{\HNOft}$}}}  \psiN \ketN{}\\
        &\stackrel{\phantom{\ref{eq:base-expansion-state}}}{=}
        \lsum[N] e^{-i E_0(N) t}  e^{\HNOft} \psiN \ketN{} = 
        \lsum[N] e^{\HeffOft} \psiN \ketN{}
    \end{split}
\end{equation}

Up to this point, the calculation is still exact. Now the effective Hamiltonian \HNOft could be approximated in different ways. Typically $e^{\HNOft}$ is \emph{Taylor-expanded}, but in this context this approximation would converge quite poorly. 
To combat this and require a smaller order of expansion, one could employ the \emph{cumulant-expansion} \cite{cumulantExpansionOriginalDerivation} to expand the object in terms of its exponent \HNOft and not the whole exponential $e^{\HNOft}$.
The first and second order cumulants for the cumulant expansion are provided in \autoref{eq:cumulant-expansion-coefficients} \cite{variationalClassicalNetworksPaper} with the \emph{cumulant-average} in respect to the state $\ketN$, $\left\langle\cdot\right\rangle_{\text{c}(N)}$.

\begin{equation}
    \label{eq:cumulant-expansion-coefficients}
    \begin{split}
        \left\langle\Aop\right\rangle_{\text{c}(N)} &= \frac{ \bracketHelper{N}{\Aop}{\picturePsi[\schroedingerPicture]{}}}{\braketHelper{N}{\picturePsi[\schroedingerPicture]{}}}\\
        \left\langle\Aop\Bop\right\rangle_{\text{c}(N)} &= \frac{\bracketHelper{N}{\Aop\Bop}{\picturePsi[\schroedingerPicture]{}}}{\braketHelper{N}{\picturePsi[\schroedingerPicture]{}}} -  \frac{\bracketHelper{N}{\Aop}{\picturePsi[\schroedingerPicture]{}} \cdot \bracketHelper{N}{\Bop}{\picturePsi[\schroedingerPicture]{}}}{\braketHelper{N}{\picturePsi[\schroedingerPicture]{}}^2}\\
    \end{split}
\end{equation}

Plugging this into the rearranged result from \autoref{eq:time-evolution-target} following the derivation in \cite{variationalClassicalNetworksPaper} and \cite{isingDynamicsWithClassicalNetworks}, an expansion for \HNOft in \autoref{eq:cumulant-expansion} is obtained.

\begin{equation}
    \label{eq:cumulant-expansion}
    \begin{split}
        \braketHelper{N}{\picturePsi[\schroedingerPicture]{}} \cdot e^{\HNOft} &\stackrel{\ref{eq:time-evolution-target}}{=} \bracketHelper{N}{\timeEvolutionOperator[\interactionPicture]}{\picturePsi[\schroedingerPicture]{}}
        \stackrel{\ref{eq:solve-equation-of-motion}}{=}
        \bracketHelper{N}{
            \timeOrderingOperator e^{-i\int\limits_0^t \text{d} t' \VhamiltonianOf[\interactionPicture]{t'}}
        }{\picturePsi[\schroedingerPicture]{}}\\
        &\stackrel{\text{\cite{variationalClassicalNetworksPaper}}}{\Longleftrightarrow}\\
        \HNOft &= \sum\limits_{v=1}^\infty \frac{\left(-i\right)^v}{v!} \int\limits_0^t\mathrm{d}t_1 \int\limits_0^t\mathrm{d}t_2 \,\,\dots \int\limits_0^t\mathrm{d}t_v \left\langle \timeOrderingOperator\, \VhamiltonianOf[\interactionPicture]{t_1}\VhamiltonianOf[\interactionPicture]{t_2} \cdots \VhamiltonianOf[\interactionPicture]{t_v} \right\rangle_{\text{c}(N)}\\
        &= -i \int\limits_0^t \mathrm{d}t_1 \frac{\bracketHelper{N}{\VhamiltonianOf[\interactionPicture]{t_1}}{\picturePsi[\schroedingerPicture]{}}}{\braketHelper{N}{\picturePsi[\schroedingerPicture]{}}}\\ 
        &- \frac{1}{2} \int\limits_0^t\mathrm{d}t_1 \int\limits_0^t\mathrm{d}t_2
        \left(
        \frac{\bracketHelper{N}{\timeOrderingOperator
        \VhamiltonianOf[\interactionPicture]{t_1}\VhamiltonianOf[\interactionPicture]{t_2}
        }{\picturePsi[\schroedingerPicture]{}}}{\braketHelper{N}{\picturePsi[\schroedingerPicture]{}}} -  \frac{\bracketHelper{N}{\VhamiltonianOf[\interactionPicture]{t_1}}{\picturePsi[\schroedingerPicture]{}} \cdot \bracketHelper{N}{\VhamiltonianOf[\interactionPicture]{t_2}}{\picturePsi[\schroedingerPicture]{}}}{\braketHelper{N}{\picturePsi[\schroedingerPicture]{}}^2}
        \right) 
        \hspace{-1.3cm} % put the equation number where there is space
        \\
        &+ \cdots
    \end{split}
\end{equation}

As $\VhamiltonianOf[\interactionPicture]{t'} \propto J$, if $J$ is a small perturbation (in comparison to $U$ and $\varepsilon$) this quickly converges, even when taking only few terms into account.
        \FloatBarrier
        
        \subsection{Operators in the Interaction Picture}
        \label{sec:theory-approximation-evaluation}
        As \autoref{eq:cumulant-expansion} reveals, by this process the problem was converted into finding $\VhamiltonianOf[\interactionPicture]{t'}$ to be able to calculate the first two orders of the expansion.

Operators in the Interaction Picture can be derived from their Schrödinger Picture variants like previously shown in \autoref{eq:time-evolution-operator-interaction}. 
With this, one can derive a useful expression for their \emph{equation of motion} (\autoref{eq:schroedinger-op-to-interaction-op}).

\begin{equation}
    \label{eq:schroedinger-op-to-interaction-op}
    \begin{split}
        \difft{}\AopOfT[\interactionPicture] &\stackrel{\phantom{\ref{eq:time-evolution-operator-interaction}}}{=} i \left[\HzeroHamiltonian[\schroedingerPicture], \AopOfT[\interactionPicture]\right]
        %
        \stackrel{\ref{eq:time-evolution-operator-interaction}}{=}
        i \left[\HzeroHamiltonian[\schroedingerPicture],\,e^{i \HzeroHamiltonian[\schroedingerPicture]{} t} \Aop[\schroedingerPicture] e^{-i \HzeroHamiltonian[\schroedingerPicture]{} t}\right]\\
        %
        &\stackrel{\phantom{\ref{eq:time-evolution-operator-interaction}}}{=}
        i e^{i \HzeroHamiltonian[\schroedingerPicture]{} t}\left[\HzeroHamiltonian[\schroedingerPicture],\, \Aop[\schroedingerPicture]\right]e^{-i \HzeroHamiltonian[\schroedingerPicture]{} t}
        %
        \stackrel{\phantom{\ref{eq:time-evolution-operator-interaction}}}{=}
        i \interactionTimeEvolution{\left[\HzeroHamiltonian[\schroedingerPicture],\, \Aop[\schroedingerPicture]\right]}
    \end{split}
\end{equation}

\autoref{eq:schroedinger-op-to-interaction-op} also defines the operator \interactionTimeEvolution{\ast}, which simply is $e^{i \HzeroHamiltonian[\schroedingerPicture]{} t} \left\{\ast \right\} e^{-i \HzeroHamiltonian[\schroedingerPicture]{} t}$.
This section uses $(t)$ for indicating that operators in the Interaction Picture have a time dependence. 
This is superficial and only for better readability, as in fact $\hopOfT[\interactionPicture]{m}{\dagger} = \hop[\interactionPicture]{m}{\dagger} = \interactionTimeEvolution{\hop[\schroedingerPicture]{m}{\dagger}}$.

If now $\left[\HzeroHamiltonian[\schroedingerPicture],\, \Aop[\schroedingerPicture]\right]$ now is a function of $\Aop[\schroedingerPicture]$, one gets a \emph{differential equation} that can be solved to obtain $\AopOfT[\interactionPicture]$.

For the repetitive calculation of the objects of type $\left[\HzeroHamiltonian[\schroedingerPicture],\, \Aop[\schroedingerPicture]\right]$ for various \Aop[\schroedingerPicture], the tool \emph{Math-Manipulator} was specifically developed. 
Like the transformations in \ref{sec:particles}, calculations in the following section have been validated or calculated from scratch in this tool.
The relevant files can be found in \filepath{\cite{selfMathManipulatorCalculations}}{/calculate-V}, and \filepath{\cite{selfMathManipulatorCalculations}}{/time-evolution}.
The equations \ref{eq:stuff-from-math-manipulator-dd} and \ref{eq:stuff-from-math-manipulator-cm} list exemplary results that were computed with the tool.

\begin{equation}
    \label{eq:stuff-from-math-manipulator-dd}
    \difft{}\left(\dop[\interactionPicture]{m}{\dagger}\dop[\interactionPicture]{m}{}\right)(t)
    = i \interactionTimeEvolution{\left[\HzeroHamiltonian[],\,\dop[\schroedingerPicture]{m}{\dagger}\dop[\schroedingerPicture]{m}{}\right]} \stackrel{\text{MM}}{=} 0 \quad \Rightarrow \quad \left(\dop[\interactionPicture]{m}{\dagger}\dop[\interactionPicture]{m}{}\right)(t) = \dop[\schroedingerPicture]{m}{\dagger}\dop[\schroedingerPicture]{m}{}
\end{equation}

\begin{equation}
    \label{eq:stuff-from-math-manipulator-cm}
    \begin{split}
        \difft{}\hopOfT[\interactionPicture]{m}{\dagger}
        &\stackrel{\phantom{\ref{eq:time-evolution-operator-interaction}}}{=} i \interactionTimeEvolution{\left[\HzeroHamiltonian[]{},\,\hop[\schroedingerPicture]{m}{\dagger}\right]} \stackrel{\text{MM}}{=} 
        i \interactionTimeEvolution{ \left(\epsl[m] + U \dop[\schroedingerPicture]{m}{\dagger}\dop[\schroedingerPicture]{m}{}\right)\hop[\schroedingerPicture]{m}{}}\\
        &\stackrel{\phantom{\ref{eq:time-evolution-operator-interaction}}}{=}i e^{i \HzeroHamiltonian[\schroedingerPicture]{} t} \left(\epsl[m] + U \dop[\schroedingerPicture]{m}{\dagger}\dop[\schroedingerPicture]{m}{}\right)\hop[\schroedingerPicture]{m}{}e^{-i \HzeroHamiltonian[\schroedingerPicture]{} t}\\
        &\stackrel{\phantom{\ref{eq:time-evolution-operator-interaction}}}{=}i \left(\epsl[m] + U e^{i \HzeroHamiltonian[\schroedingerPicture]{} t}\dop[\schroedingerPicture]{m}{\dagger}\dop[\schroedingerPicture]{m}{}e^{-i \HzeroHamiltonian[\schroedingerPicture]{} t}\right)e^{i \HzeroHamiltonian[\schroedingerPicture]{} t}\hop[\schroedingerPicture]{m}{}e^{-i \HzeroHamiltonian[\schroedingerPicture]{} t}\\
        &\stackrel{\ref{eq:time-evolution-operator-interaction}}{=}
        i \left(\epsl[m] + U \left(\dop[\interactionPicture]{m}{\dagger}\dop[\interactionPicture]{m}{}\right)(t)\right) \hopOfT[\interactionPicture]{m}{\dagger}\\
    \end{split}
\end{equation}

With this and the idempotence-relation \autoref{eq:hc-counting-op-idempotent} one can first derive \autoref{eq:operator-out-from-exponent} and finally use an exponential function as a natural ansatz to solve \autoref{eq:stuff-from-math-manipulator-cm} with \autoref{eq:the-final-time-dependence-cop}.

\begin{equation}
    \label{eq:operator-out-from-exponent}
    \begin{split}
        e^{a \cdot \withspinhcop[\schroedingerPicture]{l}{\sigma}{\dagger}\withspinhcop[\schroedingerPicture]{l}{\sigma}{}} &= \sum\limits_{m=0}^\infty \frac{a^m \left(\withspinhcop[\schroedingerPicture]{l}{\sigma}{\dagger}\withspinhcop[\schroedingerPicture]{l}{\sigma}{}\right)^m}{m!}
        \stackrel{\ref{eq:hc-counting-op-idempotent}}{=} 1 + \left[\sum\limits_{m=1}^\infty \frac{a^m}{m!}\right]\cdot \left(\withspinhcop[\schroedingerPicture]{l}{\sigma}{\dagger}\withspinhcop[\schroedingerPicture]{l}{\sigma}{}\right)\\
        &= 1 + \left[\sum\limits_{m=0}^\infty \frac{a^m}{m!}-1\right]\cdot \left(\withspinhcop[\schroedingerPicture]{l}{\sigma}{\dagger}\withspinhcop[\schroedingerPicture]{l}{\sigma}{}\right)
         = 1 + \left(e^a-1\right)\cdot \withspinhcop[\schroedingerPicture]{l}{\sigma}{\dagger}\withspinhcop[\schroedingerPicture]{l}{\sigma}{}
    \end{split}
\end{equation}


\begin{equation}
    \label{eq:the-final-time-dependence-cop}
    \begin{split}
        \stackrel{\ref{eq:stuff-from-math-manipulator-cm},\, \ref{eq:stuff-from-math-manipulator-dd}}{\Longrightarrow}
        \hopOfT[\interactionPicture]{m}{\dagger}
        &\stackrel{\phantom{\ref{eq:stuff-from-math-manipulator-dd}}}{=}
        e^{i \cdot \epsl[m]\cdot  t + i \cdot U \left(\dop[\interactionPicture]{m}{\dagger}\dop[\interactionPicture]{m}{}\right)(t) \cdot  t}  \hop[\schroedingerPicture]{m}{\dagger} 
        \stackrel{\ref{eq:stuff-from-math-manipulator-dd}}{=}
        e^{i \cdot \epsl[m]\cdot  t + i \cdot U \cdot \dop[\schroedingerPicture]{m}{\dagger}\dop[\schroedingerPicture]{m}{} \cdot  t}  \hop[\schroedingerPicture]{m}{\dagger} \\
        &\stackrel{\ref{eq:operator-out-from-exponent}}{=}
        e^{i \cdot \epsl[m]\cdot  t}\left(1 + \left(e^{i  \cdot U  \cdot  t}-1\right) \dop[\schroedingerPicture]{m}{\dagger}\dop[\schroedingerPicture]{m}{}\right)  \hop[\schroedingerPicture]{m}{\dagger}
    \end{split}
\end{equation}

From analogous calculations follows:

\begin{equation*}
    \begin{split}
        \hopOfT[\interactionPicture]{m}{\dagger}  &= e^{i \cdot \epsl[m]\cdot  t}\left(1 + \left(e^{i  \cdot U  \cdot  t}-1\right) \dop[\schroedingerPicture]{m}{\dagger}\dop[\schroedingerPicture]{m}{}\right)  \hop[\schroedingerPicture]{m}{\dagger}\\
        \hopOfT[\interactionPicture]{m}{}  &= e^{-i \cdot \epsl[m]\cdot  t}\left(1 + \left(e^{-i  \cdot U  \cdot  t}-1\right) \dop[\schroedingerPicture]{m}{\dagger}\dop[\schroedingerPicture]{m}{}\right)  \hop[\schroedingerPicture]{m}{}\\
        \dopOfT[\interactionPicture]{m}{\dagger}  &= e^{i \cdot \epsl[m]\cdot  t}\left(1 + \left(e^{i  \cdot U  \cdot  t}-1\right) \hop[\schroedingerPicture]{m}{\dagger}\hop[\schroedingerPicture]{m}{}\right)  \dop[\schroedingerPicture]{m}{\dagger}\\
        \dopOfT[\interactionPicture]{m}{}  &= e^{-i \cdot \epsl[m]\cdot  t}\left(1 + \left(e^{-i  \cdot U  \cdot  t}-1\right) \hop[\schroedingerPicture]{m}{\dagger}\hop[\schroedingerPicture]{m}{}\right)  \dop[\schroedingerPicture]{m}{}\\
    \end{split}
\end{equation*}

Injecting these derivations into \autoref{eq:main-hamiltonian-perturbation-full-sum}, finally \VhamiltonianOfT[\interactionPicture] can be derived in \autoref{eq:interaction-picture-v-ham-full}.

\begin{equation}
    \label{eq:interaction-picture-v-ham-full}
    \begin{split}
        \VhamiltonianOfT[\interactionPicture] &
        \stackrel{\phantom{\text{MM}}}{=}
        \interactionTimeEvolution{\Vhamiltonian[\schroedingerPicture]} \stackrel{\ref{eq:main-hamiltonian-perturbation-full-sum}}{=}
        - J \cdot \fullneighborsum{l}{m}  \interactionTimeEvolution{
              \left(\hop[\schroedingerPicture]{l}{\dagger}\hop[\schroedingerPicture]{m}{} + \dop[\schroedingerPicture]{l}{\dagger}\dop[\schroedingerPicture]{m}{} \right)
        } \\
        &\stackrel{\phantom{\text{MM}}}{=}
        -J \cdot \fullneighborsum{l}{m}\left(\hopOfT[\interactionPicture]{l}{\dagger}\hopOfT[\interactionPicture]{m}{} + \dopOfT[\interactionPicture]{l}{\dagger}\dopOfT[\interactionPicture]{m}{} \right)\\
        &\stackrel{\text{MM}}{=}
        %
        -J \cdot \fullneighborsum{l}{m} \left[
            \VhamiltonianAnalyticalPart{A}{l}{m}{t} \cdot \VhamiltonianOperatorPart{A}{l}{m} + 
            \VhamiltonianAnalyticalPart{B}{l}{m}{t} \cdot \VhamiltonianOperatorPart{B}{l}{m} + 
            \VhamiltonianAnalyticalPart{C}{l}{m}{t} \cdot \VhamiltonianOperatorPart{C}{l}{m} 
        \right]
    \end{split}
\end{equation}

It is possible to express these operator-strings as single hoppings, decorated with number operators. One receives \autoref{eq:interaction-picture-v-ham-parts}, that is symmetrical in \up and \down (with $\overline{\up} = \down$ and  $\overline{\down} = \up$).

\begin{equation}
    \label{eq:interaction-picture-v-ham-parts}
    \begin{split}
        \VhamiltonianAnalyticalPart{A}{l}{m}{t} \stackrel{\text{MM}}{=} e^{i\cdot \left(\epsl-\epsl[m]\right)\cdot t} \qquad
        \VhamiltonianOperatorPart{A}{l}{m} &\stackrel{\text{MM}}{=} 
        \lsum[\sigma \in \left\{\up,\,\down\right\}]
        \withspinhcop[\schroedingerPicture]{l}{\sigma}{\dagger}\withspinhcop[\schroedingerPicture]{m}{\sigma}{}
        \left(
            1+
            2 \cdot \nop[\schroedingerPicture]{l}{\overline{\sigma}}\nop[\schroedingerPicture]{m}{\overline{\sigma}}
            - \nop[\schroedingerPicture]{l}{\overline{\sigma}}
            - \nop[\schroedingerPicture]{m}{\overline{\sigma}}
        \right)
        \\
        \VhamiltonianAnalyticalPart{B}{l}{m}{t} \stackrel{\text{MM}}{=} e^{i\cdot \left(\epsl-\epsl[m] + U\right)\cdot t} \qquad
        \VhamiltonianOperatorPart{B}{l}{m} &\stackrel{\text{MM}}{=} 
        \lsum[\sigma \in \left\{\up,\,\down\right\}]        
        \withspinhcop[\schroedingerPicture]{l}{\sigma}{\dagger}\withspinhcop[\schroedingerPicture]{m}{\sigma}{}
        \left(
            \nop[\schroedingerPicture]{l}{\overline{\sigma}}
            - \nop[\schroedingerPicture]{l}{\overline{\sigma}}\nop[\schroedingerPicture]{m}{\overline{\sigma}}
        \right)
        \\
        \VhamiltonianAnalyticalPart{C}{l}{m}{t} \stackrel{\text{MM}}{=} e^{i\cdot \left(\epsl-\epsl[m] - U\right)\cdot t} \qquad
        \VhamiltonianOperatorPart{C}{l}{m} &\stackrel{\text{MM}}{=} 
        \lsum[\sigma \in \left\{\up,\,\down\right\}]
        \withspinhcop[\schroedingerPicture]{l}{\sigma}{\dagger}\withspinhcop[\schroedingerPicture]{m}{\sigma}{}
        \left(
            \nop[\schroedingerPicture]{m}{\overline{\sigma}}
            - \nop[\schroedingerPicture]{m}{\overline{\sigma}}\nop[\schroedingerPicture]{l}{\overline{\sigma}}
        \right)
        \\
    \end{split}
\end{equation}

The expressions for the dressed operators in \autoref{eq:interaction-picture-v-ham-parts} can be rearranged, to reflect more their physical effects.
\autoref{eq:interaction-picture-v-ham-parts-simplified} can be obtained and verified by comparing the measurements in a truth table (the \nop[\schroedingerPicture]{}{} measure either a $0$ or $1$ on a state, without modifying it). % TODO fix extra comma

\begin{equation}
    \label{eq:interaction-picture-v-ham-parts-simplified}
    \begin{split}
        \VhamiltonianOperatorPart{A}{l}{m} &\stackrel{\text{MM}}{=} 
        \lsum[\sigma \in \left\{\up,\,\down\right\}]
        \withspinhcop[\schroedingerPicture]{l}{\sigma}{\dagger}\withspinhcop[\schroedingerPicture]{m}{\sigma}{}
        \left(\nop[\schroedingerPicture]{l}{\overline{\sigma}}
            \equivalentOperator 
            \nop[\schroedingerPicture]{m}{\overline{\sigma}}
        \right)
        \\
        \VhamiltonianOperatorPart{B}{l}{m} &\stackrel{\text{MM}}{=} 
        \lsum[\sigma \in \left\{\up,\,\down\right\}]        
        \withspinhcop[\schroedingerPicture]{l}{\sigma}{\dagger}\withspinhcop[\schroedingerPicture]{m}{\sigma}{}
        \nop[\schroedingerPicture]{l}{\overline{\sigma}}
        \left(
            1 - \nop[\schroedingerPicture]{m}{\overline{\sigma}}
        \right)
        \\
        \VhamiltonianOperatorPart{C}{l}{m} &\stackrel{\text{MM}}{=} 
        \lsum[\sigma \in \left\{\up,\,\down\right\}]
        \withspinhcop[\schroedingerPicture]{l}{\sigma}{\dagger}\withspinhcop[\schroedingerPicture]{m}{\sigma}{}
        \nop[\schroedingerPicture]{m}{\overline{\sigma}}
        \left(
           1 - \nop[\schroedingerPicture]{l}{\overline{\sigma}}
        \right)
        \\
    \end{split}
\end{equation}

The \equivalentOperator here means \glqq $1$ if the measurements of the operators are equivalent, $0$ otherwise\grqq.

\subsubsection*{Base Energy}

In \autoref{eq:time-evolution-target}, the effective Hamiltonian $\HeffOft = -i E_0(N) t + \HNOft$ was defined.
This construct will be necessary in the following sections. 
$E_0(N)$ is luckily calculated quite swiftly, as one can see in \autoref{eq:base-energy} (with $n_{l,\,\sigma}$ - the occupation-number in $\ketN$, not the operator).

\begin{equation}
    \label{eq:base-energy}
    \begin{split}
        E_0(N) &\stackrel{\phantom{\ref{eq:main-hamiltonian-h0}}}{=} \frac{\bracketHelper{N}{\HzeroHamiltonian[]}{N}}{\braketHelper{N}{N}} \\
        %
        &\stackrel{\ref{eq:main-hamiltonian-h0}}{=} \bracketHelper{N}{U \cdot \lsum \nop{l}{\up}\nop{l}{\down}}{N} + \bracketHelper{N}{\lsum[l,\,\sigma] \epsl \nop{l}{\sigma}}{N}\\
        %
        &\stackrel{\phantom{\ref{eq:main-hamiltonian-h0}}}{=} U \cdot \lsum n_{l,\,\up}n_{l,\,\down} + \lsum[l,\,\sigma] \epsl n_{l,\,\sigma}
    \end{split}
\end{equation}

\subsubsection*{Calculate the first order of \HNOft}

The last missing element for $\HeffOft = -i E_0(N) t + \HNOft$ now is \HNOft.
At least for the first order, the integration in \autoref{eq:cumulant-expansion} is easily doable, with now knowing the form of \VhamiltonianOf[\interactionPicture]{t}.

The result of the integration can be seen in \autoref{eq:hn-integrated-first-order}.

\begin{equation}
    \label{eq:hn-integrated-first-order}
    \begin{split}
        \HNOftOrder{1} &\stackrel{\ref{eq:cumulant-expansion}}{=} -i \int\limits_0^t \mathrm{d}t' \frac{\bracketHelper{N}{\VhamiltonianOf[\interactionPicture]{t'}}{\picturePsi[\schroedingerPicture]{}}}{\braketHelper{N}{\picturePsi[\schroedingerPicture]{}}}\\
        %
        &\stackrel{\ref{eq:base-expansion-state}}{=}
        -i \frac{1}{\psiN} \int\limits_0^t \mathrm{d}t' \lsum[K] \bracketHelper{N}{\VhamiltonianOf[\interactionPicture]{t'}}{K} \psiN[K]\\
        &\stackrel{\ref{eq:interaction-picture-v-ham-full}}{=}
        i \cdot J \lsum[K] \fullneighborsum{l}{m} \frac{\psiN[K]}{\psiN} \int\limits_0^t \mathrm{d}t' 
        \left[
            e^{i\cdot \left(\epsl-\epsl[m]\right)\cdot t'} \cdot 
            \bracketHelper{N}{
                \VhamiltonianOperatorPart{A}{l}{m} 
            }{K}
            + 
            \right.\\
            &\qquad
            \left.
            e^{i\cdot \left(\epsl-\epsl[m] + U\right)\cdot t'} \cdot 
            \bracketHelper{N}{
                \VhamiltonianOperatorPart{B}{l}{m}
            }{K}
            + 
            e^{i\cdot \left(\epsl-\epsl[m] - U\right)\cdot t'} \cdot 
            \bracketHelper{N}{
                \VhamiltonianOperatorPart{C}{l}{m} 
            }{K}
        \right]
        \\
        %
        &\stackrel{\phantom{\ref{eq:interaction-picture-v-ham-full}}}{=}
        J \lsum[K] \fullneighborsum{l}{m} \frac{\psiN[K]}{\psiN}
        \left[
            \frac{e^{i\cdot \left(\epsl-\epsl[m]\right)\cdot t}-1}{\epsl-\epsl[m]} \cdot 
            \bracketHelper{N}{
                \VhamiltonianOperatorPart{A}{l}{m} 
            }{K}
            + 
            \right.\\
            &\qquad
            \left.
            \frac{
                e^{i\cdot \left(\epsl-\epsl[m] + U\right)\cdot t}-1
            }{\epsl-\epsl[m] + U}
             \cdot 
            \bracketHelper{N}{
                \VhamiltonianOperatorPart{B}{l}{m}
            }{K}
            + 
            \frac{
                e^{i\cdot \left(\epsl-\epsl[m] - U\right)\cdot t}-1
            }{\epsl-\epsl[m] - U}
             \cdot 
            \bracketHelper{N}{
                \VhamiltonianOperatorPart{C}{l}{m} 
            }{K}
        \right]\\
        %
        &\stackrel{\phantom{\ref{eq:interaction-picture-v-ham-full}}}{=}
        J \lsum[K] \fullneighborsum{l}{m} \frac{\psiN[K]}{\psiN}
        \left[
            \VhamiltonianAnalyticalPartIntegrated{A}{l}{m}{t} \cdot 
            \bracketHelper{N}{
                \VhamiltonianOperatorPart{A}{l}{m} 
            }{K}
            + 
            \right.\\
            &\qquad
            \left.
            \VhamiltonianAnalyticalPartIntegrated{B}{l}{m}{t}
             \cdot 
            \bracketHelper{N}{
                \VhamiltonianOperatorPart{B}{l}{m}
            }{K}
            + 
            \VhamiltonianAnalyticalPartIntegrated{C}{l}{m}{t}
             \cdot 
            \bracketHelper{N}{
                \VhamiltonianOperatorPart{C}{l}{m} 
            }{K}
        \right]
    \end{split}
\end{equation}

In the second to last step, the $i$ is moved into the integral.

While it seems that $\lsum[K] \fullneighborsum{l}{m}$ is even a larger summation than the previously problematic ones, because the operators $\VhamiltonianOperatorPart{A}{l}{m}$, $\VhamiltonianOperatorPart{B}{l}{m}$ and $\VhamiltonianOperatorPart{C}{l}{m}$ have interaction range limited to nearest neighbors, actually the summation is really sparse and most entries are $0$.
This summation can be efficiently evaluated in \bigo{\text{\#}(\text{sites}) \cdot \text{\#}(\text{nearest neighbors})}.
How this could be done, can be looked up in the implementation \filepath{\cite{selfCode}}{/computation-scripts/hamiltonian.py} in the specialization class \emph{FirstOrderCanonicalHamiltonian}.

\subsubsection*{Calculate the second order of \HNOft}

For the calculation of the second order, the time-ordering-operator is relevant again. 
It however is not necessary for the right side of the integral, as the two integrals can be fully localized to one factor each.
For the left side, it is necessary to integrate different cases separately.

\vspace{-0.2cm} % TODO check this page-breaks properly in the end

\begin{equation}
    \label{eq:hn-integrated-second-order}
    \begin{split}
        \HNOftOrder{2} &\stackrel{\ref{eq:cumulant-expansion}}{=} - \frac{1}{2} \int\limits_0^t\mathrm{d}t_1 \int\limits_0^t\mathrm{d}t_2
        \left(
        \frac{\bracketHelper{N}{\timeOrderingOperator
        \VhamiltonianOf[\interactionPicture]{t_1}\VhamiltonianOf[\interactionPicture]{t_2}
        }{\picturePsi[\schroedingerPicture]{}}}{\braketHelper{N}{\picturePsi[\schroedingerPicture]{}}} -  \frac{\bracketHelper{N}{\VhamiltonianOf[\interactionPicture]{t_1}}{\picturePsi[\schroedingerPicture]{}} \cdot \bracketHelper{N}{\VhamiltonianOf[\interactionPicture]{t_2}}{\picturePsi[\schroedingerPicture]{}}}{\braketHelper{N}{\picturePsi[\schroedingerPicture]{}}^2}
        \right)
        \hspace{-1.3cm} % put the equation number where there is space
        \\
        %
        &\stackrel{\ref{eq:base-expansion-state}}{=}
        - \frac{1}{2} 
        \left[
        \lsum[K] \frac{\psiN[K]}{\psiN} \int\limits_0^t \int\limits_0^t        \bracketHelper{N}{\timeOrderingOperator
        \VhamiltonianOf[\interactionPicture]{t_1}\VhamiltonianOf[\interactionPicture]{t_2}
        }{K}            \mathrm{d}t_1\mathrm{d}t_2 \right. \\
        &\stackrel{\phantom{\ref{eq:base-expansion-state}}}{-} 
        \left.
        \lsum[M]\lsum[L]
        \frac{\psiN[M]\psiN[L]}{\psiN{}^2}
        \int\limits_0^t \bracketHelper{N}{\VhamiltonianOf[\interactionPicture]{t_1}}{M}  \mathrm{d}t_1
        \cdot
        \int\limits_0^t \bracketHelper{N}{\VhamiltonianOf[\interactionPicture]{t_2}}{L}  \mathrm{d}t_2 \right]\\
        %
        &\stackrel{\ref{eq:interaction-picture-v-ham-full}}{=}
        - \frac{J^2}{2} 
        \left[
        \lsum[K]\lsum[X,Y \in \left\{\text{A}, \text{B}, \text{C}\right\}] \fullneighborsum{l}{m}\fullneighborsum{j}{k}
        \frac{\psiN[K]}{\psiN} \right. \\&\quad
        \int\limits_0^t \int\limits_0^t        \bracketHelper{N}{\timeOrderingOperator
        \VhamiltonianAnalyticalPart{$X$}{l}{m}{t_1} \VhamiltonianOperatorPart[t_1]{$X$}{l}{m}
        \cdot 
        \VhamiltonianAnalyticalPart{$Y$}{j}{k}{t_2} \VhamiltonianOperatorPart[t_2]{$Y$}{j}{k}
        }{K}            \mathrm{d}t_1\mathrm{d}t_2 \\
        &\stackrel{\phantom{\ref{eq:interaction-picture-v-ham-full}}}{-} 
        \lsum[M]\lsum[L]\lsum[V,W \in \left\{\text{A}, \text{B}, \text{C}\right\}]
        \fullneighborsum{o}{p}\fullneighborsum{q}{r}
        \frac{\psiN[M]\psiN[L]}{\psiN{}^2}\\&\quad
        \left.
        \int\limits_0^t \bracketHelper{N}{\VhamiltonianAnalyticalPart{$V$}{o}{p}{t_1} \cdot \VhamiltonianOperatorPart{$V$}{o}{p}}{M}  \mathrm{d}t_1
        \cdot
        \int\limits_0^t \bracketHelper{N}{\VhamiltonianAnalyticalPart{$W$}{q}{r}{t_2} \cdot \VhamiltonianOperatorPart{$W$}{q}{r}}{L}  \mathrm{d}t_2 \right]\\
        %
        &=
        - \frac{J^2}{2} \left[ 
        \lsum[K]\lsum[X,Y \in \left\{\text{A}, \text{B}, \text{C}\right\}] \fullneighborsum{l}{m}\fullneighborsum{j}{k}
        \frac{\psiN[K]}{\psiN} 
        \right.
        \\&\quad
        \int\limits_0^t \int\limits_0^t  
        \VhamiltonianAnalyticalPart{$X$}{l}{m}{t_1}
        \cdot 
        \VhamiltonianAnalyticalPart{$Y$}{j}{k}{t_2}
        \bracketHelper{N}{\timeOrderingOperator
            \VhamiltonianOperatorPart[t_1]{$X$}{l}{m}
        \cdot 
            \VhamiltonianOperatorPart[t_2]{$Y$}{j}{k}
        }{K}            \mathrm{d}t_1\mathrm{d}t_2 \\
        & - 
        \lsum[M]\lsum[L]\lsum[V,W \in \left\{\text{A}, \text{B}, \text{C}\right\}]
        \fullneighborsum{o}{p}\fullneighborsum{q}{r}
        \frac{\psiN[M]\psiN[L]}{\psiN{}^2}\\&\quad
        \left.
            (-1) \cdot \VhamiltonianAnalyticalPartIntegrated{$V$}{o}{p}{t} 
            \cdot
            \VhamiltonianAnalyticalPartIntegrated{$W$}{q}{r}{t}
            \bracketHelper{N}{\VhamiltonianOperatorPart{$V$}{o}{p}}{M}
            \bracketHelper{N}{\VhamiltonianOperatorPart{$W$}{q}{r}}{L}
            \vphantom{
                \lsum[M]
            }
        \right]
    \end{split}
\end{equation}

Notice in the second to last step, the \VhamiltonianOperatorPart[t_\ast]{$\ast$}{\ast}{\ast} is indicated with a $t_\ast$ to keep its time-ordering association, even after writing out the \VhamiltonianOf[\interactionPicture]{t_\ast}.
In the last step an additional -1 appears, because of a $-1 = i \cdot i$ that was introduced, to allow for integration of \VhamiltonianAnalyticalPart{$\ast$}{\ast}{\ast}{t_\ast} like in \autoref{eq:hn-integrated-first-order}.

\begin{equation}
    \label{eq:u-eps-substitution}
    \begin{split}
        \epsl-\epsl[m] &= \ueps{A}{l}{m}\\
        \epsl-\epsl[m]+U &= \ueps{B}{l}{m}\\
        \epsl-\epsl[m]-U &= \ueps{C}{l}{m}\\
    \end{split}
\end{equation}

Inverting the steps performed in \cite{dissectTimeOrderingIntegrals}, one can easily see how to split integrals with the time-ordering-operator for evaluation in \autoref{eq:split-integral-time-ordering}.

\begin{equation}
    \label{eq:split-integral-time-ordering}
    \begin{split}
        &\int\limits_0^t \mathrm{d}t_1 \int\limits_0^{t_1} \mathrm{d}t_2 A(t_1) B(t_2)  = \int\limits_0^t \mathrm{d}t_1 \int\limits_0^{t_1} \mathrm{d}t_2  \timeOrderingOperator A(t_1) B(t_2) \\
        &\int\limits_0^t \mathrm{d}t_1 \int\limits_{t_1}^t \mathrm{d}t_2  B(t_2) A(t_1)  = \int\limits_0^t \mathrm{d}t_1 \int\limits_{t_1}^t \mathrm{d}t_2  \timeOrderingOperator A(t_1) B(t_2) \\\\
        %
        &\int\limits_0^t \mathrm{d}t_1 \left( \int\limits_0^{t_1} \mathrm{d}t_2 \timeOrderingOperator A(t_1) B(t_2) +  \int\limits_{t_1}^t \mathrm{d}t_2 \timeOrderingOperator A(t_1) B(t_2)\right) =
        \int\limits_0^t \mathrm{d}t_1\int\limits_0^t \mathrm{d}t_2 \timeOrderingOperator A(t_1) B(t_2)
    \end{split}
\end{equation}

What is left, is to evaluate the integral that needs to respect the time-ordering.

\begin{equation}
    \label{eq:time-ordered-integral}
    \begin{split}
        &\int\limits_0^t \int\limits_0^t  
        \VhamiltonianAnalyticalPart{$X$}{l}{m}{t_1}
        \cdot 
        \VhamiltonianAnalyticalPart{$Y$}{j}{k}{t_2}
        \bracketHelper{N}{\timeOrderingOperator
            \VhamiltonianOperatorPart[t_1]{$X$}{l}{m}
        \cdot 
            \VhamiltonianOperatorPart[t_2]{$Y$}{j}{k}
        }{K}            \mathrm{d}t_1\mathrm{d}t_2\\
        \stackrel{\ref{eq:u-eps-substitution},\,\ref{eq:interaction-picture-v-ham-parts}}{=}
        %
        &\int\limits_0^t \int\limits_0^t  
        e^{
            i \cdot \ueps{$X$}{l}{m} \cdot t_1
        }
        \cdot 
        e^{
            i \cdot \ueps{$Y$}{j}{k} \cdot t_2
        }
        \bracketHelper{N}{\timeOrderingOperator
            \VhamiltonianOperatorPart[t_1]{$X$}{l}{m}
        \cdot 
            \VhamiltonianOperatorPart[t_2]{$Y$}{j}{k}
        }{K}            \mathrm{d}t_1\mathrm{d}t_2\\
        \stackrel{\ref{eq:split-integral-time-ordering}}{=}
        %
        &\int\limits_0^t  \int\limits_0^{t_1}
        e^{
            i \cdot \ueps{$X$}{l}{m} \cdot t_1
        }
        e^{
            i \cdot \ueps{$Y$}{j}{k} \cdot t_2
        }
        \mathrm{d}t_2
        \mathrm{d}t_1
        \cdot 
        \bracketHelper{N}{
            \VhamiltonianOperatorPart[]{$X$}{l}{m}
            \VhamiltonianOperatorPart[]{$Y$}{j}{k}
        }{K}           \\ +
        &\int\limits_0^t  \int\limits_{t_1}^t
        e^{
            i \cdot \ueps{$X$}{l}{m} \cdot t_1
        }
        e^{
            i \cdot \ueps{$Y$}{j}{k} \cdot t_2
        }
        \mathrm{d}t_2
        \mathrm{d}t_1
        \cdot 
        \bracketHelper{N}{
            \VhamiltonianOperatorPart[]{$Y$}{j}{k}
            \VhamiltonianOperatorPart[]{$X$}{l}{m}
        }{K} 
    \end{split}
\end{equation}

\begin{equation}
    \label{eq:time-ordered-integral-part1}
    \begin{split}
        \Gamma_1 &\left(X,\,l,\,m,\,Y,\,j,\,k\right) = \\=&\int\limits_0^t  \int\limits_0^{t_1}
        e^{
            i \cdot \ueps{$X$}{l}{m} \cdot t_1
        }
        e^{
            i \cdot \ueps{$Y$}{j}{k} \cdot t_2
        }
        \mathrm{d}t_2
        \mathrm{d}t_1 = 
        %
        \int\limits_0^t 
        e^{
            i \cdot \ueps{$X$}{l}{m} \cdot t_1
        }
        \int\limits_0^{t_1}
        e^{
            i \cdot \ueps{$Y$}{j}{k} \cdot t_2
        }
        \mathrm{d}t_2
        \mathrm{d}t_1 \\=
        %
        &\int\limits_0^t 
        e^{
            i \cdot \ueps{$X$}{l}{m} \cdot t_1
        }
        \frac{
            e^{
                i \cdot \ueps{$Y$}{j}{k} \cdot t
            }
            -1
        }{
            i \cdot\ueps{$Y$}{j}{k}
        }
        \mathrm{d}t_1 =
        %
        \frac{
            e^{
                i \cdot \ueps{$X$}{l}{m} \cdot t
            }
            -1
        }{
            \ueps{$X$}{l}{m} \cdot\ueps{$Y$}{j}{k}
        } + 
        \int\limits_0^t 
        \frac{
            e^{
                i \cdot \left(\ueps{$X$}{l}{m} + \ueps{$Y$}{j}{k} \right) \cdot t_1
            }
        }{
            i \cdot\ueps{$Y$}{j}{k}
        }
        \mathrm{d}t_1\\=
        %
        &\begin{cases}
            \ueps{$X$}{l}{m} + \ueps{$Y$}{j}{k} = 0: \quad
            \frac{
                e^{
                    i \cdot \ueps{$X$}{l}{m} \cdot t
                }
                -1
            }{
                \ueps{$X$}{l}{m} \cdot\ueps{$Y$}{j}{k}
            } - \frac{i \cdot t}{\ueps{$Y$}{j}{k}}            \\
            %%%%%
            \ueps{$X$}{l}{m} + \ueps{$Y$}{j}{k} \neq 0: \quad
            \frac{
                e^{
                    i \cdot \ueps{$X$}{l}{m} \cdot t
                }
                -1
            }{
                \ueps{$X$}{l}{m} \cdot\ueps{$Y$}{j}{k}
            } - \frac{
                e^{
                    i \cdot \left(\ueps{$X$}{l}{m} + \ueps{$Y$}{j}{k} \right) \cdot t
                }
                -1
            }{\ueps{$Y$}{j}{k} \cdot \left(\ueps{$X$}{l}{m} + \ueps{$Y$}{j}{k} \right)}        
        \end{cases}
    \end{split}
\end{equation}


\begin{equation}
    \label{eq:time-ordered-integral-part2}
    \begin{split}
        \Gamma_2&\left(X,\,l,\,m,\,Y,\,j,\,k\right) =\\ =&\int\limits_0^t  \int\limits_{t_1}^t
        e^{
            i \cdot \ueps{$X$}{l}{m} \cdot t_1
        }
        e^{
            i \cdot \ueps{$Y$}{j}{k} \cdot t_2
        }
        \mathrm{d}t_2
        \mathrm{d}t_1 = 
        %
        \int\limits_0^t 
        e^{
            i \cdot \ueps{$X$}{l}{m} \cdot t_1
        }
        \int\limits_{t_1}^t
        e^{
            i \cdot \ueps{$Y$}{j}{k} \cdot t_2
        }
        \mathrm{d}t_2
        \mathrm{d}t_1 \\=
        %
        &\int\limits_0^t 
        e^{
            i \cdot \ueps{$X$}{l}{m} \cdot t_1
        }
        \frac{
            e^{
                i \cdot \ueps{$Y$}{j}{k} \cdot t
            }
            -
            e^{
                i \cdot \ueps{$Y$}{j}{k} \cdot t_1
            }
        }{
            i \cdot\ueps{$Y$}{j}{k}
        }
        \mathrm{d}t_1 \\=
        %
        &\frac{
            e^{
                i \cdot \ueps{$Y$}{j}{k} \cdot t
            }
            -
            e^{
                i \cdot \left(\ueps{$X$}{l}{m} + \ueps{$Y$}{j}{k} \right) \cdot t
            }
        }{
            \ueps{$X$}{l}{m} \cdot\ueps{$Y$}{j}{k}
        } -
        \int\limits_0^t 
        \frac{
            e^{
                i \cdot \left(\ueps{$X$}{l}{m} + \ueps{$Y$}{j}{k} \right) \cdot t_1
            }
        }{
            i \cdot\ueps{$Y$}{j}{k}
        }
        \mathrm{d}t_1\\=
        %
        &\begin{cases}
            \ueps{$X$}{l}{m} + \ueps{$Y$}{j}{k} = 0: \quad
            \frac{
                e^{
                    i \cdot \ueps{$Y$}{j}{k} \cdot t
                }
                -
                e^{
                    i \cdot \left(\ueps{$X$}{l}{m} + \ueps{$Y$}{j}{k} \right) \cdot t
                }
            }{
                \ueps{$X$}{l}{m} \cdot\ueps{$Y$}{j}{k}
            } + \frac{i \cdot t}{\ueps{$Y$}{j}{k}}            \\
            %%%%%
            \ueps{$X$}{l}{m} + \ueps{$Y$}{j}{k} \neq 0: \quad
            \frac{
                e^{
                    i \cdot \ueps{$Y$}{j}{k} \cdot t
                }
                -
                e^{
                    i \cdot \left(\ueps{$X$}{l}{m} + \ueps{$Y$}{j}{k} \right) \cdot t
                }
            }{
                \ueps{$X$}{l}{m} \cdot\ueps{$Y$}{j}{k}
            } + \frac{
                e^{
                    i \cdot \left(\ueps{$X$}{l}{m} + \ueps{$Y$}{j}{k} \right) \cdot t
                }
                -1
            }{\ueps{$Y$}{j}{k} \cdot \left(\ueps{$X$}{l}{m} + \ueps{$Y$}{j}{k} \right)}  
        \end{cases}
    \end{split}
\end{equation}

Adding \autoref{eq:time-ordered-integral-part1} and \autoref{eq:time-ordered-integral-part2}, \autoref{eq:time-ordered-integral-parts-sum} can be verified, which is necessary when comparing with the integration result of the second half of \autoref{eq:hn-integrated-second-order}.

\begin{equation}
    \label{eq:time-ordered-integral-parts-sum}
    \begin{split}
        &\Gamma_1 \left(X,\,l,\,m,\,Y,\,j,\,k\right) + \Gamma_2 \left(X,\,l,\,m,\,Y,\,j,\,k\right) = \\
        &=
        \frac{
            e^{
                i \cdot \ueps{$X$}{l}{m} \cdot t
            }
            -1
        }{
            \ueps{$X$}{l}{m} \cdot\ueps{$Y$}{j}{k}
        }
         + 
        \frac{
            e^{
                i \cdot \ueps{$Y$}{j}{k} \cdot t
            }
            -
            e^{
                i \cdot \left(\ueps{$X$}{l}{m} + \ueps{$Y$}{j}{k} \right) \cdot t
            }
        }{
            \ueps{$X$}{l}{m} \cdot\ueps{$Y$}{j}{k}
        } \\
        & = 
        -\frac{
            e^{
                i \cdot \ueps{$X$}{l}{m} \cdot t
            }
            -
            1
        }{
            \ueps{$X$}{l}{m}
        }
        \cdot 
        \frac{
            e^{
                i \cdot \ueps{$Y$}{j}{k} \cdot t
            }
            -
            1
        }{
            \ueps{$Y$}{j}{k}
        }
         = 
         - \VhamiltonianAnalyticalPartIntegrated{$X$}{l}{m}{t} 
         \cdot
         \VhamiltonianAnalyticalPartIntegrated{$Y$}{j}{k}{t}
    \end{split}
\end{equation}

With the integrals rewritten, \autoref{eq:hn-integrated-second-order} can be re-arranged to obtain the final result: 

\begin{equation}
    \label{eq:hn-integrated-second-order-final}
    \begin{split}
        &\HNOftOrder{2} \stackrel{\ref{eq:time-ordered-integral},\,\ref{eq:time-ordered-integral-part1},\,\ref{eq:time-ordered-integral-part2}}{=} 
        - \frac{J^2}{2} 
        \lsum[X,Y \in \left\{\text{A}, \text{B}, \text{C}\right\}] \fullneighborsum{l}{m}\fullneighborsum{j}{k}\\
        %
        \left[         \vphantom{\lsum[K]}        \right.
        &\phantom{+}
        \lsum[K]
        \frac{\psiN[K]}{\psiN} 
        \Gamma_1 \left(X,\,l,\,m,\,Y,\,j,\,k\right)
        \cdot 
        \bracketHelper{N}{
            \VhamiltonianOperatorPart[]{$X$}{l}{m}
            \cdot 
            \VhamiltonianOperatorPart[]{$Y$}{j}{k}
            }{K}
        \\
        &+ 
        \lsum[K]
        \frac{\psiN[K]}{\psiN} 
        \Gamma_2 \left(X,\,l,\,m,\,Y,\,j,\,k\right)
        \cdot 
        \bracketHelper{N}{
            \VhamiltonianOperatorPart[]{$Y$}{j}{k}
        \cdot 
            \VhamiltonianOperatorPart[]{$X$}{l}{m}
        }{K} \\
        &-
        \left.
        \lsum[M]\lsum[L]
        \frac{\psiN[M]\psiN[L]}{\psiN{}^2}
            \left[\Gamma_1 \left(X,\,l,\,m,\,Y,\,j,\,k\right) + \Gamma_2 \left(X,\,l,\,m,\,Y,\,j,\,k\right)\right] \cdot
            \bracketHelper{N}{\VhamiltonianOperatorPart{$X$}{l}{m}}{M}
            \cdot
            \bracketHelper{N}{\VhamiltonianOperatorPart{$Y$}{j}{k}}{L}
        \right]
        \hspace{-1cm} % put the equation number where there is space
    \end{split}
\end{equation}

As for the first order perturbation, this seems to have made this computationally more expensive than the full diagonalization.
But again, because of the limited interaction range (for second order the range is a sphere of at most two nearest-neighbor-hops), this can be efficiently evaluated in 
\bigo{\text{\#}(\text{sites}) \cdot \text{\#}(\text{nearest neighbors})^2} by using the sparse nature of the operators.

The Operator $\VhamiltonianOperatorPart[]{$X$}{l}{m}\VhamiltonianOperatorPart[]{$Y$}{j}{k}$ is actually already much too complicated, to write down any of its mixed terms in a ordered, closed form.
To get around this, in the implementation, the first half is applied to the left, the state modified accordingly and then the second half is applied - as this is much more efficient than calculating the closed form.

See \filepath{\cite{selfCode}}{/calculation-helpers/generatesecondorder.py}.

Finally - upon close inspection - one can spot, that \autoref{eq:hn-integrated-second-order-final} is $0$ for all cases, where $l \neq j \land l \neq k \land m \neq j \land n \neq k$, if specific properties hold for \psiN[\ast].
In a subsequent section this optimization will be used.
        \FloatBarrier

    \section{Observable Measurements}
    \label{sec:theory-observables}
    When performing time-evolution with a perturbative approach, one ought to watch out for resonances.
In the case of degenerate energy solutions, so-called \emph{secular terms} \cite{secularTermsPerturbation} might be introduced, if non-degenerate perturbation theory is used to find the solution to a degenerate problem. 
Often inaccuracies manifest as linear terms in a function where only oscillating terms would be expected.
The linear term constantly grows over time, guaranteeing the approximation is only valid for short times and then diverges further and further from the exact solution.

For the examined Hamiltonian the first occurrence of such terms can be spotted in the second order cumulant expansion of the Hamiltonian.
\autoref{eq:time-ordered-integral-part1} and \autoref{eq:time-ordered-integral-part2} specifically list extra solutions for the cases where the two energies are the same.
This resonance causes a linear term to appear in the expansion, while all other pre-factors so far were complex exponentials -- meaning they are oscillating.

The idea to combat this is to replace the perturbative parameters that were calculated analytically with variational parameters instead.
This has the motive to generate a physically motivated structure, but the fine details of the parameter weights are adjusted by optimization.
E.g. for the \emph{neural-network quantum state} \cite{neuralNetworkQuantumStates} approach, a neural network structure is used to represent a wave-function.
Reason for this being the hope that the exponentially complicated wave-function can be expressed as a \glqq higher truth\grqq{} only requiring few parameters to describe a seemingly complicated wave-function.
While it would be basically impossible to always find the few-parameter-solution, \emph{approximating} it close enough with optimization akin to machine-learning should be a suitable replacement.

The following section aims to achieve better approximations than the ones found by cumulant expansion by letting the parameters be optimized during the step to reflect the situation more closely.
    \FloatBarrier

        \subsection{Observable Sampling}
        \label{sec:theory-observables-sampling}
        The evaluation of observables is a primary requirement for this application. 
First it is necessary to acquire a more general formula, that one later can plug specific observables into.

From \autoref{eq:time-evolution-target} and applying \emph{dagger} $\dagger$ to said equation, \autoref{eq:time-evolved-state-and-dagger} follows.


\begin{equation}
    \label{eq:time-evolved-state-and-dagger}
    \begin{split}
        \braketHelper{N}{\psiOfT[\schroedingerPicture]} &\stackrel{\ref{eq:time-evolution-target}}{=} \braN{} \lsum[K] e^{\HeffOft[K]} \psiN[K] \ketN[K]{}\\
        & \stackrel{\phantom{\ref{eq:time-evolution-target}}}{=} 
        \lsum[K]   e^{\HeffOft[K]} \psiN[K] \underbrace{\braketHelper{N}{K}}_\text{$\delta_{N,K}$} = e^{\HeffOft} \psiN{} \\
        %
        \stackrel{\dagger}{\Rightarrow} \braketHelper{\psiOfT[\schroedingerPicture]}{N} &\stackrel{\phantom{\ref{eq:time-evolution-target}}}{=} e^{\HeffOftStar} \psiNStar
    \end{split}
\end{equation}

In \autoref{eq:expectation-value} by starting with the default way of calculating a normalized \emph{expectation-value} for an observable \ObservableOp \cite{monteCarloObservableSampling} and applying multiple modifications, one gets a representation that can be effectively sampled later. 

\begin{equation}
    \label{eq:expectation-value}
    \begin{split}
        &\frac{\bracketHelper{\psiOfT[\schroedingerPicture]}{\ObservableOp}{\psiOfT[\schroedingerPicture]}}{\braketHelper{\psiOfT[\schroedingerPicture]}{\psiOfT[\schroedingerPicture]}}
        %
        \stackrel{\phantom{\ref{eq:time-evolution-target}}}{=} 
        \frac{
            \lsum[N] \braketHelper{\psiOfT[\schroedingerPicture]}{N} \bracketHelper{N}{\ObservableOp}{\psiOfT[\schroedingerPicture]}
        }{
            \lsum[K] \braketHelper{\psiOfT[\schroedingerPicture]}{K} \braketHelper{K}{\psiOfT[\schroedingerPicture]}
        }\\
        %
        &\stackrel{\ref{eq:time-evolved-state-and-dagger}}{=} 
        \frac{
            \lsum[N] e^{\HeffOftStar} \psiNStar \bracketHelper{N}{\ObservableOp}{\psiOfT[\schroedingerPicture]}
        }{
            \lsum[K] e^{\HeffOftStar[K]} \psiNStar[K] e^{\HeffOft[K]} \psiN[K] 
        }\\
        & \stackrel{\phantom{\ref{eq:time-evolved-state-and-dagger}}}{=} 
        \frac{
            \lsum[N] e^{\HeffOftStar} \psiNStar \bracketHelper{N}{\ObservableOp}{\psiOfT[\schroedingerPicture]}
        }{
            \lsum[K] \absSquare{e^{\HeffOft[K]}} \absSquare{\psiN[K]} 
        } \cdot 
        \frac{\braketHelper{N}{\psiOfT[\schroedingerPicture]}}{\braketHelper{N}{\psiOfT[\schroedingerPicture]}}\\
        & \stackrel{\phantom{\ref{eq:time-evolved-state-and-dagger}}}{=} 
        \lsum[N]
        \frac{
            e^{\HeffOftStar} \psiNStar e^{\HeffOft} \psiN{} \braketHelper{N}{\psiOfT[\schroedingerPicture]}
        }{
            \lsum[K] \absSquare{e^{\HeffOft[K]}} \absSquare{\psiN[K]} 
        }  \cdot 
        \frac{\bracketHelper{N}{\ObservableOp}{\psiOfT[\schroedingerPicture]}}{\braketHelper{N}{\psiOfT[\schroedingerPicture]}}\\
        & \stackrel{\ref{eq:time-evolved-state-and-dagger}}{=} 
        \lsum[N]
        \underbrace{\frac{
            \absSquare{e^{\HeffOft}} \absSquare{\psiN} 
        }{
            \lsum[K] \absSquare{e^{\HeffOft[K]}} \absSquare{\psiN[K]} 
        }}_{\probabilityOf{N}{t}} \cdot 
        \frac{\bracketHelper{N}{\ObservableOp}{\psiOfT[\schroedingerPicture]}}{\braketHelper{N}{\psiOfT[\schroedingerPicture]}}\\
        & \stackrel{\phantom{\ref{eq:time-evolved-state-and-dagger}}}{=} \lsum[N]
        \probabilityOf{N}{t}
        \frac{\bracketHelper{N}{\ObservableOp}{\psiOfT[\schroedingerPicture]}}{\braketHelper{N}{\psiOfT[\schroedingerPicture]}}
        %
        \stackrel{\phantom{\ref{eq:time-evolved-state-and-dagger}}}{=} \lsum[N]
        \probabilityOf{N}{t}
        \frac{
            \lsum[K] \bracketHelper{N}{\ObservableOp}{K} \braketHelper{K}{\psiOfT[\schroedingerPicture]} 
        }{
            \braketHelper{N}{\psiOfT[\schroedingerPicture]}
        }\\
        & \stackrel{\ref{eq:time-evolved-state-and-dagger}}{=} \lsum[N]
        \probabilityOf{N}{t}
        \frac{
            \lsum[K] \bracketHelper{N}{\ObservableOp}{K} e^{\HeffOft[K]} \psiN[K]
        }{
            e^{\HeffOft[N]} \psiN[N]
        }\\
        %
        &\stackrel{\phantom{\ref{eq:time-evolved-state-and-dagger}}}{=} \lsum[N]
        \probabilityOf{N}{t}
        \underbrace{\lsum[K] \bracketHelper{N}{\ObservableOp}{K} e^{\HeffOft[K]-\HeffOft[N]}
        \frac{
            \psiN[K]
        }{
            \psiN[N]
        }}_{\localObservable{N}{t}}
        %
        \stackrel{\phantom{\ref{eq:time-evolved-state-and-dagger}}}{=} \lsum[N]
        \probabilityOf{N}{t}
        \localObservable{N}{t}
    \end{split}
\end{equation}

By inserting \autoref{eq:time-evolved-state-and-dagger} into \autoref{eq:expectation-value}, a probability-dependent formula is obtained, that can be used to plug in specific observables.  
The probability of a base-state at a specified time \probabilityOf{N}{t} is important for this calculation. 
The \fullref{sec:theory-monte-carlo} will discuss the viability of obtaining these probabilities and strategies of doing so.
The matrix-element $\bracketHelper{N}{\ObservableOp}{K}$ has the important property of being $0$ for almost all combinations of $\braN$ and $\ketN[K]$, depending on the chosen observable \ObservableOp, which is what makes the calculation of the \emph{local-observable} \localObservable{N}{t} possible in the first place.

        \FloatBarrier

        \subsection{Occupation Types}
        \label{sec:theory-observables-occupations}
        % \subsubsection*{Double-Occupation}

One possibly interesting observable would be the measurement of \emph{double-occupation} (one particle of spin-up and one of spin-down) per site.

The operator $\doubleOccupationOperator{l}$ that measure the double-occupation on site $l$ can be easily written as $\doubleOccupationOperator{l} = \nop{l}{\up}\nop{l}{\down}$.

As the used basis-states are eigenstates of the number-operators, the evaluation of this operator is quite straight forward. 
$\bracketHelper{N}{\doubleOccupationOperator{l}}{K}$ becomes $\delta_{N,\,K}\cdot n_{l,\,\up}\cdot n_{l,\,\down}$ with the occupation-number (not operator) of the specific site and spin in \braN, $n_{l,\,\sigma}$. 
With that \localObservable{N}{t} can be simply evaluated to the result in \autoref{eq:double-occupation-loc}.

\begin{equation}
    \label{eq:double-occupation-loc}
    \localObservable{N}{t} = \lsum[K] \bracketHelper{N}{\doubleOccupationOperator{l}}{K} e^{\HeffOft[K]-\HeffOft[N]}
    \frac{
        \psiN[K]
    }{
        \psiN[N]
    } = n_{l,\,\up}\cdot n_{l,\,\down}
\end{equation}

The operator, that measures \emph{single-occupation} $\singleOccupationOperator{l}{\sigma}$ is derived exactly in the same way.

\begin{equation}
    \label{eq:single-occupation-loc}
    \localObservable{N}{t} = \lsum[K] \bracketHelper{N}{\singleOccupationOperator{l}{\sigma}}{K} e^{\HeffOft[K]-\HeffOft[N]}
    \frac{
        \psiN[K]
    }{
        \psiN[N]
    } = n_{l,\,\sigma}
\end{equation}
        \FloatBarrier

        \subsection{Spin-Current}
        \label{sec:theory-observables-current}
        % \subsubsection*{Spin-Polarized Kinetics}

For measuring the spin-polarized flow of particles, a slightly more complex observable must be employed. 
The current operator is dependent on the hamiltonian and must be derived from the \emph{continuity equation} (\autoref{eq:continuity-equation}) in the \emph{Heisenberg Picture} (\autoref{eq:heisenberg-picture}) \cite{discreteDivergenceAndOtherCurrentDerivation}.

\begin{equation}
    \label{eq:heisenberg-picture}
    \begin{split}
        \AopOfT[\heisenbergPicture] &= e^{i \pictureHamiltonian[\schroedingerPicture]{} t} \Aop[\schroedingerPicture] e^{-i \pictureHamiltonian[\schroedingerPicture]{} t}
        =  \heisenbergTimeEvolution{\Aop[\schroedingerPicture]}
        \quad \Rightarrow \quad 
        \pictureHamiltonian[\schroedingerPicture]{} = \pictureHamiltonian[\heisenbergPicture]{}\\
    \end{split}
\end{equation}

It is important to notice, that the time-evolution in the Heisenberg Picture \heisenbergTimeEvolution{\ast} is a different one, than the previous \interactionTimeEvolution{\ast}.
For the discrete gradient \discreteGradient it is important to remember, that the system lives on a discrete, 2-dimensional lattice.
There, $\partial_\text{x} f(k) = f(k+1) - f(k)$, with any site $k$ and $k+1$ being the \glqq{}next\grqq{} site in positive $\text{x}$-direction.

\begin{equation}
    \label{eq:continuity-equation}
    \begin{split}
        \difft{\nopOfT[\heisenbergPicture]{k}{\sigma}} &= 
        i \left[\pictureHamiltonian[\heisenbergPicture],\, \nopOfT[\heisenbergPicture]{k}{\sigma}\right] = 
        i \heisenbergTimeEvolution{
            \left[\pictureHamiltonian[\schroedingerPicture],\, \nop[\schroedingerPicture]{k}{\sigma}\right]
        }\\ 
        &= - \discreteGradient \cdot \currentvecOft[\heisenbergPicture]{k}{\sigma}
         = - \heisenbergTimeEvolution{\discreteGradient \cdot \currentvec[\schroedingerPicture]{k}{\sigma} } \\
        \stackrel{\text{2-dim.}}{\Rightarrow}
        \left(\begin{matrix}
            \partial_\text{x}\\
            \partial_\text{y}\\
        \end{matrix}\right) \cdot 
        \left(\begin{matrix}
            j_{\text{x}\,\,k,\,\sigma}^{\,\,\schroedingerPicture}\\
            j_{\text{y}\,\,k,\,\sigma}^{\,\,\schroedingerPicture}
        \end{matrix}\right) &=
        \partial_\text{x}
        j_{\text{x}\,\,k,\,\sigma}^{\,\,\schroedingerPicture}+
        \partial_\text{y} 
        j_{\text{y}\,\,k,\,\sigma}^{\,\,\schroedingerPicture}
         = - i \cdot \left[\pictureHamiltonian[\schroedingerPicture],\, \nop[\schroedingerPicture]{k}{\sigma}\right]
    \end{split}
\end{equation}

The commutator can swiftly be evaluated, a reference calculation can also be found in 

\filepath{\cite{selfMathManipulatorCalculations}}{/current}. The calculation is somewhat specific to the dimension and geometry, in terms of what neighbor indices are summed over.
For arguing about indices and directions, without worrying about geometric details, \biggerNeighbor{l}{m} will describe all possible values of $m$ that fulfill both the properties that the site $m$ is a nearest neighbor of the site $l$ and the site of $m$ lies either in positive $\text{x}$-direction, or $\text{y}$-direction of the other site, but not in negative direction.
This bisects the sites along the primary diagonal of the coordinate system and makes the calculation work for all lattices with symmetry around the $\text{x}$- and $\text{y}$-axis.

\begin{equation}
    \label{eq:calculation-current-commutator}
    \begin{split}
         \left[\pictureHamiltonian[\schroedingerPicture],\, \nop[\schroedingerPicture]{k}{\sigma}\right] &\stackrel{\ref{eq:main-hamiltonian}}{=}
         \left[\HzeroHamiltonian[\schroedingerPicture],\, \nop[\schroedingerPicture]{k}{\sigma}\right] + \left[\Vhamiltonian[\schroedingerPicture],\, \nop[\schroedingerPicture]{k}{\sigma}\right] \stackrel{\ref{eq:stuff-from-math-manipulator-dd}}{=}
         \left[\Vhamiltonian[\schroedingerPicture],\, \nop[\schroedingerPicture]{k}{\sigma}\right]\\
         &\stackrel{\ref{eq:main-hamiltonian-perturbation-full-sum}}{=}
         - J \cdot \fullneighborsum{l}{m} 
        \left(
            \left[\hop[\schroedingerPicture]{l}{\dagger}\hop[\schroedingerPicture]{m}{},\, \nop[\schroedingerPicture]{k}{\sigma}\right]
            + 
            \left[\dop[\schroedingerPicture]{l}{\dagger}\dop[\schroedingerPicture]{m}{},\, \nop[\schroedingerPicture]{k}{\sigma}\right]
        \right)\\
        &\stackrel{\text{MM}}{=}
        J \cdot \fullneighborsum{l}{m} 
        \left(
            \withspinhcop[\schroedingerPicture]{l}{\sigma}{\dagger}
            \withspinhcop[\schroedingerPicture]{m}{\sigma}{}
            \cdot \delta_{k,\,l}
            -
            \withspinhcop[\schroedingerPicture]{l}{\sigma}{\dagger}
            \withspinhcop[\schroedingerPicture]{m}{\sigma}{}
            \cdot \delta_{k,\,m}
        \right)\\
        &\stackrel{\phantom{\text{MM}}}{=}
        J \cdot \neighborsum{l}{m} 
        \left(
            \withspinhcop[\schroedingerPicture]{l}{\sigma}{\dagger}
            \withspinhcop[\schroedingerPicture]{m}{\sigma}{}
            \cdot \delta_{k,\,l}
            -
            \withspinhcop[\schroedingerPicture]{l}{\sigma}{\dagger}
            \withspinhcop[\schroedingerPicture]{m}{\sigma}{}
            \cdot \delta_{k,\,m}
            +
            \withspinhcop[\schroedingerPicture]{m}{\sigma}{\dagger}
            \withspinhcop[\schroedingerPicture]{l}{\sigma}{}
            \cdot \delta_{k,\,m}
            -
            \withspinhcop[\schroedingerPicture]{m}{\sigma}{\dagger}
            \withspinhcop[\schroedingerPicture]{l}{\sigma}{}
            \cdot \delta_{k,\,l}
        \right)\\
        &\stackrel{\phantom{\text{MM}}}{=} J \cdot \lsum \sum\limits_{\biggerNeighbor{l}{m}}
        \left(
            \withspinhcop[\schroedingerPicture]{l}{\sigma}{\dagger}
            \withspinhcop[\schroedingerPicture]{m}{\sigma}{}
            \cdot \delta_{k,\,l}
            -
            \withspinhcop[\schroedingerPicture]{m}{\sigma}{\dagger}
            \withspinhcop[\schroedingerPicture]{l}{\sigma}{}
            \cdot \delta_{k,\,l}
        \right)\\
        &\stackrel{\phantom{\text{MM}}}{+} J \cdot \lsum[m] \sum\limits_{\smallerNeighbor{m}{l}}
        \left(
            \withspinhcop[\schroedingerPicture]{m}{\sigma}{\dagger}
            \withspinhcop[\schroedingerPicture]{l}{\sigma}{}
            \cdot \delta_{k,\,m}
            -
            \withspinhcop[\schroedingerPicture]{l}{\sigma}{\dagger}
            \withspinhcop[\schroedingerPicture]{m}{\sigma}{}
            \cdot \delta_{k,\,m}
        \right)\\
        &\stackrel{\phantom{\text{MM}}}{=} J \cdot \sum\limits_{\biggerNeighbor{k}{m}}
        \left(
            \withspinhcop[\schroedingerPicture]{k}{\sigma}{\dagger}
            \withspinhcop[\schroedingerPicture]{m}{\sigma}{}
            -
            \withspinhcop[\schroedingerPicture]{m}{\sigma}{\dagger}
            \withspinhcop[\schroedingerPicture]{k}{\sigma}{}
        \right)
        + J \cdot \sum\limits_{\smallerNeighbor{k}{l}}
        \left(
            \withspinhcop[\schroedingerPicture]{k}{\sigma}{\dagger}
            \withspinhcop[\schroedingerPicture]{l}{\sigma}{}
            -
            \withspinhcop[\schroedingerPicture]{l}{\sigma}{\dagger}
            \withspinhcop[\schroedingerPicture]{k}{\sigma}{}
        \right)\\
        &\stackrel{\phantom{\text{MM}}}{=} J \cdot \sum\limits_{l: \,\left\langle k,\,l\right\rangle}
        \left(
            \withspinhcop[\schroedingerPicture]{k}{\sigma}{\dagger}
            \withspinhcop[\schroedingerPicture]{l}{\sigma}{}
            -
            \withspinhcop[\schroedingerPicture]{l}{\sigma}{\dagger}
            \withspinhcop[\schroedingerPicture]{k}{\sigma}{}
        \right)
    \end{split}
\end{equation}

% Define colors
\definecolor{arrow1}{HTML}{DB10E9} % pink
\definecolor{arrow1back}{HTML}{AA0000} % red
\definecolor{arrow2}{HTML}{000080} % blue
\definecolor{arrow2back}{HTML}{007700} % cyan
\definecolor{arrow3}{HTML}{FF8E00} % orange
\definecolor{arrow3back}{HTML}{EEBE00} % yellow
\definecolor{arrow4}{HTML}{2DCF2D} % mint / light green
\definecolor{arrow4back}{HTML}{007700} % green

\begin{SCfigure}[1.5][htbp]
    \centering
    \makebox[0.33\textwidth][c]{
        \begin{tikzpicture}
            % Central node
            \node[circle, draw, fill=white, minimum size=0.6cm] (center) at (0,0) {$k$};
        
            % Surrounding nodes
            \node[circle, draw, fill=white, minimum size=0.6cm] (n1) at (0,2) {1};
            \node[circle, draw, fill=white, minimum size=0.6cm] (n2) at (2,0) {2};
            \node[circle, draw, fill=white, minimum size=0.6cm] (n3) at (0,-2) {3};
            \node[circle, draw, fill=white, minimum size=0.6cm] (n4) at (-2,0) {4};
            
            % Arrows
            \draw[->, color=arrow1, bend left=20] (center) to (n1); % Center to 1
            \draw[<-, color=arrow1back, bend right=20] (center) to (n1); % 1 to Center

            \draw[->, color=arrow2, bend left=20] (center) to (n2); % Center to 2
            \draw[<-, color=arrow2back, bend right=20] (center) to (n2); % 2 to Center

            \draw[->, color=arrow3, bend left=20] (center) to (n3); % Center to 3
            \draw[<-, color=arrow3back, bend right=20] (center) to (n3); % 3 to Center

            \draw[->, color=arrow4, bend left=20] (center) to (n4); % Center to 4
            \draw[<-, color=arrow4back, bend right=20] (center) to (n4); % 4 to Center
            
        \end{tikzpicture}
    }
    \caption{Exemplary cutout of a square lattice. The nearest neighbors to a central node with index $k$ are depicted. The numeration is arbitrary, but matches the indices that are used in \autoref{eq:current-geometry-example}.
    The colors match the terms that correspond to their transition.
    }
    \label{fig:current-geometry-example}
\end{SCfigure}

\begin{equation}
    \label{eq:current-geometry-example}
    \begin{split}
        \partial_\text{x}
        j_{\text{x}\,\,k,\,\sigma}^{\,\,\schroedingerPicture}+
        \partial_\text{y} 
        j_{\text{y}\,\,k,\,\sigma}^{\,\,\schroedingerPicture} 
        &\stackrel{\phantom{\ref{eq:continuity-equation}}}{=}
        \textcolor{arrow2}{
            \spinPolarizedKineticsOperatorDir{k}{2}{\sigma}
        }
        -
        \textcolor{arrow4back}{
            \spinPolarizedKineticsOperatorDir{4}{k}{\sigma}
        }\\
        &\stackrel{\phantom{\ref{eq:continuity-equation}}}{+}
        \textcolor{arrow3}{
            \spinPolarizedKineticsOperatorDir{k}{3}{\sigma}
        }
        -
        \textcolor{arrow1back}{
            \spinPolarizedKineticsOperatorDir{1}{k}{\sigma}
        }\\
        & \stackrel{\ref{eq:continuity-equation}}{=}
         - i \cdot \left[\pictureHamiltonian[\schroedingerPicture],\, \nop[\schroedingerPicture]{k}{\sigma}\right]
         \stackrel{\ref{eq:calculation-current-commutator}}{=}\\
         & \hspace{-3cm} i\cdot J \left(
            \textcolor{arrow2}{
                \withspinhcop[\schroedingerPicture]{2}{\sigma}{\dagger}
                \withspinhcop[\schroedingerPicture]{k}{\sigma}{}
            }
            -
            \textcolor{arrow2back}{
                \withspinhcop[\schroedingerPicture]{k}{\sigma}{\dagger}
                \withspinhcop[\schroedingerPicture]{2}{\sigma}{}
            }
            -
            \textcolor{arrow4back}{
                \withspinhcop[\schroedingerPicture]{k}{\sigma}{\dagger}
                \withspinhcop[\schroedingerPicture]{4}{\sigma}{}
            }
            +
            \textcolor{arrow4}{
                \withspinhcop[\schroedingerPicture]{4}{\sigma}{\dagger}
                \withspinhcop[\schroedingerPicture]{k}{\sigma}{}
            }
            +
            \textcolor{arrow3}{
                \withspinhcop[\schroedingerPicture]{3}{\sigma}{\dagger}
                \withspinhcop[\schroedingerPicture]{k}{\sigma}{}
            }
            -
            \textcolor{arrow3back}{
                \withspinhcop[\schroedingerPicture]{k}{\sigma}{\dagger}
                \withspinhcop[\schroedingerPicture]{3}{\sigma}{}
            }
            -
            \textcolor{arrow1back}{
                \withspinhcop[\schroedingerPicture]{k}{\sigma}{\dagger}
                \withspinhcop[\schroedingerPicture]{1}{\sigma}{}
            }
            +
            \textcolor{arrow1}{
                \withspinhcop[\schroedingerPicture]{1}{\sigma}{\dagger}
                \withspinhcop[\schroedingerPicture]{k}{\sigma}{}
            }
         \right)
    \end{split}
\end{equation}


\autoref{eq:spin-polarized-kinetics-operator-definition} describes the resulting observable, that measures such kinetics, depending on the direction \spinPolarizedKineticsOperatorDir{l}{m}{\sigma}.

\begin{equation}
    \label{eq:spin-polarized-kinetics-operator-definition}
    \begin{split}
        \spinPolarizedKineticsOperatorDir{l}{m}{\sigma} &= i J \left(\withspinhcop{m}{\sigma}{\dagger}\withspinhcop{l}{\sigma}{} - \withspinhcop{l}{\sigma}{\dagger}\withspinhcop{m}{\sigma}{}\right)\\
    \end{split}
\end{equation}

In this case, the used basis-states are not eigenstates of the operators. 
$\bracketHelper{N}{\withspinhcop{m}{\sigma}{\dagger}\withspinhcop{l}{\sigma}{}}{K}$ becomes $\delta_{N,\,\vphantom{N}\smash{\tilde{N}}}\cdot n_{m,\,\sigma} \cdot (1-n_{l,\,\sigma} )$, where \ketN[{\vphantom{N}\smash{\tilde{N}}}] is the state obtained when the particle number on site $m,\, \sigma$ and the one of site $l,\, \sigma$ are swapped (this \emph{hopping} only has an effect, when there is a particle on the original site and no particle yet on the target site, or vice versa).
Evaluating the whole operator with the signs correctly and in a efficient manner can best be looked up in the implementation 

\filepath{\cite{selfCode}}{/computation-scripts/observables.py}.

Overall, the evaluation of this observable requires knowing the value of an object of the form presented in \autoref{eq:form-heff-difference}.

\begin{equation}
    \label{eq:form-heff-difference}
    i \cdot 
    \left[
        n_{m,\,\sigma} \cdot (1-n_{l,\,\sigma})
        -
        n_{l,\,\sigma} \cdot (1-n_{m,\,\sigma})
    \right]
    \cdot 
    \frac{\psiN[\vphantom{N}\smash{\tilde{N}}]}{\psiN}
    \cdot 
     e^{\HeffOft[\vphantom{N}\smash{\tilde{N}}]-\HeffOft[N]}
\end{equation}

The \fullref{sec:theory-optimizations-analytical} will go into how to compute this efficiently for states \ketN[N] and \ketN[{\vphantom{N}\smash{\tilde{N}}}] that are connected via the hopping between nearest-neighbor lattice sites.

It is finally important to notice, that the objects from \autoref{eq:form-heff-difference} are not real-valued, but complex.
Only the complete observable from the full sum over all basis-states \ketN over \localObservable{N}{t} has a fully vanishing imaginary part. 
Especially when the observable is approximated with an incomplete set of basis-states, computationally an imaginary component remains. 
It is important to monitor the magnitude of this value and monitor it going to $0$.
If it doesn't fully vanish, this indicates an error in the sampling strategy or the implementation or the number of sampled states is not great enough.

One could entertain the question, whether is is guaranteed for the imaginary part to vanish with a large enough sample size.
On first thought, if the hamiltonian is only approximated, might there be a possibility for asymmetrical terms to introduce a non-vanishing imaginary part?

Conveniently, rewriting \autoref{eq:expectation-value} in the manner of \autoref{eq:imaginary-part-vanishes}, yields two expressions for the value of the observation that are complex conjugations of each other.
This again means, that no matter how \glqq bad\grqq{} the \psiOfT[\schroedingerPicture] approximates the real wave-function, for full sampling the result must have an imaginary part of 0.

\begin{equation}
    \label{eq:imaginary-part-vanishes}
    \begin{split}
        \frac{\bracketHelper{\psiOfT[\schroedingerPicture]}{\ObservableOp}{\psiOfT[\schroedingerPicture]}}{\braketHelper{\psiOfT[\schroedingerPicture]}{\psiOfT[\schroedingerPicture]}}
        &\cdot
        \left(
            \lsum[K] \absSquare{e^{\HeffOft[K]}} \absSquare{\psiN[K]} 
        \right) =\\
        %
        \lsum[N]
        \absSquare{e^{\HeffOft}} \absSquare{\psiN} 
        \frac{\bracketHelper{N}{\ObservableOp}{\psiOfT[\schroedingerPicture]}}{\braketHelper{N}{\psiOfT[\schroedingerPicture]}}
        &\stackrel{\ast}{=}
        \lsum[N]
        \absSquare{e^{\HeffOft}} \absSquare{\psiN} 
        \frac{\bracketHelper{\psiOfT[\schroedingerPicture]}{\ObservableOp}{N}}{\braketHelper{\psiOfT[\schroedingerPicture]}{N}}
    \end{split}
\end{equation}
        \FloatBarrier

        \subsection{Reduced Density Matrix}
        \label{sec:theory-observables-density-matrix}
        Typically, the \emph{density matrix} $\rho = \ketpsiof[]{t}\brapsiof[]{t}$ encodes the full quantum information of the state. 
Because the number of entries in the matrix is the number of base-states squared - which is exponential in the number of sites, exactly what the approximation is supposed to avoid - it is not feasible to access this quantity.

While it is not possible to access the complete matrix, there exists one closely related entity, that has multiple uses that are going to be utilized in the subsequent section.
The \emph{reduced density matrix} $\rho_A$ is obtained upon \emph{tracing out} (\autoref{eq:tracing-out}, \cite{partialTraceEntanglementOfSubsystemsBlochVector}) the part $B$ from a system that is completely partitioned into two parts $A$ and $B$.
The basis of $A$ is $\left\{\ketN[\Phi_k]\right\}_{k=0}^{2^{2}-1 = 3 = N} = \left\{\ketN[11],\,\ketN[10],\,\ketN[01],\,\ketN[00]\right\}$ and the basis of $B$ $\left\{\ketN[\Chi_k]\right\}_{k=0}^{2^{2 \cdot \#\left(\text{sites}\right) - 2}-1 = M}$.
Notice, the order of basis states starting with $\ketN[1\dots1]$ then $\ketN[1\dots10]$ and ending with $\ketN[0\dots0]$. As stated in \autoref{eq:spin-hcb-state-mapping}, this is the convention that is necessary to keep consistent sings, as $\ketN[0\dots0] \rightarrow \ketN[1\dots1]$ would introduce exactly one $-1$ into all measurements with \pauli{y}.

\begin{equation}
    \label{eq:tracing-out}
    \rho_A = \partialTrace{B}{\rho_{AB}} = \sum\limits_{k=0}^{M} \left(\one[A]\otimes\braN[\Chi_k] \right)\rho_{AB}\left(\one[A]\otimes\ketN[\Chi_k]\right)
\end{equation}

If - for this application - one reduces the part $A$ to contain only two particles $l, \sigma$ and $m, \mu$, then $\rho_A$ will be a complex $4\times 4$ matrix.
As it is again not feasible to obtain via tracing out, the matrix can be calculated by rewriting it in the basis of Pauli matrices (\ref{sec:particles}):

\begin{equation}
    \label{eq:reduced-density-matrix-via-paulis}
    \rho_A (t) = \rho_{l,\,\sigma,\,m,\,\mu} (t) = \frac{1}{4} \sum\limits_{\alpha,\beta\in\left\{0,\text{x},\text{y},\text{z}\right\}}
    {\left\langle \pauli[l,\,\sigma]{$\alpha$}\pauli[m,\,\mu]{$\beta$} \right\rangle}_{\psiOfT} \left(\pauli[]{$\alpha$}\otimes\pauli[]{$\beta$}\right)
\end{equation}

\autoref{eq:reduced-density-matrix-via-paulis} uses the property, that the 16 matrices $\pauli{$\alpha$}\otimes\pauli{$\beta$}$ are a complete basis of the complex $4\times 4$ matrices.
The observable $\pauli[l,\,\sigma]{$\alpha$}\pauli[m,\,\mu]{$\beta$}$ can be translated into hard-core bosonic operators, thanks to the previously derived \autoref{eq:jordan-wigner-spin-hcb} and with that can easily be measured like all other observables.
The factor $\frac{1}{4}$ follows from a normalization argument in \autoref{eq:reduced-density-matrix-normalization}, that uses properties of density matrices \cite{partialTraceEntanglementOfSubsystemsBlochVector} and algebraic properties of the Pauli matrices \cite{pauliPropertiesAndBlochVector}.

\begin{equation}
    \label{eq:reduced-density-matrix-normalization}
    \begin{split}
        {\left\langle \pauli[l,\,\sigma]{$\alpha$}\pauli[m,\,\mu]{$\beta$} \right\rangle}_{\psiOfT} &= \partialTrace{}{\rho_{l,\,\sigma,\,m,\,\mu} (t) \left(\pauli[l,\,\sigma]{$\alpha$}\otimes\pauli[m,\,\mu]{$\beta$}\right)}  \\
        &\stackrel{\ref{eq:reduced-density-matrix-via-paulis}}{=}\sum\limits_{\alpha',\beta'\in\left\{0,\text{x},\text{y},\text{z}\right\}} \frac{{\left\langle \pauli[l,\,\sigma]{$\alpha'$}\pauli[m,\,\mu]{$\beta'$} \right\rangle}_{\psiOfT}}{4}
        \partialTrace{}{\left(\pauli{$\alpha'$}\otimes\pauli{$\beta'$}\right)\left(\pauli{$\alpha$}\otimes\pauli{$\beta$}\right)}\\
        &\stackrel{\phantom{\ref{eq:reduced-density-matrix-via-paulis}}}{=}\sum\limits_{\alpha',\beta'\in\left\{0,\text{x},\text{y},\text{z}\right\}} \frac{{\left\langle \pauli[l,\,\sigma]{$\alpha'$}\pauli[m,\,\mu]{$\beta'$} \right\rangle}_{\psiOfT}}{4}
        \delta_{\alpha',\,\alpha}\delta_{\beta',\,\beta}\partialTrace{}{\one[4\times 4]}\\
        &\stackrel{\phantom{\ref{eq:reduced-density-matrix-via-paulis}}}{=}\frac{4}{4} \cdot {\left\langle \pauli[l,\,\sigma]{$\alpha$}\pauli[m,\,\mu]{$\beta$} \right\rangle}_{\psiOfT}
    \end{split}
\end{equation}

\autoref{eq:reduced-density-matrix-via-paulis} now has the advantage, that the ${\left\langle \pauli[l,\,\sigma]{$\alpha$}\pauli[m,\,\mu]{$\beta$} \right\rangle}_{\psiOfT}$ can be evaluated by sampling from a state \psiOfT, that has been acquired by any desired means or sampling methods.

The operators $\left(\pauli[]{$\alpha$}\otimes\pauli[]{$\beta$}\right)$ were manually translated into the required complex base-matrices in 
\filepath{\cite{selfCode}}{/calculation-helpers/concurrence/spin-to-z-basis-transformation.py}, which is the same as the \emph{Kronecker product} \cite{kroneckerProduct} of the known matrix-representations (\autoref{eq:pauli-matrices}), precisely because of the consistent basis-convention.

As hinted, one can now employ \autoref{eq:jordan-wigner-spin-hcb} to translate the operators $\pauli[l,\,\sigma]{$\alpha$}\pauli[m,\,\mu]{$\beta$}$ into a form that may be inserted into \localObservable{N}{t} in \autoref{eq:expectation-value}. The non-trivial translations are listed in \autoref{eq:reduced-density-matrix-final-operators} (given $l\neq m$, the rest of the 16 factors can be trivially mirrored or multiplied from the give ones, because the $\pauli[]{0} = \one[]$ and \pauli[]{z} operators are diagonal in the used basis and the eigenvalue of all \one[] is 1).

\begin{equation}
    \label{eq:reduced-density-matrix-final-operators}
    \begin{split}
        \pauli[l,\,\sigma]{x} \pauli[m,\,\mu]{0}  &\hat{=}
        \frac{\psiN[\vphantom{N}\smash{\tilde{N}}]}{\psiN}        \cdot         e^{\HeffOft[\vphantom{N}\smash{\tilde{N}}]-\HeffOft[N]}
        \\
        \pauli[l,\,\sigma]{y} \pauli[m,\,\mu]{0}  &\hat{=}
        \frac{\psiN[\vphantom{N}\smash{\tilde{N}}]}{\psiN}        \cdot         e^{\HeffOft[\vphantom{N}\smash{\tilde{N}}]-\HeffOft[N]} \cdot i \cdot (1-2 \cdot n_{l,\,\sigma}) 
        \\
        \pauli[l,\,\sigma]{z} \pauli[m,\,\mu]{0}  &\hat{=}
        \frac{\psiN[\vphantom{N}\smash{\tilde{N}}]}{\psiN}        \cdot         e^{\HeffOft[\vphantom{N}\smash{\tilde{N}}]-\HeffOft[N]} \cdot (2 \cdot n_{l,\,\sigma} - 1) 
        \\
        \pauli[l,\,\sigma]{x} \pauli[m,\,\mu]{x}  &\hat{=}
        \frac{\psiN[\vphantom{N}\smash{\tildetilde{N}}]}{\psiN}        \cdot         e^{\HeffOft[\vphantom{N}\smash{\tildetilde{N}}]-\HeffOft[N]}
        \\
        \pauli[l,\,\sigma]{y} \pauli[m,\,\mu]{y}  &\hat{=}
        \frac{\psiN[\vphantom{N}\smash{\tildetilde{N}}]}{\psiN}        \cdot         e^{\HeffOft[\vphantom{N}\smash{\tildetilde{N}}]-\HeffOft[N]} \cdot (2 \cdot n_{l,\,\sigma} - 1) \cdot (1- 2 \cdot n_{m,\,\mu})
        \\
        \pauli[l,\,\sigma]{x} \pauli[m,\,\mu]{y}  &\hat{=}
        \frac{\psiN[\vphantom{N}\smash{\tildetilde{N}}]}{\psiN}        \cdot         e^{\HeffOft[\vphantom{N}\smash{\tildetilde{N}}]-\HeffOft[N]} \cdot i \cdot (1-2 \cdot n_{m,\,\mu}) 
        \\
        \pauli[l,\,\sigma]{y} \pauli[m,\,\mu]{x}  &\hat{=}
        \frac{\psiN[\vphantom{N}\smash{\tildetilde{N}}]}{\psiN}        \cdot         e^{\HeffOft[\vphantom{N}\smash{\tildetilde{N}}]-\HeffOft[N]} \cdot i \cdot (1-2 \cdot n_{l,\,\sigma}) 
    \end{split}
\end{equation}

Based on the sampled state \ketN[N], in \autoref{eq:reduced-density-matrix-final-operators} the state \ketN[\vphantom{N}\smash{\tilde{N}}] is the same state as \ketN[N], but with the occupation of site $l$ and spin direction $\sigma$ \emph{flipped}, meaning an occupation of 0 in \ketN[N] is 1 in \ketN[\vphantom{N}\smash{\tilde{N}}] and vice versa.
Furthermore, \ketN[\vphantom{N}\smash{\tildetilde{N}}] indicates a flipping of both the occupation of $l,\, \sigma$ and $m,\, \mu$.

The difference of effective hamiltonian with either one or two flipped occupations can be evaluated efficiently and is the last necessary optimization that will be presented in \fullref{sec:theory-optimizations}.
        \FloatBarrier
        
        \subsection{Entanglement Measurements}
        \label{sec:theory-observables-entanglement}
        The $\rho_{l,\,\sigma,\,m,\,\mu} (t)$ from \autoref{eq:reduced-density-matrix-via-paulis} has been specifically chosen to give the full quantum information of a subsystem of two arbitrary site-indices and spin-directions.
Most importantly, as the research is performed on quantum systems, it is a goal to extract quantum properties and the prime candidate for this is extracting the \emph{entanglement}.

Already research in \cite{isingDynamicsWithClassicalNetworks} used the same way of obtaining the reduced density matrix as presented in \autoref{eq:reduced-density-matrix-via-paulis} with the goal of measuring the entanglement.
There the number of sampled states, that is required to obtain the measurements of the observables, was called out as a limiting factor.
Monte-Carlo sampling was suggested as a fix for this limitation, which will actually be introduced in \fullref{sec:theory-optimizations-monte-carlo}.

In this work, the \emph{concurrence} was chosen as the \emph{entanglement monotone} to provide insight into the entanglement between different sites.
The concurrence has been introduced in \cite{concurrenceMainPaper} as a measure to obtain the entanglement of two spins from the reduced density matrix that contains these two spins.
A concurrence measurement does not give the value for the \emph{entanglement of formation} directly, but it needs to be rescaled as shown in \autoref{fig:entanglement}.
As one can see there, the rescaling does not introduce major additional features and for that reason the concurrence will be studied without rescaling to the entanglement.
The point of minimal entanglement (maximum separability) and maximal entanglement map to the same values and the intermediate values communicate the same qualitative information.

\begin{SCfigure}[2.0][htbp]
    \centering
    \includegraphics[width=0.35\textwidth]{plotgeneration/binary-entropy-function/binary-entropy-function.pdf}
    \caption{
        Plot of the function that connects the entanglement $\text{E}_{\psiOfT}$ of a subsystem with the concurrence of a subsystem $C(\psiOfT)$ in the relevant region $0\leq x \leq 1$.
        According to \cite{concurrenceMainPaper}: $\text{E}_{\psiOfT} = \text{E}(C(\psiOfT))$, with the function $\text{E}(x) = H(\frac{1}{2} + \frac{1}{2} \sqrt{1-x^2})$ and the \emph{binary entropy function} $H(x) = -\left[x \cdot \log_2(x) + (1-x) \cdot \log_2(1-x)\right]$.
        Comparing with the identity $y=x$, it shows that the values at the borders of the region $x=0$ and $x=1$ are the same for both functions. 
        The examined function $\text{E}(x)$ increases monotonically for the relevant input range and lies reasonably close to the identity. 
    }
    \label{fig:entanglement}
\end{SCfigure}

\filepath{\cite{selfCode}}{/calculation-helpers/concurrence/} holds multiple scripts used to validate the calculations in the following section give the same result, no matter which of the alternative implementation-/calculation-strategies are used.

\cite{concurrenceMainPaper} introduces the \emph{magic basis} that is shown in \autoref{eq:magic-basis} as opposed to the counting basis that was introduced in the beginning of \fullref{sec:theory-observables-density-matrix}. The magic basis consists of \emph{Bell states} with particular phases.

\begin{equation}
    \label{eq:magic-basis}
    \begin{split}
        \ketN[m_1] &= \frac{1}{\sqrt{2}} \left(\ketN[11] + \ketN[00]\right) \\
        \ketN[m_2] &= \frac{1}{\sqrt{2}} i \left(\ketN[11] - \ketN[00]\right) \\
        \ketN[m_3] &= \frac{1}{\sqrt{2}} i \left(\ketN[10] + \ketN[01]\right) \\
        \ketN[m_4] &= \frac{1}{\sqrt{2}} \left(\ketN[10] - \ketN[01]\right) \\
    \end{split}
\end{equation}

The Bell state is a construct tightly linked to entanglement, consequently it is by design straight forward, to write down the concurrence of a pure state expressed in the magic basis.
For $\ketN[\psiOfT] = \sum\limits_{i=1}^4 \alpha_i \ketN[m_i]$ the concurrence $C(\psiOfT)$ is $C(\psiOfT) = \left|\sum\limits_{i=1}^4 \alpha_i^2 \right|$.
This is equivalent to saying for a pure state expressed in the counting basis $\ketN[\psiOfT] = \alpha \ketN[00] +  \beta \ketN[01] +  \gamma \ketN[10] +  \delta\ketN[11]$, the concurrence is calculated $C(\psiOfT) = 2 \left|\alpha \delta - \beta \gamma\right|$.
This gives the same value as for all calculations that extract the concurrence from the density matrix that is calculated from these pure state coefficients.
But, while all states obtained via time evolution $\ketN[{\psiOfT[\schroedingerPicture]}]$ - as described in previous chapters - are pure states, through the reduction to a \glqq two-spin-system\grqq{} of particles $l,\,\sigma$ and $m,\,\mu$ to $\rho_{l,\,\sigma,\,m,\,\mu} (t)$ the remaining system is likely to be a mixed state \cite{mixedStateFromPureState}.

The calculation necessarily needed to be generalized to take the density matrix as an input.
\cite{concurrenceRewording} gives two alternate ways of calculating the concurrence from the density matrix.
Both require the \emph{spin-flip} operation $\tilde{\rho}$.
The original paper \cite{concurrenceMainPaper} describes this operation for density matrix $\rho$ expressed in the counting basis to be: transforming it into to magic basis of \autoref{eq:magic-basis}, taking the complex conjugate, then transforming it back.
As per \cite{concurrenceRewording} this is equivalent to the operation described in \autoref{eq:spin-flip}, where $\rho$ and \pauli[]{y} (\autoref{eq:pauli-matrices}) are expressed in the standard counting basis.

\begin{equation}
    \label{eq:spin-flip}
    \begin{split}
        \tilde{\rho} = \left(\pauli[]{y} \otimes \pauli[]{y}\right) \rho^\ast  \left(\pauli[]{y} \otimes \pauli[]{y}\right)
    \end{split}
\end{equation}

\filepath{\cite{selfCode}}{/calculation-helpers/concurrence/basis-transformation.py} verifies this equivalence and shows the ${\tilde{\ast}}$-operation modifies a hermitian matrix like shown in \autoref{eq:demonstration-spin-flip}.

\begin{equation}
    \label{eq:demonstration-spin-flip}
    \begin{split}
        \rho = \left(\begin{matrix}
            a& b & c &d \\
            b^\ast& e & f &g \\
            c^\ast& f^\ast & h &k \\
            d^\ast& g^\ast & k^\ast &l \\
        \end{matrix}\right)  \quad\Rightarrow\quad\tilde{\rho} = 
        \left(\begin{matrix}
            l& -k & -g &d \\
            -k^\ast& h & f &-c \\
            -g^\ast& f^\ast & e &-b \\
            d^\ast& -c^\ast & -b^\ast &a \\
        \end{matrix}\right) 
    \end{split}
\end{equation}

With that knowledge \autoref{eq:concurrence-calculation} states how to calculate the concurrence $C(\rho)$ from the reduced density matrix $\rho$.

\begin{equation}
    \label{eq:concurrence-calculation}
    \begin{gathered}
        R = \sqrt{\sqrt{\rho} \cdot \tilde{\rho} \cdot \sqrt{\rho}}\\
        \lambda_1 \geq \lambda_2 \geq \lambda_3 \geq \lambda_4, \quad \left\{\lambda_i\right\} = \text{Eigenvalues}(R)\\
        C(\rho) = \max (0, \lambda_1 - \lambda_2 - \lambda_3 - \lambda_4)
    \end{gathered}
\end{equation}

In this formula, the square root $\sqrt{M}$ of a square matrix $M$ is a square matrix of the same dimensionality, that fulfills the property $\sqrt{M}\cdot \sqrt{M} = M$.
A candidate can be obtained numerically, by calculating the eigenvalues of the matrix (which for a hermitian matrix are all positive real numbers), taking their square root and re-assembling the matrix with the eigenvectors like in \autoref{eq:matrix-square-root} ($\lambda_i$ different values than in \autoref{eq:concurrence-calculation}).

\begin{equation}
    \label{eq:matrix-square-root}
    \begin{gathered}
        M  = Q \Lambda Q^\text{T}\\
        \Lambda = \text{diag}(\lambda_1,\, \lambda_2,\, \lambda_3,\, \lambda_4)\\
        \sqrt{\Lambda} = \text{diag}(\sqrt{\lambda_1},\, \sqrt{\lambda_2},\,\sqrt{\lambda_3},\,\sqrt{\lambda_4})\\
        \sqrt{M}  = Q \sqrt{\Lambda} Q^\text{T}
    \end{gathered}
\end{equation}

While this operation is quite complicated, it is required only once per calculation and only for a matrix of fixed size $4\times 4$.
To avoid it, \cite{concurrenceRewording} alternatively provides that the $\lambda_i$ from \autoref{eq:concurrence-calculation} are the square roots of the non-hermitian matrix $\rho \tilde{\rho}$.
For exact calculations both ways of calculating yield the same results, which is implemented and checked. 
If there are any inconsistencies in the order of basis vectors however, the results diverge. Which in fact is a good check that the bases are consistent across the whole implementation.

        \FloatBarrier
        
        \subsection{Energy}
        \label{sec:theory-observables-energy}
        The energy \energyOft (\autoref{eq:energy-main}) is a key measure of many quantum-mechanical experiments, with the observable being the hamiltonian that governs the time-evolution.
The variance \varianceOft of this quantity (\autoref{eq:variance-main}) describes the energy-fluctuations in the system.

\begin{equation}
    \label{eq:energy-main}
    \energyOft = \frac{\bracketHelper{\psiOfT[\schroedingerPicture]}{\pictureHamiltonian[\schroedingerPicture]}{\psiOfT[\schroedingerPicture]}}{\braketHelper{\psiOfT[\schroedingerPicture]}{\psiOfT[\schroedingerPicture]}}
       = \timeExpectationVal{\pictureHamiltonian[\schroedingerPicture]}
\end{equation}
\begin{equation}
    \label{eq:variance-main}
    \varianceOft = \frac{\bracketHelper{\psiOfT[\schroedingerPicture]}{{\pictureHamiltonian[\schroedingerPicture]}^2}{\psiOfT[\schroedingerPicture]}}{\braketHelper{\psiOfT[\schroedingerPicture]}{\psiOfT[\schroedingerPicture]}}-\left(\frac{\bracketHelper{\psiOfT[\schroedingerPicture]}{\pictureHamiltonian[\schroedingerPicture]}{\psiOfT[\schroedingerPicture]}}{\braketHelper{\psiOfT[\schroedingerPicture]}{\psiOfT[\schroedingerPicture]}}\right)^2
       = \timeExpectationVal{{\pictureHamiltonian[\schroedingerPicture]}^2} - {\timeExpectationVal{\pictureHamiltonian[\schroedingerPicture]}}^2
\end{equation}

Without any additional constrictions (like e.g. conservation of particle number) the energy - as well as all other cumulants of the hamiltonian operator - must be conserved under normal time evolution, as \autoref{eq:energy-conserved} shows. This uses the fact that an exact hamiltonian commutes with itself in all instances and allows to show that the energy $\energyOft = \energyOft[t=0]$.
Because equivalent reasoning, all higher order cumulants must necessarily equal $0$ for all times, while the first cumulant % https://en.wikipedia.org/wiki/Cumulant (yes, energy is the first, not the zeroth)
is constant (the first order being 0 comes only down to the choice of energy reference).

\begin{equation}
    \label{eq:energy-conserved}
    \begin{split}
        \energyOft &\stackrel{\ref{eq:energy-main}}{=}
        \frac{\bracketHelper{\psiOfT[\schroedingerPicture]}{\pictureHamiltonian[\schroedingerPicture]}{\psiOfT[\schroedingerPicture]}}{\braketHelper{\psiOfT[\schroedingerPicture]}{\psiOfT[\schroedingerPicture]}}
        \stackrel{\ref{eq:time-evolution-schroedinger}}{=}
        \frac{\bracketHelper{\picturePsi[\schroedingerPicture]{}}{e^{i \pictureHamiltonian[\schroedingerPicture] t}{\pictureHamiltonian[\schroedingerPicture]}e^{-i \pictureHamiltonian[\schroedingerPicture] t}}{\picturePsi[\schroedingerPicture]{}}}{\bracketHelper{\picturePsi[\schroedingerPicture]{}}{e^{i \pictureHamiltonian[\schroedingerPicture] t}e^{-i \pictureHamiltonian[\schroedingerPicture] t}}{\picturePsi[\schroedingerPicture]{}}}\\
        &\stackrel{\phantom{\ref{eq:time-evolution-schroedinger}}}{=}
        \frac{\bracketHelper{\picturePsi[\schroedingerPicture]{}}{e^{i \pictureHamiltonian[\schroedingerPicture] t}e^{-i \pictureHamiltonian[\schroedingerPicture] t}{\pictureHamiltonian[\schroedingerPicture]}}{\picturePsi[\schroedingerPicture]{}}}{\bracketHelper{\picturePsi[\schroedingerPicture]{}}{e^{i \pictureHamiltonian[\schroedingerPicture] t}e^{-i \pictureHamiltonian[\schroedingerPicture] t}}{\picturePsi[\schroedingerPicture]{}}} = 
        \frac{\bracketHelper{\picturePsi[\schroedingerPicture]{}}{\pictureHamiltonian[\schroedingerPicture]}{\picturePsi[\schroedingerPicture]{}}}{\braketHelper{\picturePsi[\schroedingerPicture]{}}{\picturePsi[\schroedingerPicture]{}}}
         = \energyOft[t=0]
    \end{split}
\end{equation}

While it might not seem to give much insight to calculate the time evolved measurements of these values if they are constant, in reality this serves the purpose of \emph{checking} the accuracy of potential approximations.
If the time evolution is performed with an approximated, effective hamiltonian, then the exact one will no longer commute with it and the energy will vary over time and the variance will have nonzero readings.

Measuring \localObservable{N}{t} for $\ObservableOp = \pictureHamiltonian[\schroedingerPicture]$ is straight forward, as \autoref{eq:energy-local-observable} shows.

\begin{equation}
    \label{eq:energy-local-observable}
    \begin{split}
        \localObservable{N}{t} &\stackrel{\ref{eq:expectation-value}}{=}
        \lsum[K] \bracketHelper{N}{\pictureHamiltonian[\schroedingerPicture]}{K} e^{\HeffOft[K]-\HeffOft[N]}
        \frac{
            \psiN[K]
        }{
            \psiN[N]
        } \\
        &\stackrel{\ref{eq:main-hamiltonian-h0},\,\ref{eq:main-hamiltonian-perturbation}}{=}
        %
        \lsum[K] \bracketHelper{N}{U \cdot \lsum \nop{l}{\up}\nop{l}{\down} + \lsum[l,\,\sigma] \epsl \nop{l}{\sigma}}{K} e^{\HeffOft[K]-\HeffOft[N]}
        \frac{
            \psiN[K]
        }{
            \psiN[N]
        }
        %
        \\&\stackrel{\phantom{\ref{eq:base-energy}}}{-}
        %
        J \cdot 
        \lsum[K] \bracketHelper{N}{\neighborsumWSpin{l}{m}{\sigma} \left(\withspinhcop{l}{\sigma}{\dagger} \withspinhcop{m}{\sigma}{} + \withspinhcop{m}{\sigma}{\dagger} \withspinhcop{l}{\sigma}{} \right)}{K} e^{\HeffOft[K]-\HeffOft[N]}
        \frac{
            \psiN[K]
        }{
            \psiN[N]
        }
        \\&\stackrel{\phantom{\ref{eq:base-energy}}}{=}
        %
        \lsum \lsum[K] \bracketHelper{N}{U \cdot  \nop{l}{\up}\nop{l}{\down} + \epsl \left(\nop{l}{\up} + \nop{l}{\down}\right)}{K} e^{\HeffOft[K]-\HeffOft[N]}
        \frac{
            \psiN[K]
        }{
            \psiN[N]
        }
        %
        \\&\stackrel{\phantom{\ref{eq:base-energy}}}{-}
        %
        J \cdot \neighborsumWSpin{l}{m}{\sigma}
        \lsum[K] \bracketHelper{N}{ \withspinhcop{l}{\sigma}{\dagger} \withspinhcop{m}{\sigma}{} + \withspinhcop{m}{\sigma}{\dagger} \withspinhcop{l}{\sigma}{} }{K} e^{\HeffOft[K]-\HeffOft[N]}
        \frac{
            \psiN[K]
        }{
            \psiN[N]
        }
        \\&\stackrel{\phantom{\ref{eq:base-energy}}}{=}
        %
        \lsum \left(U \cdot  n_{l,\,\up}n_{l,\,\down} + \epsl \left(n_{l,\,\up} + n_{l,\,\down}\right)\right)
        %
        \\&\stackrel{\phantom{\ref{eq:base-energy}}}{-}
        %
        J \cdot \neighborsumWSpin{l}{m}{\sigma}
        \lsum[K] \bracketHelper{N}{ \withspinhcop{l}{\sigma}{\dagger} \withspinhcop{m}{\sigma}{} + \withspinhcop{m}{\sigma}{\dagger} \withspinhcop{l}{\sigma}{} }{K} e^{\HeffOft[K]-\HeffOft[N]}
        \frac{
            \psiN[K]
        }{
            \psiN[N]
        }\\
        %
        &\stackrel{\ref{eq:base-energy}}{=} E_0(N,\,t)-
        %
        J \cdot \neighborsumWSpin{l}{m}{\sigma}
        \left(n_{l,\,\sigma} (1- n_{m,\,\sigma}) + (1-n_{l,\,\sigma}) n_{m,\,\sigma}\right)
         e^{\HeffOft[\vphantom{N}\smash{\tilde{N}}]-\HeffOft[N]}
        \frac{
            \psiN[\vphantom{N}\smash{\tilde{N}}]
        }{
            \psiN[N]
        }\\
        %
        &\stackrel{\phantom{\ref{eq:base-energy}}}{=} E_0(N,\,t)-
        %
        J \cdot \neighborsumWSpin{l}{m}{\sigma}
        \left(n_{l,\,\sigma} \neq n_{m,\,\sigma}\right)
         e^{\HeffOft[\vphantom{N}\smash{\tilde{N}}]-\HeffOft[N]}
        \frac{
            \psiN[\vphantom{N}\smash{\tilde{N}}]
        }{
            \psiN[N]
        }\\
        %
        &\stackrel{\phantom{\ref{eq:base-energy}}}{=} E_0(N,\,t) + E_\text{V}(N,\,t)
    \end{split}
\end{equation}

\ketN[\vphantom{N}\smash{\tilde{N}}] is here the state that is obtained, when starting on state \ketN and performing a swapping of the occupations of sites $l,\,\sigma$ and $m,\,\sigma$.
On comparison with \autoref{eq:base-energy} one can see, that the definition for $E_0(N,\,t)$ is consistent for the expansion and the observable calculation.
Newly defined is the energy $E_\text{V}(N,\,t)$, originating from \Vhamiltonian[].
As hinted for the other observables earlier, the difference of hamiltonians with local modifications $\HeffOft[\vphantom{N}\smash{\tilde{N}}]-\HeffOft[N]$ can be evaluated in constant time.
However the energy is the first observable that requires \bigo{\text{\#}(\text{sites})} on its own to evaluate. This can not be reduced, however as will become apparent in the calculation of the variance, this should be an upper bound for observable-calculations.


        \FloatBarrier

        \subsection{Monte Carlo Sampling}
        \label{sec:theory-monte-carlo}
        Mathematically, the problem of finding the solution to such a quantum-mechanical many-body model is solved by the calculations provided up to this point. 
However when looking at \autoref{eq:expectation-value}, one can easily see why in practice this is not viable with larger systems. 
The calculation requires summing over all base-states. 
Their number is given in \autoref{eq:system-size} (factor 2 in the exponent for the spin-degree) to be of exponential size in regards to the number of lattice-sites.
Considering a physical material has at least in the order of $10^{23}$ (from \emph{Avogadros-Constant}) degrees of freedom, an \bigo{2^{\text{\#}(\text{sites})}} computational-complexity is not good enough to simulate interestingly-sized systems. 

\begin{equation}
    \label{eq:system-size}
    \text{\#}(\text{states}) = 2^{\text{\#}(\text{sites}) \cdot 2}
\end{equation}

Different strategies exist, to circumvent this limitation.
Here, the focus will be put onto randomized algorithms that use \emph{Monte-Carlo sampling} in order generate a limited number of states, representative of the real probability-distribution, in order to not having to sum over the complete \emph{Hilbert-Space}.
The strategy of solving problems of this kind with the strategy of \emph{variational Monte-Carlo} is e.g. suggested in \cite{metropolisAlgorithmAndVariationalMonteCarlo}.
This means starting with a random state and applying modifications to it (= variation).
When the accepting/rejecting of these modifications is coupled to the probability the state has, based on its energy in the system's energy-distribution, the resulting states will have an energy/probability distribution, that resembles the one they had if they were sampled from the full Hilbert-Space.
The derivation for the whole process is shown in \cite{monteCarloObservableSampling}.

However computing the normalized probability for the states is still required, in order to compare against a (e.g. thermal) reference-distribution.
As one sees in the probability in \autoref{eq:probability-isolated}, that was defined in \autoref{eq:expectation-value}, still a full sum over all states is required to get the normalization factor.
\begin{equation}
    \label{eq:probability-isolated}
    \probabilityOf{N}{t} = \frac{
        \absSquare{e^{\HeffOft}} \absSquare{\psiN} 
    }{
        \lsum[K] \absSquare{e^{\HeffOft[K]}} \absSquare{\psiN[K]} 
    }
\end{equation}
However by employing the \emph{Metropolis-Hastings Algorithm} instead of only the \emph{Metropolis-Algorithm} to decide on the acceptance/rejection of proposed states, one can drop the requirement of normalizing the probability.
In \autoref{eq:transition-probability-final} the \emph{acceptance-ratio} $\alpha$ - that is necessary for the Metropolis-Hastings Algorithm to work - is derived (as per \cite{metropolisHastingsAlgorithmGeneral}) and decide on the transition from state \ketN[N] to \ketN[{\vphantom{N}\smash{\tilde{N}}}].
The algorithm proposes to generate a new state \ketN[{\vphantom{N}\smash{\tilde{N}}}] and calculate the corresponding $\alpha$. Then it requires generating a random number $u$ between $0$ and $1$, pulled from a uniform distribution.
If $u \leq \alpha$, \ketN[{\vphantom{N}\smash{\tilde{N}}}] is accepted and will be used as the new current state in the next step of the sampling. 
If $u>\alpha$, the proposed state is rejected and \ketN[N] will be used again as the next state.
The function $f(N,t)$, provided in \autoref{eq:proportional-function}, fulfills the necessary property of being proportional to the probability  \probabilityOf{N}{t}.

\begin{equation}
    \label{eq:proportional-function}
    f(N,t) = \absSquare{e^{\HeffOft}} \absSquare{\psiN} \propto \probabilityOf{N}{t}
\end{equation}

\begin{equation}
    \label{eq:transition-probability-final}
    \begin{split}
        \alpha &\stackrel{\phantom{\ref{eq:proportional-function}}}{=} \frac{\probabilityOf{\vphantom{N}\smash{\tilde{N}}}{t}}{\probabilityOf{N}{t}} =  \frac{f(\vphantom{N}\smash{\tilde{N}},t)}{f(N,t)}
        \stackrel{\ref{eq:proportional-function}}{=}
        \frac{
            \absSquare{e^{\HeffOft[\vphantom{N}\smash{\tilde{N}}]}} \absSquare{\psiN[\vphantom{N}\smash{\tilde{N}}]}
        }{
            \absSquare{e^{\HeffOft}} \absSquare{\psiN}
        }\\
        &\stackrel{\phantom{\ref{eq:proportional-function}}}{=}
        \frac{\absSquare{\psiN[\vphantom{N}\smash{\tilde{N}}]}}{\absSquare{\psiN}}
        \frac{
            e^{\real\left(\HeffOft[\vphantom{N}\smash{\tilde{N}}]\right) + i\imaginary\left(\HeffOft[\vphantom{N}\smash{\tilde{N}}]\right) +\real\left(\HeffOft[\vphantom{N}\smash{\tilde{N}}]\right)-i \imaginary\left(\HeffOft[\vphantom{N}\smash{\tilde{N}}]\right)}
        }{
            e^{\real\left(\HeffOft\right) + i\imaginary\left(\HeffOft\right) +\real\left(\HeffOft\right)-i \imaginary\left(\HeffOft\right)}
        }\\
        &\stackrel{\phantom{\ref{eq:proportional-function}}}{=}
        \frac{\absSquare{\psiN[\vphantom{N}\smash{\tilde{N}}]}}{\absSquare{\psiN}}
        e^{2\cdot \real\left(\HeffOft[\vphantom{N}\smash{\tilde{N}}]\right) - 2\cdot \real\left(\HeffOft\right)}\\
        &\stackrel{\phantom{\ref{eq:proportional-function}}}{=}
        \frac{\absSquare{\psiN[\vphantom{N}\smash{\tilde{N}}]}}{\absSquare{\psiN}}
        e^{2\cdot \real\left(\HeffOft[\vphantom{N}\smash{\tilde{N}}] - \HeffOft\right)}
    \end{split}
\end{equation}

Multiple factors must be considered, which kind of variation is applied to the states in this schema. 
However probably the most important criteria is how quickly one can evaluate \autoref{eq:transition-probability-final}, as this will be done tens to hundreds of times between each sampled state.

Finally, the method of sampling a selection of distributed states results in \autoref{eq:final-sampling-equation} as a rewrite to \autoref{eq:expectation-value} \cite{monteCarloObservableSampling}.

\begin{equation}
    \label{eq:final-sampling-equation}
    \frac{\bracketHelper{\psiOfT[\schroedingerPicture]}{\ObservableOp}{\psiOfT[\schroedingerPicture]}}{\braketHelper{\psiOfT[\schroedingerPicture]}{\psiOfT[\schroedingerPicture]}} \stackrel{\ref{eq:expectation-value}}{=} 
    \lsum[N]
    \probabilityOf{N}{t}
    \localObservable{N}{t} \cdot
    \stackrel{\text{\cite{monteCarloObservableSampling}}}{\approx} \frac{1}{\left|\left\{N\right\}_\text{MC}\right|} \lsum[\left\{N\right\}_\text{MC}]\localObservable{N}{t}
\end{equation}
        \FloatBarrier
        
    \section{Time-Complexity Optimizations}
    \label{sec:theory-optimizations}
    When performing time-evolution with a perturbative approach, one ought to watch out for resonances.
In the case of degenerate energy solutions, so-called \emph{secular terms} \cite{secularTermsPerturbation} might be introduced, if non-degenerate perturbation theory is used to find the solution to a degenerate problem. 
Often inaccuracies manifest as linear terms in a function where only oscillating terms would be expected.
The linear term constantly grows over time, guaranteeing the approximation is only valid for short times and then diverges further and further from the exact solution.

For the examined Hamiltonian the first occurrence of such terms can be spotted in the second order cumulant expansion of the Hamiltonian.
\autoref{eq:time-ordered-integral-part1} and \autoref{eq:time-ordered-integral-part2} specifically list extra solutions for the cases where the two energies are the same.
This resonance causes a linear term to appear in the expansion, while all other pre-factors so far were complex exponentials -- meaning they are oscillating.

The idea to combat this is to replace the perturbative parameters that were calculated analytically with variational parameters instead.
This has the motive to generate a physically motivated structure, but the fine details of the parameter weights are adjusted by optimization.
E.g. for the \emph{neural-network quantum state} \cite{neuralNetworkQuantumStates} approach, a neural network structure is used to represent a wave-function.
Reason for this being the hope that the exponentially complicated wave-function can be expressed as a \glqq higher truth\grqq{} only requiring few parameters to describe a seemingly complicated wave-function.
While it would be basically impossible to always find the few-parameter-solution, \emph{approximating} it close enough with optimization akin to machine-learning should be a suitable replacement.

The following section aims to achieve better approximations than the ones found by cumulant expansion by letting the parameters be optimized during the step to reflect the situation more closely.
    \FloatBarrier

        \subsection{Analytical Simplifications for Special Cases}
        \label{sec:theory-optimizations-analytical}
        Generally, handling the computation of the difference of the effective Hamiltonian for two states $\HeffOft[\tilde{N}]-\HeffOft[N]$ is the most expensive computation per Monte-Carlo-step.
It is required in \autoref{eq:form-heff-difference} for the calculation of the spin-polarized kinetics.
As stated in \fullref{sec:theory-monte-carlo} it is also needed to compute the transition probability between two sampled states in \autoref{eq:transition-probability-final}.
The calculation of $\HeffOft[N]$ on its own requires an effort of \bigo{\text{\#}(\text{sites}) \cdot \text{\#}(\text{nearest neighbors})}.
However, for the states \ketN[N] and \ketN[{\vphantom{N}\smash{\tilde{N}}}] that are connected with only small changes, most of the elements in the sum cancel and the complexity may be here even reduced to \bigo{$\text{\#}(\text{nearest neighbors})$}.
This is quite attractive, because it de-couples the per-step computational costs from the number of lattice sites - a vital step for simulating large systems efficiency - however this requires extra analytical calculations.
In this section analytical simplifications are presented to get the objects $E_0(N)-E_0(\tilde{N})$ and $\HNOft - \HNOft[\tilde{N}]$ (Caution, do not miss the extra \emph{minus} required, to calculate the required $\HeffOft[\tilde{N}]-\HeffOft[N]$ from this!).

\paragraph*{Initial State} \makebox{}\\

The initial state of the system before the start of the time-evolution is encoded in $\psiN$, defined by \autoref{eq:base-expansion-state}. Different configurations can be considered here.

One natural choice would be the preparation of a single defined base-state which would result in all $\psiN$ being $0$ but one that would have the value of $1$. 
This case needs to be considered with a specialized analytical calculation, because there are many instances in this calculation, where a division by $\psiN$ is performed. 

However the current calculation is compatible with the concept of one base-state having a probability $y$ that is $x$ times (e.g. $x=100$-times) as large as the probability $z = y/x$ of all the other states.
This effectively also concentrates the initial state to one base-state.

With the additional restriction of normalization from \autoref{eq:normalization} the individual values for \psiN{} can be calculated (see \autoref{eq:psi-n-concentrated}).

\begin{equation}
    \label{eq:normalization}
    \lsum[N] \psiN{}^2 = 1
\end{equation}

\begin{equation}
    \label{eq:psi-n-concentrated}
    \begin{split}
        \stackrel{\ref{eq:normalization}}{\Rightarrow} y^2 + \left[\text{\#}(\text{states}) - 1 \right] z^2 = 1\\
        y^2 + \left[\text{\#}(\text{states}) - 1 \right] \frac{y^2}{x^2} = 1\\
        y = \sqrt{\left(1+\left[\text{\#}(\text{states}) - 1 \right]\frac{1}{x^2}\right)^{-1}}
    \end{split}
\end{equation}

This configuration was tested and verified for small systems with summation over the full Hilbert-space. 
However when Monte-Carlo sampled, the configuration produced different results than for exact sampling.
This can probably be explained by the fact that the energy landscape of base-states for this configuration is not smooth enough. 
Causing the Metropolis-Algorithm not to perform as intended for the number of samples.
As stated, treating non-uniform distribution requires a different analytical approach. Therefore no further experimentation on this was done.\\

A second natural choice would be the perfectly uniform distribution of probability for all the base-states. 
This results in a smooth energy-landscape, that can be properly sampled. 

The calculation for this is straight-forward and listed in \autoref{eq:psi-n-homogenous}

\begin{equation}
    \label{eq:psi-n-homogenous}
    \psiN{} = \frac{1}{\sqrt{\text{\#}(\text{states})}} \stackrel{\ref{eq:system-size}}{=} \frac{1}{\sqrt{2^{\text{\#}(\text{sites}) \cdot 2}}} = \frac{1}{2^{\text{\#}(\text{sites})}}
\end{equation}

One last advantage of the uniform distribution of base-state-probabilities is, that all $\psiN$ are equal, which makes it possible to optimize sums, as $\psiN{}/\psiN[\tilde{N}]{} = 1$.
The following two calculations require this assumption.

As the states under consideration are \ketN[N] and \ketN[{\vphantom{N}\smash{\tilde{N}}}], the occupation-numbers $n_{l,\,\sigma}$ will describe the occupation of \ketN[N] and $\tilde{n}_{l,\,\sigma}$ the occupation of \ketN[{\vphantom{N}\smash{\tilde{N}}}].
Both the following modifications restrict the values $\tilde{n}_{l,\,\sigma}$ to make them dependent on $n_{l,\,\sigma}$.

\newpage 
\paragraph*{Hopping-only modification} \makebox{}\\

In this simplification, the maximum difference that can happen, is that a \emph{hopping} event occurs from site $i$ and spin $\sigma_i$ to the site $j$ and spin $\sigma_j$. 
This can only happen if $n_{i,\,\sigma_i} \neq n_{j,\,\sigma_j}$, because if they are the same, the swapping doesn't change the state. 
This check can be performed beforehand, to save on computational resources (if they are equal, $\HeffOft[\tilde{N}]-\HeffOft[N] = 0$).
Because the inverse (hopping backwards) is also an equivalently valid change of state, this really looks at the \emph{swapping} of the two occupations.

The hopping results in the values for $\tilde{n}_{l,\,\sigma}$ are given in \autoref{eq:new-n-hopping}.

\begin{equation}
    \label{eq:new-n-hopping}
    \tilde{n}_{l,\,\sigma} = \begin{cases}
        n_{i,\,\sigma_i}&\text{ : } l = j \land \sigma = \sigma_j   \\
        n_{j,\,\sigma_j}&\text{ : } l = i \land \sigma = \sigma_i   \\
        n_{l,\,\sigma} &\text{ : else}
    \end{cases}
\end{equation}

This simplifies the energy difference $E_0(N)-E_0(\tilde{N})$ like equation \autoref{eq:simplified-base-energy-hopping} presents.

\begin{equation}
    \label{eq:simplified-base-energy-hopping}
    \begin{split}
        E_0(N)-E_0(\tilde{N}) 
        &\stackrel{\phantom{\ref{eq:new-n-hopping}}}{=} U\lsum n_{l,\,\down}n_{l,\,\up}-U\lsum \tilde{n}_{l,\,\down}\tilde{n}_{l,\,\up} 
        + \lsum[l,\,\sigma] \epsl n_{l,\,\sigma} - \lsum[l,\,\sigma] \epsl \tilde{n}_{l,\,\sigma}\\
        &\stackrel{\ref{eq:new-n-hopping}}{=} \left(\epsl[i]-\epsl[j]\right)\left(n_{i,\,\sigma_i} - n_{j,\,\sigma_j}\right) +
        U\cdot \begin{cases}
            \left(n_{i,\,\up}-n_{j,\,\up}\right) \left(n_{i,\,\down}-n_{j,\,\down}\right) &\text{ : } \sigma_i = \sigma_j   \\
            \left(n_{i,\,\up}-n_{j,\,\down}\right) \left(n_{i,\,\down}-n_{j,\,\up}\right) &\text{ : } \sigma_i \neq \sigma_j
        \end{cases}\\
        &\stackrel{\phantom{\ref{eq:new-n-hopping}}}{=} \left(\epsl[i]-\epsl[j]\right)\left(n_{i,\,\sigma_i} - n_{j,\,\sigma_j}\right) +
        U \left(n_{i,\,\sigma_i}-n_{j,\,\sigma_j}\right) \left(n_{i,\,\overline{\sigma_i}}-n_{j,\,\overline{\sigma_j}}\right)
    \end{split}
\end{equation}

The simplification of $\HNOft - \HNOft[\tilde{N}]$ can not be calculated as quickly as the difference of base-energies.
Primarily this is because of the number of terms in this calculation. Each of the three \VhamiltonianPartA{N}{t}, \VhamiltonianPartB{N}{t} and \VhamiltonianPartC{N}{t} is a sum of 8 summands, as one can see in \autoref{eq:interaction-picture-v-ham-parts}.
Each of those must be treated for for cases: $(\sigma_i = \up \land \sigma_j = \up)$, $(\sigma_i = \up \land \sigma_j = \down)$, $(\sigma_i = \down \land \sigma_j = \up)$ and $(\sigma_i = \down \land \sigma_j = \down)$.
This makes 32 differently-shaped terms, or combined with the pre-factors 96 terms to fully write down.

Different notation could be used, to lessen the number of calculations, but this seemed unnecessarily complex.

A simpler solution was in fact, to let a python script generate all the final simplified analytical terms and output them as usable functions into a python script to include in the rest of the code.
The generating script can be found here: 

\filepath{\cite{selfCode}}{/calculation-helpers/simplificationtermhelper.py}

and the output with the analytically fully optimized calculations in a unified format here:
\filepath{\cite{selfCode}}{/computation-scripts/analyticalcalcfunctions.py}.

As this needs to be run quite often, it is useful to reduce the number of calculations as much as possible and avoid running unnecessary calculations on every iteration (e.g. $n_{l,\,\up} + 1 - (1 + n_{l,\,\up})$ which is definitely 0, no matter the value of $n_{l,\,\up}$).
For this reason, the terms were pre-optimized as much as possible, by trying to pre-evaluate to definite zeros or constants terms. 
This was performed automatically with the python package \emph{SymPy} \cite{sympyPackage}.

More details on how this was calculated can be found in \fullref{sec:implementation-analytical-calculations} and further resources, that that section redirects towards.

\paragraph*{Flipping-only modification} \makebox{}\\

In this simplification, the maximum difference that can happen, is that a \emph{flipping} event occurs on site $i$ and spin $\sigma_i$. 

The flipping results in the values for $\tilde{n}_{l,\,\sigma}$ are given in \autoref{eq:new-n-flipping}.

\begin{equation}
    \label{eq:new-n-flipping}
    \tilde{n}_{l,\,\sigma} = \begin{cases}
        (1 - n_{i,\,\sigma_i})&\text{ : } l = i \land \sigma = \sigma_i   \\
        n_{l,\,\sigma} &\text{ : else}
    \end{cases}
\end{equation}

This simplifies the energy difference $E_0(N)-E_0(\tilde{N})$ like equation \autoref{eq:simplified-base-energy-flipping} presents.

\begin{equation}
    \label{eq:simplified-base-energy-flipping}
    \begin{split}
        E_0(N)-E_0(\tilde{N}) 
        &\stackrel{\phantom{\ref{eq:new-n-flipping}}}{=} U\lsum n_{l,\,\down}n_{l,\,\up}-U\lsum \tilde{n}_{l,\,\down}\tilde{n}_{l,\,\up} 
        + \lsum[l,\,\sigma] \epsl n_{l,\,\sigma} - \lsum[l,\,\sigma] \epsl \tilde{n}_{l,\,\sigma}\\
        &\stackrel{\ref{eq:new-n-flipping}}{=} \epsl[i] \left(2 n_{i,\,\sigma_i} - 1\right) +
        U\cdot \begin{cases}
            n_{i,\,\down} (2 n_{i,\,\up} - 1) &\text{ : } \sigma_i = \up   \\
            n_{i,\,\up} (2 n_{i,\,\down} - 1) &\text{ : } \sigma_i = \down 
        \end{cases}
    \end{split}
\end{equation}

Like for the hopping problem, the calculation of $\HNOft - \HNOft[\tilde{N}]$ is more involved in the number of terms.

More details on how this was calculated can be found in \fullref{sec:implementation-analytical-calculations} and further resources, that that section redirects towards.

\vspace{1cm}
All presented simplifications to the Hamiltonian and corresponding MC-modification strategies in this section are implemented as options that are selectable in the control-script   

\filepath{\cite{selfCode}}{/computation-scripts/script.py}.
 
        \FloatBarrier

        \subsection{Geometry-Dependent Interaction-Range}
        \label{sec:theory-optimizations-geometry}
        
\def\getFromArray#1{%
    \ifcase#1 0\or
    1\or
    2\or
    \dots\or
    (M-1)\or
    M\or
    M+1\or
    M+2\or
    \dots\or
    M+(M-1)\or
    2\cdot M\or
    2\cdot M+1\or
    2\cdot M+2\or
    \dots\or
    2\cdot M+(M-1)\or
    \dots\or
    \dots\or
    \dots\or
    \dots\or
    \dots\or
    (M-1)\cdot M\or
    M^2-M+1\or
    M^2-M+2\or
    \dots\or
    M^2 -1\or
    \else\fi
}


\begin{tikzpicture}[scale=0.38]
    \def\m{5} % Change this value to adjust the grid size (m x m)
    \def\side{6} % Change this value to adjust the square side length
    \def\labelsize{\side/6} % Adjust label size
    
    % Draw dots and labels with custom text using a loop
    \foreach \x in {0,...,\numexpr\m-1\relax} {
        \foreach \y in {0,...,\numexpr\m-1\relax} {
            \pgfmathtruncatemacro{\index}{(\m-\y-1) * \m + \x}
            \node[scale=\labelsize] at (\x*\side + \side/2,\y*\side + \side/2) {$\getFromArray{\index}$};
            \fill[color=dblue] (\x*\side + \side/2,\y*\side + \side/2 + \side/4) circle (0.4);
        }
    }

    % Store coordinates for reference
    \coordinate (top) at (current bounding box.north west);

    % Draw horizontal black line
    \draw[line width=2pt,overlay,lightgray,dashed] (top) ++ (0,-0.55*\side) -- ++(\m*\side,0);
        
    % Draw green arrow pointing downwards
    \draw[green, ->, overlay, line width=2pt,>=stealth] (top) ++ (1.3*\side,-0.55*\side) -- ++(3*\side,-3*\side) node[pos=0.05,right] {$\varphi$};

    % Draw angle circle
    \draw[overlay, green, line width=1.5pt] (top) ++ (1.3*\side,-0.55*\side) +(0:0.5*\side) arc (0:-45:0.5*\side);
\end{tikzpicture}
        \FloatBarrier

    \section{Variational Classical Networks}
        \cite{neuralNetworkQuantumStates} % neural quantum states general
        
        \subsection{General Theory of VCN}
        \cite{originalDerivationTimeDependendVariationalPrinciple} % tdvp historic source
        \cite{TDVPcomplexPrefactors} % how to treat complex pre-factors, complex differentiation avoided through real differentiation
        \cite{complexDifferentiation} % in this case the complex derivative is fine
        \cite{probabilitySamplingRequirementVCN} % VCN general derivation, eta dot derivation and that it needs to sample from the probability distribution of the current state(t)
        \cite{VCNsolutionForRBM} % VCN solution completely derived for a RBM

        \subsection{Application to this Task}
        

\chapter{Implementation and Calculation Details}
\label{sec:implementation-details}
When performing time-evolution with a perturbative approach, one ought to watch out for resonances.
In the case of degenerate energy solutions, so-called \emph{secular terms} \cite{secularTermsPerturbation} might be introduced, if non-degenerate perturbation theory is used to find the solution to a degenerate problem. 
Often inaccuracies manifest as linear terms in a function where only oscillating terms would be expected.
The linear term constantly grows over time, guaranteeing the approximation is only valid for short times and then diverges further and further from the exact solution.

For the examined Hamiltonian the first occurrence of such terms can be spotted in the second order cumulant expansion of the Hamiltonian.
\autoref{eq:time-ordered-integral-part1} and \autoref{eq:time-ordered-integral-part2} specifically list extra solutions for the cases where the two energies are the same.
This resonance causes a linear term to appear in the expansion, while all other pre-factors so far were complex exponentials -- meaning they are oscillating.

The idea to combat this is to replace the perturbative parameters that were calculated analytically with variational parameters instead.
This has the motive to generate a physically motivated structure, but the fine details of the parameter weights are adjusted by optimization.
E.g. for the \emph{neural-network quantum state} \cite{neuralNetworkQuantumStates} approach, a neural network structure is used to represent a wave-function.
Reason for this being the hope that the exponentially complicated wave-function can be expressed as a \glqq higher truth\grqq{} only requiring few parameters to describe a seemingly complicated wave-function.
While it would be basically impossible to always find the few-parameter-solution, \emph{approximating} it close enough with optimization akin to machine-learning should be a suitable replacement.

The following section aims to achieve better approximations than the ones found by cumulant expansion by letting the parameters be optimized during the step to reflect the situation more closely.
\FloatBarrier

    \section{Math-Manipulator}
    \cite{compilersDragonBook}

    \section{Programming Style Optimizations}

    \section{Python-Script Generation}

    \section{Numerical Alternatives}

    \section{Simplification Verifications}

\chapter{Numerical Experiments}
\label{sec:numerical-experiments}
When performing time-evolution with a perturbative approach, one ought to watch out for resonances.
In the case of degenerate energy solutions, so-called \emph{secular terms} \cite{secularTermsPerturbation} might be introduced, if non-degenerate perturbation theory is used to find the solution to a degenerate problem. 
Often inaccuracies manifest as linear terms in a function where only oscillating terms would be expected.
The linear term constantly grows over time, guaranteeing the approximation is only valid for short times and then diverges further and further from the exact solution.

For the examined Hamiltonian the first occurrence of such terms can be spotted in the second order cumulant expansion of the Hamiltonian.
\autoref{eq:time-ordered-integral-part1} and \autoref{eq:time-ordered-integral-part2} specifically list extra solutions for the cases where the two energies are the same.
This resonance causes a linear term to appear in the expansion, while all other pre-factors so far were complex exponentials -- meaning they are oscillating.

The idea to combat this is to replace the perturbative parameters that were calculated analytically with variational parameters instead.
This has the motive to generate a physically motivated structure, but the fine details of the parameter weights are adjusted by optimization.
E.g. for the \emph{neural-network quantum state} \cite{neuralNetworkQuantumStates} approach, a neural network structure is used to represent a wave-function.
Reason for this being the hope that the exponentially complicated wave-function can be expressed as a \glqq higher truth\grqq{} only requiring few parameters to describe a seemingly complicated wave-function.
While it would be basically impossible to always find the few-parameter-solution, \emph{approximating} it close enough with optimization akin to machine-learning should be a suitable replacement.

The following section aims to achieve better approximations than the ones found by cumulant expansion by letting the parameters be optimized during the step to reflect the situation more closely.
\FloatBarrier

    \section{Time-Complexity Verification}

    \section{Perturbation-Approximation Convergence}

    \cite{variationalClassicalNetworksPaper}
    % TODO -> time order to where this should converge in VCN for dynamics paper
    % Let us remark that the cumulant expansion (11) goes
    % beyond conventional time-dependent perturbation the-
    % ory, since the corrections considered here effectively ac-
    % count for a resummation of several terms that appear in a
    % standard perturbative expansion31 . However, pCNs face
    % their own limitations, too. In particular, they are inher-
    % ently restricted to weak quantum fluctuations (small γ)
    % to ensure that we can safely truncate the expansion (11).
    % Besides, the description of the evolution of observables
    % will eventually break down, since resonant processes may
    % be present, giving rise to secular terms that limit a correct
    % description to timescales of order O(1/γ)31 . Nonetheless,
    % one can still benefit from the framework introduced here,
    % while mitigating the drawbacks mentioned before. This
    % is achieved by constructing adequate variational wave
    % functions with a network architecture that is inherited
    % from a corresponding pCN, as argued in the following.
    
    \cite{starkManyBodyLocalization} % what we want to observe for electrical field on a discrete lattice: fluctuations instead of transport

    \section{Monte-Carlo-Sampling Convergence}

    \section{Observable-Dependent Noise}

    \section{Improvements through VCN}

    \section{System-Size Dependency of Measurements}

\chapter{Conclusion}
\label{sec:conclusion}
While working on this research subject over the course of my practical training, my project work and finally the master thesis, the material turned out to be more challenging than expected.
Initially, work advanced with steady progress and clear goals.
Because of the seemingly very complicated commutation relations, I spent a lot of time on the development of the Math-Manipulator to streamline similar work.
I had the feeling that several people at the chair would benefit from a tool to validate and improve the calculations that are traditionally done on paper, without resorting to the full adoption of a computer-algebra-system.
Looking back at the calculations eventually needed for the thesis, this tool seems to have been under-utilized.

However, during the development I was able to significantly improve my understanding for the different types of operators, their properties and differences.
This not only directly lead to the calculations now being as straight forward as they are, but also - as my understanding increased - numerous recalculations were required.
In that context, having this consistent tool saved many hours.
So I can confidently say that the time for the side-project was not wasted but well invested.

Consequently, when it came down to practical realization of the mathematical derivations, I spent a lot of time writing extra alternative calculations and programmed checks to verify the correctness of the analytical calculations and numerical implementations.
This lead to many bugs and quirks being found - many of which with only small effects on the result. 
Still, this strategy (as opposed to taking the first \glqq working\grqq{} code) made me highly confident to state that the final implementation is bug-free and correctly reflects the described methods accurately to a high degree.
This would be especially important, as a logical next step would be to produce an even more efficient implementation in a lower-level, compiled language like C(++) or Rust to improve base level performance.

While this thesis and the corresponding code aspires to provide adequate comparisons between the methods and validate their viability, because of the amount of inefficient small operations in Python it lacks the efficiency required for taking on experiments on very large systems.
This shows that the described computational complexity class can be reached and the method can be used to handle the underlying problem - yet there still are a lot of micro-optimizations left to tackle.
Also, several generalizations needed in previous revisions/comparisons could be taken out, if but the final iteration of the method is needed.

Getting all processes to work as intended, turned out to be more challenging than first anticipated.
For one, a problem in two dimensions and non-periodic boundary conditions has a lot more edge-cases than the same method on a one-dimensional, linear chain with periodic boundary conditions.
When tracking which interactions cancel, the overlapping \glqq spheres of influence\grqq{} produce a significantly higher number of interaction terms in this problem compared to descriptions of previous applications to simpler geometric lattices.

This, combined with two unexpectedly complex road-blocks caused my available time to erode.
Nearing the end of the allocated time-frame, the work was still containing less \glqq Physics\grqq{} (Concrete new insights, obtained by experimentally observing the effects of parameter variations) than we were hoping.
In the words of my supervisor Markus Heyl, in the cases where we can not imagine what is stopping the experiment from behaving as expected: \glqq That's science!\grqq{}.
Now at the end I can proudly state, that this motivated me to finally solve all remaining road-blocks before the deadline.

Which also is the reason I am happy to be able to provide a fully tested and detailed implementation and verification of the hypothesized solution.
To provide the biggest contribution to science, I believe in providing clean final code and transparent plotting process for making numerical methods reproducible and expandable.
And I hope the code can be build upon in future research, that may focus on honing the experiment's parameters and not being held up with debugging while replicating theoretical methods.

\FloatBarrier

% ! Addendum
\newpage
\nocite{*}

{
    \setlength{\emergencystretch}{5em} % should stop urls from going over the right side of the textwidth
    \sloppy                             % looser line-breaking rules
    \setcounter{biburllcpenalty}{9000}  % Try not to break urls at lower letters / in words (preferably more at /:-+= or other signs)

    \printbibliography[title={Bibliography}]
}
\chapter{Appendix}
\label{sec:appendix}

% minted to properly import and style code. ! Needs Python libraries

\section{Code-Optimization: Default} % X.1
    \label{code:if-optimized}
    \filepath{\cite{selfCode}}{/calculation-helpers/timing-efficiency/analyticallogicalfunctions.py}
    \inputminted[firstline=61, lastline=92]{python}{./../physics-code/calculation-helpers/timing-efficiency/analyticallogicalfunctions.py}

\newpage
\section{Code-Optimization: If-free} % X.2
    \label{code:if-free-optimized}
    \filepath{\cite{selfCode}}{/calculation-helpers/timing-efficiency/analyticallogicalfunctions.py}
    \inputminted[firstline=5, lastline=20]{python}{./../physics-code/calculation-helpers/timing-efficiency/analyticallogicalfunctions.py}

\newpage
\section{Code-Optimization: Espresso} % X.3
    \label{code:logically-optimized}
    \filepath{\cite{selfCode}}{/calculation-helpers/timing-efficiency/logicalcalcfunctions.py}
    \inputminted[firstline=6, lastline=27]{python}{./../physics-code/calculation-helpers/timing-efficiency/logicalcalcfunctions.py}
    
\markboth{\leftmark}{\thechapter.4\,\,\,Small Interaction-Term Limit: Extra Plots}
\begin{minipage}{\textwidth}
    \section{Small Interaction-Term Limit: Extra Plots} % X.4
        \label{appendix:more-jsweep-plots}
        Referenced by \autoref{fig:j-sweep-single-occ-center}.\\
        \vbox{
            \makebox[\textwidth][c]{
                \includegraphics[width=0.85\textwidth]{plotgeneration/j-sweep/current-border.pdf}
            }
            \vspace{-0.8cm}

            \makebox[\textwidth][c]{
                \includegraphics[width=0.85\textwidth]{plotgeneration/j-sweep/current-center.pdf}
            }
            \vspace{-0.8cm}

            \makebox[\textwidth][c]{
                \includegraphics[width=0.85\textwidth]{plotgeneration/j-sweep/double-occ-border.pdf}
            }
        }
\end{minipage}

\markboth{\leftmark}{\thechapter.5\,\,\,Density Matrix Measurement: Pauli Observables}
\begin{minipage}{\textwidth}
    \section{Density Matrix Measurement: Pauli Observables} % X.5
        \label{appendix:pauli-measurements}
        Referenced by \autoref{fig:concurrence-first-experiment}.\\
        \vbox{
            \makebox[\textwidth][c]{
                \includegraphics[width=1.0\textwidth]{plotgeneration/concurrence-from-spin/pauli-measurements.pdf}
            }
        }
\end{minipage}

\markboth{\leftmark}{\thechapter.6\,\,\,Monte-Carlo-Sampling Convergence: Current}
\begin{minipage}{\textwidth}
    \section{Monte-Carlo-Sampling Convergence: Current} % X.6
        \label{appendix:mc-convergence-spin-current}
        Referenced by \autoref{fig:mc-convergence-occupation}.\\
        \vbox{
            \makebox[\textwidth][c]{
                \includegraphics[width=1.0\textwidth]{plotgeneration/monte-carlo-variance-test/spin-current.pdf}
            }
        }
\end{minipage}

\markboth{\leftmark}{\thechapter.7\,\,\,VCN Problem: Parameters for Step Size}
\begin{minipage}{\textwidth}
    \section{VCN Problem: Parameters for Step Size} % X.7
        \label{appendix:vcn-fail-parameters}
        Referenced by \autoref{fig:vcn-fail-energy-variance}.\\
        \vbox{
            \makebox[\textwidth][c]{
                \includegraphics[width=1.0\textwidth]{plotgeneration/vcn-eff-stepsize/vcn-parameters.pdf}
            }
        }
\end{minipage}

\markboth{\leftmark}{\thechapter.8\,\,\,VCN Problem: Parameters for Starting Values}
\begin{minipage}{\textwidth}
    \section{VCN Problem: Parameters for Starting Values} % X.8
        \label{appendix:vcn-fail-starting-parameters}
        Referenced by \autoref{fig:vcn-fail-starting-energy-variance}. Note the imaginary part (represented by the dotted line) of parameters 1,\,2,\,4 and 5 diverging nearly instantly.\\
        \vbox{
            \makebox[\textwidth][c]{
                \includegraphics[width=1.0\textwidth]{plotgeneration/vcn-starting-values/vcn-parameters.pdf}
            }
        }
\end{minipage}

\markboth{\leftmark}{\thechapter.9\,\,\,VCN Application: Eff. Step Size for Square System}
\begin{minipage}{\textwidth}
    \section{VCN Application: Eff. Step Size for Square System} % X.9
        \label{appendix:vcn-success-small}
        Referenced by \autoref{fig:vcn-energy-conservation-small-system}.\\
        \vbox{
            \makebox[\textwidth][c]{
                \includegraphics[width=1.0\textwidth]{plotgeneration/vcn-square-small/energy-variance-step-dependent.pdf}
            }
        }
\end{minipage}

\markboth{\leftmark}{\thechapter.10\,\,\,VCN Application: Larger Square System}
\begin{minipage}{\textwidth}
    \section{VCN Application: Larger Square System} % X.10
        \label{appendix:vcn-success-big}
        Referenced at the end of \autoref{sec:experiments-vcn-successful}. The parameters are identical to the experiment in \autoref{fig:vcn-energy-conservation-small-system}, just with a larger system.\\
        \vbox{
            \makebox[\textwidth][c]{
                \includegraphics[width=1.0\textwidth]{plotgeneration/vcn-square-comparison/energy-variance.pdf}
            }
        }
\end{minipage}
\newpage
\begin{minipage}{\textwidth}
    \vbox{
        \makebox[\textwidth][c]{
            \includegraphics[width=1.0\textwidth]{plotgeneration/vcn-square-comparison/energy-variance-step-dependent.pdf}
        }
    }
\end{minipage}

\markboth{\leftmark}{\thechapter.11\,\,\,VCN Application: System Size Dependency}
\begin{minipage}{\textwidth}
    \section{VCN Application: System Size Dependency} % X.11
        \label{appendix:vcn-system-size-dependency}
        Referenced at the end of \autoref{sec:experiments-size-dependency-measurement} and further details in \autoref{fig:final-measurement-energy}\\
        \vbox{
            \makebox[\textwidth][c]{
                \includegraphics[width=1.0\textwidth]{plotgeneration/system-size-dependency/system-size-purity.pdf}
            }
        }
\end{minipage}
\newpage
\begin{minipage}{\textwidth}
    \vbox{
        \makebox[\textwidth][c]{
            \includegraphics[width=1.0\textwidth]{plotgeneration/system-size-dependency/system-size-concurrence.pdf}
        }
    }
\end{minipage}


\end{document}





